%% Generated by blueprint/scripts/render_registers.py from doc/SHIPPING-CHANGES.md. Do not edit.
\chapter{Shipping changes}\label{chap:shipping-changes}

\section*{At a glance}
\begin{itemize}
\item \lead{Source} \code{doc/SHIPPING-CHANGES.md}, rendered by \code{blueprint/scripts/render\_registers.py}; every section below is the register's own text.
\item \lead{Changes} 24 entries, one commit each, subject \code{shipping(S-<n>): ...}.
\item \lead{Reported, not fixed} 28 entries.
\end{itemize}

\begin{tabular}{lp{0.8\linewidth}}
\toprule
Entry & Title \\
\midrule
\ref{reg:S-1} & Change register, corpus and regression harness \\
\ref{reg:S-2} & Merge of \texttt{dev/\allowbreak{}zulip-issues} \\
\ref{reg:S-3} & Nested products are parenthesized in the \texttt{.mli} signature \\
\ref{reg:S-4} & The benchmark Makefile compiles a \texttt{.mli} that names \texttt{LeanArray} \\
\ref{reg:S-5} & \texttt{String} is no longer printed as \texttt{string} in the \texttt{.mli} signature \\
\ref{reg:S-6} & The benchmark Makefile regenerates a missing \texttt{.ast.inlinings} before running peregrine \\
\ref{reg:S-7} & Auto-inlining marks a constant only if its inlined body is small \\
\ref{reg:S-8} & Auto-inlining marks only constants whose erased body is a value \\
\ref{reg:S-9} & The dead \texttt{containsFix} guard of auto-inlining is removed \\
\ref{reg:S-10} & Auto-inlining does not mark constants tagged \texttt{@[noinline]} \\
\ref{reg:S-11} & The documentation of auto-inlining describes the shapes the code accepts \\
\ref{reg:S-12} & Port to Lean v4.33.0-rc2, with a Lake dependency on \texttt{lean4lean} \\
\ref{reg:S-13} & Sparse \texttt{casesOn} and per-constructor eliminators are erased by constructor name \\
\ref{reg:S-14} & A catch-all no longer re-evaluates a discriminant that is not a variable \\
\ref{reg:S-15} & The corpus gains examples inside and near the verification scope, irreducible-alias programs and unsafe recursion \\
\ref{reg:S-16} & The erasure traversal is total and generic over its backend \\
\ref{reg:S-17} & The traversal's context lists its locals, and binder names come from that list \\
\ref{reg:S-18} & The dependency closure of a program is computed by \texttt{collectDeps}, which \texttt{\#erase} does not call yet \\
\ref{reg:S-19} & A pure erasability oracle and the monad of a pure backend, which \texttt{\#erase} does not use yet \\
\ref{reg:S-20} & The pure backend is complete, and \texttt{erasePure} runs the traversal with it, which \texttt{\#erase} does not call yet \\
\ref{reg:S-21} & \texttt{\#erase} erases programs in the fragment with the pure path \\
\ref{reg:S-22} & Comments, docstrings and the register cite only what the repository contains \\
\ref{reg:S-23} & CI builds and tests the \texttt{shipping} branch when it is pushed \\
\ref{reg:S-24} & The docstring of \texttt{binderNameOf} gives the reason that holds \\
\bottomrule
\end{tabular}

\begin{tabular}{lp{0.8\linewidth}}
\toprule
Entry & Title \\
\midrule
\ref{reg:R-1} & Primitive integers are printed as quoted strings \\
\ref{reg:R-2} & String atoms are not escaped \\
\ref{reg:R-3} & Kernames are not injective \\
\ref{reg:R-4} & Unsupported literals are erased to \ensuremath{\Box} after a panic \\
\ref{reg:R-5} & An alternative whose \ensuremath{\Pi}-type is not syntactic gets a branch without binders \\
\ref{reg:R-6} & Peano literals exhaust the stack \\
\ref{reg:R-7} & Inductive types are always declared non-propositional \\
\ref{reg:R-8} & Informative index fields of large-eliminating Prop inductives are lost \\
\ref{reg:R-9} & Recursors and other constants without a value become axioms \\
\ref{reg:R-10} & Recursive structures get no projection declarations \\
\ref{reg:R-11} & Fixpoints have recursive-argument index 0 and bodies that are not \ensuremath{\eta}-expanded \\
\ref{reg:R-12} & Mutual blocks skip the \texttt{@[extern]} and \texttt{@[inline]} handling \\
\ref{reg:R-13} & \texttt{@[implemented\_\allowbreak{}by]} is ignored \\
\ref{reg:R-14} & Erasability is decided by the elaborator at default transparency \\
\ref{reg:R-15} & \texttt{\#erase} without \texttt{to} logs the program in place of the attributes \\
\ref{reg:R-16} & A log message lacks string interpolation \\
\ref{reg:R-17} & The \texttt{.mli} signature falls back to \texttt{unit} \\
\ref{reg:R-18} & The benchmark Makefile does not re-erase after a frontend change \\
\ref{reg:R-19} & \texttt{lake build} in \texttt{benchmarks/} builds nothing \\
\ref{reg:R-20} & The benchmark runtime disagrees with Lean on some operations \\
\ref{reg:R-21} & Axioms without an OCaml implementation fail only at link time \\
\ref{reg:R-22} & The benchmark \texttt{list\_\allowbreak{}sum\_\allowbreak{}rev} does not reverse \\
\ref{reg:R-23} & Documentation drift \\
\ref{reg:R-24} & Auto-inlining cannot be turned on in the benchmarks and does not shrink the \texttt{.ast} \\
\ref{reg:R-25} & The inlining diagnosis note is not committed, and its candidate causes are wrong \\
\ref{reg:R-26} & \texttt{[@inlined hint]} silences warning 55 without inlining anything \\
\ref{reg:R-27} & The benchmark Makefile builds \texttt{axioms.cmx} without waiting for \texttt{LeanArray.cmx} \\
\ref{reg:R-28} & The benchmark Makefile needs GNU make 4.4 and does not say so \\
\bottomrule
\end{tabular}

This register lists every change to the shipping code on the \texttt{shipping} branch: the eraser (\leanfile{LeanToLambdaBox.lean}, \texttt{LeanToLambdaBox/}), the package files (\texttt{lakefile.toml}, \texttt{lean-toolchain}, \texttt{lake-manifest.json}), the benchmarks (\texttt{benchmarks/}) and the tooling that checks them (\texttt{scripts/}, \texttt{tests/}, \texttt{.github/\allowbreak{}workflows/}). The branch starts from \texttt{main} at \texttt{58701f8}.

Each change has an id \texttt{S-\textless{}n\textgreater{}} and exactly one commit, whose subject is \texttt{shipping(S-\textless{}n\textgreater{}): \textless{}what changed and its observable effect\textgreater{}}. Its entry has these fields:

\begin{itemize}
\item \textbf{Commit:} the commit whose subject starts with \texttt{shipping(S-\textless{}n\textgreater{}):} (find it with \texttt{git log -{}-grep='\textasciicircum{}shipping(S-\textless{}n\textgreater{}):'}); for a merge, say so.
\item \textbf{Files and functions:} every file and function/declaration touched.
\item \textbf{Why necessary:} the concrete reason (a reproduced defect, a compile error, a requirement of the spec).
\item \textbf{Behaviour before:} observable behaviour, with the reproduction.
\item \textbf{Behaviour after:} observable behaviour, with the same reproduction.
\item \textbf{Effect on emitted .ast (corpus):} byte-level: which corpus files change (or "byte-identical for all N files"), and the nature of every difference.
\item \textbf{Regression test:} the test under \texttt{tests/\allowbreak{}regress/} that fails before and passes after (or, for a non-behavioural change, what it guards).
\end{itemize}

Defects found but not fixed are listed at the end, under "Reported, not fixed", as \texttt{R-\textless{}n\textgreater{}} entries.

\section{\texorpdfstring{The fixed corpus}{The fixed corpus}}\label{reg:the-fixed-corpus}

The corpus is everything the eraser emits for a fixed set of inputs:

\begin{itemize}
\item \textbf{Benchmarks.} The 20 programs of \texttt{benchmarks/\allowbreak{}TESTS} whose harness is \texttt{natio}, erased through \texttt{benchmarks/\allowbreak{}via\_\allowbreak{}malfunction/\allowbreak{}Makefile} as the benchmark pipeline does it, once with \texttt{PRUNE\_\allowbreak{}CONSTRUCTORS=1} (the Makefile default, used by 12 of the 14 benchmark backends) and once with \texttt{PRUNE\_\allowbreak{}CONSTRUCTORS=0}. Each run yields \texttt{\textless{}test\textgreater{}.ast}, \texttt{\textless{}test\textgreater{}.ast.inlinings} and \texttt{\textless{}test\textgreater{}.mli}: 120 files.
\item \textbf{Examples.} The files \texttt{tests/\allowbreak{}corpus/\allowbreak{}*.lean}. Each one imports \texttt{LeanToLambdaBox} and writes its outputs with \texttt{\#erase ... to \textquotedbl{}\textless{}file\textgreater{}\textquotedbl{}}. A file must elaborate without error; an \texttt{\#erase} that fails is wrapped in \texttt{\#guard\_\allowbreak{}msgs}, which pins its error, and writes no file.

\begin{itemize}
\item \texttt{PortProbe.lean}: Nat functions, a Prop-carrying argument, a subtype, a structure with a Prop field, wildcard matches, derived \texttt{BEq}/\texttt{DecidableEq} (3 and 10 constructors); each erased with the default configuration, with constructor pruning, and applied to a literal under \texttt{\{nat := .peano, extern := .preferLogical\}}.
\item \texttt{Scope.lean}: programs with no inductive type in their dependency closure (Church numerals and booleans, polymorphic combinators, proof and type arguments, \texttt{let} of values, types and proofs, universe polymorphism, a Prop-typed axiom used as an argument).
\item \texttt{Defects.lean}: reproductions of the \texttt{R-\textless{}n\textgreater{}} entries below.
\item \texttt{Examples.lean}, \texttt{Names.lean}, \texttt{NearMiss.lean}: example programs inside and near the verification scope. \texttt{Examples.lean} has 61 programs without an inductive type in their dependency closure (Church numerals and booleans, universe-polymorphic combinators at \texttt{Sort 2}, \texttt{Sort 1} and \texttt{Prop}, type and type-former arguments, sort and type aliases, proof arguments by axiom, theorem and \ensuremath{\lambda}, \texttt{let} and \texttt{have}, an opaque, \texttt{@[implemented\_\allowbreak{}by]}, \texttt{@[inline]} and \texttt{@[macro\_\allowbreak{}inline]}, a \texttt{Prop} alias made \texttt{@[irreducible]}, relevant axioms) and 48 readouts, which apply such a program to \texttt{Nat}, \texttt{Nat.succ}, \texttt{Nat.zero} (or \texttt{Bool}, \texttt{true}, \texttt{false}). \texttt{Names.lean} has 7 programs of the same kind whose constant names need escaping or collide as kernames, and 2 readouts. \texttt{NearMiss.lean} has 8 programs that each add one feature outside that fragment (a literal, a projection, a match, structural recursion, a quotient, a \texttt{partial def}, a term metavariable, a universe metavariable), and 2 readouts. Each program is erased under \texttt{\{nat := .peano\}} to \texttt{\textless{}name\textgreater{}.peano.ast} and under the default configuration to \texttt{\textless{}name\textgreater{}.default.ast}; each readout under \texttt{\{nat := .peano\}} only.
\item \texttt{IrrAlias.lean}: 3 programs without an inductive type whose types are propositions or \ensuremath{\Pi}-types only through an \texttt{@[irreducible]} alias, erased as the programs of \texttt{Examples.lean}.
\item \texttt{UnsafeRec.lean}: 3 programs without an inductive type that use \texttt{unsafe} recursion (a recursive constant, a two-member mutual block, a recursive constant whose value is not a \ensuremath{\lambda}), erased as the programs of \texttt{Examples.lean}.
\end{itemize}
\end{itemize}

Regenerate it from the current checkout with

\begin{quote}
\texttt{scripts/\allowbreak{}corpus.sh~OUTDIR}
\end{quote}

which writes \texttt{OUTDIR/\allowbreak{}benchmarks/\allowbreak{}\{prune,noprune\}/} and \texttt{OUTDIR/\allowbreak{}examples/\allowbreak{}\textless{}Stem\textgreater{}/}, plus logs and run information under \texttt{OUTDIR/\allowbreak{}\_\allowbreak{}meta/} (not part of the corpus; \texttt{\_\allowbreak{}meta/\allowbreak{}panics} counts \texttt{PANIC} messages per log). The script uses only \texttt{lake} and \texttt{make} from \texttt{PATH}; the toolchain is the one the checkout's \texttt{lean-toolchain} files name. Compare two corpora with

\begin{quote}
\texttt{scripts/\allowbreak{}corpus-diff.sh~[-{}-normalize]~[-{}-summary]~A~B}
\end{quote}

It prints the files that are identical, differing, or present on one side only, then a unified diff per differing file (with a line break before every constructor). \texttt{-{}-normalize} also compares the files after replacing hygienic name suffixes (\texttt{x.\_\allowbreak{}@.M.\_\allowbreak{}hyg.N}, \texttt{x.\_\allowbreak{}@.M.\textless{}hash\textgreater{}.\_\allowbreak{}hygCtx.\_\allowbreak{}hyg.N}) by a fixed token, to separate toolchain-induced renamings from real changes. The exit status is 0 when the corpora agree.

\section{\texorpdfstring{Regression tests}{Regression tests}}\label{reg:regression-tests}

A regression test is a file \texttt{tests/\allowbreak{}regress/\allowbreak{}\textless{}name\textgreater{}.lean} that writes its outputs into its working directory, normally with \texttt{\#erase ... to \textquotedbl{}\textless{}file\textgreater{}\textquotedbl{}}. Its expected outputs are \texttt{tests/\allowbreak{}regress/\allowbreak{}expected/\allowbreak{}\textless{}name\textgreater{}/}. Run

\begin{quote}
\texttt{scripts/\allowbreak{}regress.sh~[TEST...]}
\end{quote}

It builds the package, elaborates each test in a fresh directory, and compares the files written with the expected ones byte for byte. A test fails if Lean reports an error, if a \texttt{PANIC} message appears (unless the test contains the line \texttt{-{}- regress: allow-panic}), or if any file is missing, extra or different. With \texttt{PEREGRINE=\textless{}path to the peregrine binary\textgreater{}}, it also runs the lines \texttt{-{}- peregrine: validate \textless{}file\textgreater{} [\textless{}option\textgreater{}...]} and \texttt{-{}- peregrine: eval \textless{}file\textgreater{} [\textless{}option\textgreater{}...]} of each test; each must succeed and, when \texttt{tests/\allowbreak{}regress/\allowbreak{}expected-peregrine/\allowbreak{}\textless{}name\textgreater{}/\allowbreak{}\textless{}file\textgreater{}.\textless{}verb\textgreater{}} exists, print exactly that. \texttt{scripts/\allowbreak{}regress.sh -{}-update [TEST...]} rewrites the expected files from the current checkout. CI (\texttt{.github/\allowbreak{}workflows/\allowbreak{}build.yml}, on pushes to \texttt{main} and \texttt{shipping} and on pull requests) runs \texttt{scripts/\allowbreak{}regress.sh} after the build, without peregrine.

The pure path of the eraser (\leandecl{Erasure.erasePure}, S-20), which \texttt{\#erase} takes on the programs in the fragment (S-21), is compared with the \texttt{Meta} path (\leandecl{Erasure.erase}) on the corpus examples by

\begin{quote}
\texttt{scripts/\allowbreak{}pure-harness.sh~[-{}-update]~[-{}-keep~DIR]}
\end{quote}

It elaborates each file \texttt{tests/\allowbreak{}corpus/\allowbreak{}\textless{}Stem\textgreater{}.lean} with the harness \leanfile{tests/\allowbreak{}harness/\allowbreak{}PureHarness.lean} imported. For each \texttt{\#erase ... to \textquotedbl{}\textless{}f\textgreater{}\textquotedbl{}}, the harness computes \texttt{collectDeps} over the elaboration environment; for a program in the fragment it runs \texttt{erasePure} and the \texttt{Meta} path (\leandecl{Erasure.erase}), and, at every call of the pure oracle in the pure run, the \texttt{Meta} oracle on the same term in the same local context. Then \texttt{\#erase} runs as usual. The summary has one line per \texttt{\#erase}: out of the fragment and why; or the results of both paths, whether their outputs are byte-identical, which of them the file written by \texttt{\#erase} equals, the number of oracle calls, and the calls at which the two oracles answer differently. It must equal \texttt{tests/\allowbreak{}harness/\allowbreak{}expected/\allowbreak{}pure-harness.txt}; \texttt{-{}-update} rewrites that file. CI runs the harness after \texttt{scripts/\allowbreak{}regress.sh}.

\section{\texorpdfstring{S-1: Change register, corpus and regression harness}{S-1: Change register, corpus and regression harness}}\label{reg:S-1}

\begin{itemize}
\item \textbf{Commit:} the commit whose subject starts with \texttt{shipping(S-1):} (\texttt{git log -{}-grep='\textasciicircum{}shipping(S-1):'}).
\item \textbf{Files and functions:}

\begin{itemize}
\item new \texttt{doc/\allowbreak{}SHIPPING-CHANGES.md} (this register);
\item new \texttt{scripts/\allowbreak{}common.sh} (\texttt{die}, \texttt{build\_\allowbreak{}frontend}, \texttt{run\_\allowbreak{}lean}, \texttt{tokenize}, \texttt{normalize\_\allowbreak{}hygiene}), \texttt{scripts/\allowbreak{}corpus.sh}, \texttt{scripts/\allowbreak{}corpus-diff.sh}, \texttt{scripts/\allowbreak{}regress.sh};
\item new \leanfile{tests/\allowbreak{}corpus/\allowbreak{}PortProbe.lean}, \leanfile{tests/\allowbreak{}corpus/\allowbreak{}Scope.lean}, \leanfile{tests/\allowbreak{}corpus/\allowbreak{}Defects.lean};
\item new \leanfile{tests/\allowbreak{}regress/\allowbreak{}smoke.lean}, \texttt{tests/\allowbreak{}regress/\allowbreak{}expected/\allowbreak{}smoke/} (5 files), \texttt{tests/\allowbreak{}regress/\allowbreak{}expected-peregrine/\allowbreak{}smoke/} (3 files), \texttt{tests/\allowbreak{}regress/\allowbreak{}.gitignore} (re-includes the \texttt{*.ast} files that the root \texttt{.gitignore} excludes), \texttt{tests/\allowbreak{}regress/\allowbreak{}.gitattributes} (expected files are kept verbatim: no end-of-line conversion, no whitespace check);
\item \texttt{.github/\allowbreak{}workflows/\allowbreak{}build.yml}: new step "Shipping regression tests".
\end{itemize}

No file under \texttt{LeanToLambdaBox/}, no package file and no file under \texttt{benchmarks/} changes.
\item \textbf{Why necessary:} the specification (\S{}5.1) requires this register, the byte-level effect of each change on a fixed corpus of the benchmark programs plus own examples, and a regression test per change; \S{}5.3 asks CI to run the shipping regression tests.
\item \textbf{Behaviour before:} the eraser as on \texttt{main}; there was no register, no corpus tooling and no test, and CI ran only \texttt{lake build}.
\item \textbf{Behaviour after:} the eraser is unchanged. \texttt{scripts/\allowbreak{}regress.sh} prints \texttt{ok   smoke} and \texttt{all 1 tests passed}; with \texttt{PEREGRINE} set, \texttt{fact3.ast} validates and evaluates to \texttt{Nat.succ} applied six times to \texttt{Nat.zero} (\texttt{fact 3 = 6}). \texttt{scripts/\allowbreak{}corpus.sh OUTDIR} writes 284 files.
\item \textbf{Effect on emitted .ast (corpus):} byte-identical for all 284 files (116 \texttt{.ast}, 116 \texttt{.ast.inlinings}, 52 \texttt{.mli}), since no shipping code changes. The benchmark part (120 files) is byte-identical to the outputs of an unmodified \texttt{58701f8} checkout erased through the same Makefile rules, and two runs of \texttt{scripts/\allowbreak{}corpus.sh} give byte-identical corpora.
\item \textbf{Regression test:} \leanfile{tests/\allowbreak{}regress/\allowbreak{}smoke.lean} guards the README example under the default configuration and a closed recursive program under Peano naturals (\texttt{fact 3}).
\end{itemize}

\section{\texorpdfstring{S-2: Merge of \texttt{dev/\allowbreak{}zulip-issues}}{S-2: Merge of dev/zulip-issues}}\label{reg:S-2}

\begin{itemize}
\item \textbf{Commit:} the merge commit whose subject starts with \texttt{shipping(S-2):} (\texttt{git log -{}-grep='\textasciicircum{}shipping(S-2):'}). It is a \texttt{-{}-no-ff} merge of \texttt{dev/\allowbreak{}zulip-issues} at \texttt{28955eb} (\texttt{main} plus \texttt{94a429c}, \texttt{e46224f} and \texttt{8db2dd9}), and it also adds the tests and register text listed below.
\item \textbf{Files and functions:}

\begin{itemize}
\item \leanfile{LeanToLambdaBox/\allowbreak{}Erasure.lean}: \texttt{ErasureConfig} (new field \texttt{auto\_\allowbreak{}inline\_\allowbreak{}typeclass\_\allowbreak{}dispatch}, default \texttt{false}); new \texttt{LBTerm.stripLambdas}, \texttt{LBTerm.containsFix}, \texttt{LBTerm.isTrivialAlias}; \texttt{erase.visitMutual} (the \texttt{@[inline]} check is computed first, as \texttt{leanInline}; new post-erasure auto-inline step in the non-recursive branch); \texttt{MLType} (new constructors \texttt{string}, \texttt{option}, \texttt{array}, \texttt{prod}); \texttt{MLType.toString} (now \texttt{partial}; \texttt{protArrow} and \texttt{protCtor} replace \texttt{toStringProtected}); \texttt{to\_\allowbreak{}ml\_\allowbreak{}type} (new cases \texttt{Int}, \texttt{String}, \texttt{Option}, \texttt{Array}, \texttt{Prod}).
\item \texttt{benchmarks/\allowbreak{}via\_\allowbreak{}malfunction/\allowbreak{}Makefile}: the \texttt{lake lean} rule also names \texttt{\%.ast.inlinings} as a target.
\item \texttt{benchmarks/\allowbreak{}via\_\allowbreak{}malfunction/\allowbreak{}nat-inline.ml} (15 attributes) and \texttt{int-inline.ml} (9 attributes): \texttt{[@inlined]} becomes \texttt{[@inlined hint]} on Zarith calls.
\item new \leanfile{tests/\allowbreak{}regress/\allowbreak{}mli\_\allowbreak{}types.lean}, \leanfile{tests/\allowbreak{}regress/\allowbreak{}auto\_\allowbreak{}inline.lean}, \leanfile{tests/\allowbreak{}regress/\allowbreak{}runtime\_\allowbreak{}inlined\_\allowbreak{}hint.lean}, their expected outputs in \texttt{tests/\allowbreak{}regress/\allowbreak{}expected/\allowbreak{}\textless{}test\textgreater{}/}, and \texttt{tests/\allowbreak{}regress/\allowbreak{}expected-peregrine/\allowbreak{}auto\_\allowbreak{}inline/};
\item \texttt{tests/\allowbreak{}regress/\allowbreak{}expected/\allowbreak{}smoke/\allowbreak{}fact3.ast}: re-baselined; its 6 hygienic suffixes grow by 1, as in the corpus below (\texttt{fact3.ast} is written with a \texttt{config} clause);
\item \texttt{scripts/\allowbreak{}regress.sh}: header comment only (a test may write its outputs without \texttt{\#erase});
\item \texttt{doc/\allowbreak{}SHIPPING-CHANGES.md}: this entry; the paragraph on regression tests; R-17 and R-18 updated to the merged code; new R-24, R-25, R-26.
\end{itemize}
\item \textbf{Why necessary:} \S{}5.1 of the specification requires merging \texttt{dev/\allowbreak{}zulip-issues} after checking, before and after, each issue it claims to address. Its commit messages make five claims; each was reproduced on this branch before and after the merge, with these verdicts:

\begin{itemize}
\item C1 (\texttt{94a429c}): the \texttt{.mli} printer covers \texttt{Int}, \texttt{String}, \texttt{Option}, \texttt{Array} and \texttt{Prod}, with OCaml-correct precedence. \textbf{Partially correct}; defects D1, D2, D3.
\item C2 (\texttt{94a429c}, \texttt{e46224f}, \texttt{8db2dd9}): \texttt{auto\_\allowbreak{}inline\_\allowbreak{}typeclass\_\allowbreak{}dispatch} marks typeclass-dispatch constants for inlining. \textbf{Correct as an option that is off by default and leaves the default output unchanged; defective when turned on}: defects D5 to D10.
\item C3 (\texttt{94a429c}): \texttt{[@inlined hint]} silences OCaml warning 55. \textbf{Correct; the generated code is unchanged} (D12).
\item C4 (\texttt{94a429c}): declaring \texttt{\%.ast.inlinings} as an output of the \texttt{lake lean} rule prevents peregrine's "no \ldots{} inlinings file" failure. \textbf{Not fixed by this change alone} (D4).
\item C5 (\texttt{94a429c}): the note \texttt{references/\allowbreak{}inlining\_\allowbreak{}diagnosis.md} shows that the kernames match and that the \texttt{instDecidableEqNat} symptom lies in peregrine. \textbf{The conclusion holds; the note's candidate causes are wrong, and no commit contains the note} (D11).
\end{itemize}

The defects found in the branch are numbered D1 to D12 below. D1 to D9 remain at this commit and are fixed in later entries of this register. D10, D11 and D12 are informational and are reported as R-24, R-25 and R-26.
\item \textbf{Behaviour before:} the native checks of C1 compile the erased program with peregrine (\texttt{unbox.config}, after rewriting \texttt{(primInt \textquotedbl{}N\textquotedbl{})} to \texttt{(primInt N)} because of R-1), malfunction and \texttt{ocamlopt} 4.14.2 without flambda, and link it with the \texttt{via\_\allowbreak{}malfunction} runtime and an OCaml harness written against the emitted \texttt{.mli}.

\begin{itemize}
\item C1: a result of type \texttt{Int}, \texttt{Option Nat}, \texttt{Nat \ensuremath{\times} Bool}, \texttt{List (Nat \ensuremath{\times} Nat)}, \texttt{Option (Nat \ensuremath{\to} Nat)}, \texttt{(Nat \ensuremath{\to} Nat) \ensuremath{\times} Nat}, \texttt{Array Nat}, \texttt{(Nat \ensuremath{\times} Nat) \ensuremath{\times} Nat}, \texttt{Nat \ensuremath{\times} (Nat \ensuremath{\times} Nat)} or \texttt{String} gets the signature \texttt{unit} (\texttt{unit list} for the list), with the warning \texttt{failed to translate \ldots{} into ML type, emitting unit instead}. \texttt{List (Nat \ensuremath{\to} Nat)} is printed \texttt{Z.t -\textgreater{} Z.t -\textgreater{} Z.t list}, which OCaml reads as a function of two arguments; a harness that matches a list does not type-check (\texttt{This pattern should not be a list literal}).
\item C2: \texttt{config \{auto\_\allowbreak{}inline\_\allowbreak{}typeclass\_\allowbreak{}dispatch := true\}} is an error: \texttt{'auto\_\allowbreak{}inline\_\allowbreak{}typeclass\_\allowbreak{}dispatch' is not a field of structure 'Erasure.ErasureConfig'}.
\item C3: \texttt{make build/\allowbreak{}\textless{}id\textgreater{}/\allowbreak{}axioms.cmx FLAMBDA=0 MALFUNCTION\_\allowbreak{}NO\_\allowbreak{}FLAMBDA\_\allowbreak{}SWITCH=peregrine ARRAYML=JCFArrayOCaml4.ml} in \texttt{benchmarks/\allowbreak{}via\_\allowbreak{}malfunction} prints warning 55 (\texttt{Cannot inline: Function information unavailable}) 20 times, for \texttt{nat.ml} and \texttt{int.ml}.
\item C4: in \texttt{benchmarks/\allowbreak{}via\_\allowbreak{}malfunction}, after \texttt{make build/\allowbreak{}\textless{}id\textgreater{}/\allowbreak{}even.mlf} (with the options of C3 and \texttt{PEREGRINE} set to a wrapper applying the R-1 rewrite), delete \texttt{even.ast.inlinings} and \texttt{even.mlf} and run the same command: peregrine fails with \texttt{option '-{}-attributes': invalid element in list (build/\allowbreak{}\textless{}id\textgreater{}/\allowbreak{}even.ast.inlinings): no build/\allowbreak{}\textless{}id\textgreater{}/\allowbreak{}even.ast.inlinings file or directory}, the error reported on Zulip on Feb 17. With a copy of the Makefile whose \texttt{INLINING=1} rule for \texttt{\%.mlf} also lists \texttt{\$(build)/\allowbreak{}\%.ast.inlinings} as a prerequisite, make stops with \texttt{No rule to make target 'build/\allowbreak{}\textless{}id\textgreater{}/\allowbreak{}even.mlf'}, the failure reported on Zulip on Feb 13 (\texttt{No rule to make target 'bin/\allowbreak{}even'}).
\item C5: \texttt{scr n := isZeroIf n + (match isZeroDec n with \textbar{} .isTrue \_\allowbreak{} =\textgreater{} 10 \textbar{} .isFalse \_\allowbreak{} =\textgreater{} 20)}, with \texttt{isZeroIf n := if n = 0 then 1 else 2} and \texttt{isZeroDec n : Decidable (n = 0) := inferInstance}, erased with \texttt{\{remove\_\allowbreak{}irrel\_\allowbreak{}constr\_\allowbreak{}args := true\}}: \texttt{instDecidableEqNat} has the same kername in its call sites, its declaration and \texttt{scr.ast.inlinings}. After \texttt{peregrine compile scr.ast unbox.config -{}-attributes=scr.ast.inlinings}, \texttt{isZeroDec} calls \texttt{\$def\_\allowbreak{}\_\allowbreak{}Nat\_\allowbreak{}decEq}, but \texttt{isZeroIf} still computes \texttt{(let (\$discr (apply \$def\_\allowbreak{}\_\allowbreak{}instDecidableEqNat \$n \ldots{})) (switch \$discr \ldots{}))}.
\end{itemize}
\item \textbf{Behaviour after:}

\begin{itemize}
\item C1: the same results get \texttt{Z.t}, \texttt{Z.t option}, \texttt{Z.t * bool}, \texttt{(Z.t * Z.t) list}, \texttt{(Z.t -\textgreater{} Z.t) option}, \texttt{(Z.t -\textgreater{} Z.t) * Z.t}, \texttt{(Z.t -\textgreater{} Z.t) list} and \texttt{Z.t LeanArray.array}, without warning, and the harnesses print the Lean values (\texttt{-7}, \texttt{None /\allowbreak{} Some 5}, \texttt{0 true}, \texttt{5 6}, \texttt{6}, \texttt{6 5}, \texttt{6}, \texttt{2}). Remaining defects:

\begin{itemize}
\item D1: both nested products are printed \texttt{Z.t * Z.t * Z.t}, an OCaml triple, while the value is a pair with a pair inside. The harness written against it crashes (exit status 139) for \texttt{Nat \ensuremath{\times} (Nat \ensuremath{\times} Nat)} and prints garbage for \texttt{(Nat \ensuremath{\times} Nat) \ensuremath{\times} Nat}; with the signatures corrected by hand to \texttt{Z.t * (Z.t * Z.t)} and \texttt{(Z.t * Z.t) * Z.t}, both print \texttt{5 6 7}. Before the merge both were a warning and \texttt{unit}. \texttt{MLType.toString} prints the components of \texttt{prod} with \texttt{protArrow}, which does not parenthesize products.
\item D2: the benchmark Makefile cannot compile the \texttt{.mli} of an \texttt{Array} result: its \texttt{\%.cmi} rule runs \texttt{ocamlfind ocamlopt -package zarith -c build/\allowbreak{}\textless{}id\textgreater{}/\allowbreak{}\textless{}test\textgreater{}.mli} without \texttt{-I \$(build)} and fails with \texttt{Unbound module LeanArray}. The same file compiles with \texttt{-I}.
\item D3: \texttt{String} is printed \texttt{string}, but nothing represents a Lean \texttt{String} as an OCaml string: a program returning \texttt{String.mk \ldots{}} references axioms that \texttt{axioms.ml} does not implement (\texttt{malfunction cmx} fails with \texttt{Unbound value Axioms.def\_\allowbreak{}\_\allowbreak{}String\_\allowbreak{}mk}), and string literals become \texttt{\ensuremath{\Box}} (R-4).
\item A user inductive result type, the case raised on Zulip on Feb 10, still gets \texttt{unit} (R-17).
\end{itemize}
\item C2: the option exists and is \texttt{false} by default. With it on, the 20 natio benchmarks erased with \texttt{\{remove\_\allowbreak{}irrel\_\allowbreak{}constr\_\allowbreak{}args := true, auto\_\allowbreak{}inline\_\allowbreak{}typeclass\_\allowbreak{}dispatch := true\}} have the same \texttt{.ast} as without the option, up to hygienic suffixes (see below), and longer \texttt{.ast.inlinings} (binarytrees 3 to 19 entries, unionfind 16 to 37) that add instances and projection chains (\texttt{instHAdd}, \texttt{instAddNat}, \texttt{HAdd.hAdd}, \texttt{OfNat.ofNat}, \ldots{}). \texttt{peregrine eval -{}-attributes=\ldots{}} of \texttt{binarytrees 4}, \texttt{triangle\_\allowbreak{}rec 12}, \texttt{iflazy 7}, \texttt{even 9} and \texttt{const\_\allowbreak{}fold 3} (Peano naturals, logical externs) gives 610, 66, 42, 0 and 22, the values of Lean's \texttt{\#eval}, before the merge, after it, and with the option on. Remaining defects, all with the option on:

\begin{itemize}
\item D5: unbounded code growth. The \texttt{.mlf} of unionfind grows from 126 374 to 14 510 809 bytes and peregrine's compile time from 0.04 s to 3.27 s (unionfind\_noinline: 89 106 to 5 818 655 bytes); among the marked constants are the monad instances \texttt{Id.instMonad}, \texttt{UnionFind.StateT'.instMonad} and \texttt{UnionFind.ExceptT'.instMonad}. \texttt{e46224f} limited inlined instances to 40 erased nodes; \texttt{8db2dd9} removed that limit (\texttt{autoInlineMaxBodySize}, \texttt{LBTerm.size}) without saying so.
\item D6: repeated work. Every instance is marked, whatever its body. With \texttt{instance instTbl : Tbl := let s := slowSum 100000; \ensuremath{\langle}fun i =\textgreater{} s + i\ensuremath{\rangle}} and \texttt{useTbl n := (List.range n).foldl (fun acc i =\textgreater{} acc + Tbl.get i) 0}, the native \texttt{useTbl 1000} takes 0.007 s with the option off and 0.764 s with it on (\texttt{useTbl 3000}: 0.010 s and 2.257 s): \texttt{slowSum 100000} is recomputed at every use.
\item D7: the guard \texttt{!t.containsFix}, meant to refuse recursive bodies, never applies: \texttt{LBTerm.fix} is built only in the recursive branch of \texttt{erase.visitMutual}, and the guard is in the non-recursive branch.
\item D8: \texttt{@[noinline]} is ignored: \texttt{@[noinline] instance instBar : Inhabited Nat := \ensuremath{\langle}42\ensuremath{\rangle}} is logged \texttt{Auto-inlining typeclass instance instBar.} and listed in \texttt{.ast.inlinings}.
\item D9: the documentation does not match the code. The docstring of the option promises "a single-ctor structure literal whose fields are shallow", and that of \texttt{LBTerm.isTrivialAlias} "a single-ctor structure literal (the usual \texttt{Foo.mk arg\ensuremath{_1} \ldots{} arg\ensuremath{_n}} shape \ldots{})"; the code accepts \texttt{.construct \_\allowbreak{} 0 \_\allowbreak{}} after stripping \ensuremath{\lambda}s, which matches only argument-less constructors of index 0, since constructors are emitted in applied form, and it also accepts constant functions. With the option on, \texttt{fLit : Bool := false} and \texttt{kfun (\_\allowbreak{} : Nat) : Nat := seven} are marked; \texttt{tLit : Bool := true} and the structure literal \texttt{pLit : P := \ensuremath{\langle}1, 2\ensuremath{\rangle}} are not.
\item D10 (R-24): the benchmark Makefile cannot turn the option on, and it does not make the \texttt{.ast} smaller.
\end{itemize}
\item C3: the same \texttt{make} prints no warning 55. The Cmm of \texttt{nat.ml} and \texttt{int.ml} (\texttt{-dcmm}) is the same as before up to source locations: every Zarith call is still an out-of-line call (\texttt{app \textquotedbl{}camlZ\_\allowbreak{}\_\allowbreak{}\ldots{}\textquotedbl{}}, 12 in \texttt{nat.ml}, 8 in \texttt{int.ml}). D12 (R-26).
\item C4: the same experiment with the repository Makefile fails with the same peregrine error. With the prerequisite added, make now re-runs \texttt{lake lean} and builds \texttt{even.mlf}: the co-target is one half of the fix. Remaining defect:

\begin{itemize}
\item D4: the \texttt{\%.mlf} rules do not list \texttt{\%.ast.inlinings} as a prerequisite, so make does not regenerate a missing \texttt{.ast.inlinings} before running peregrine.
\end{itemize}
\item C5: the same \texttt{.mlf}. Peregrine's inlining pass (MetaRocq 1.5.1 \texttt{EInlining.inline}, extracted to \texttt{peregrine-tool/\allowbreak{}\_\allowbreak{}build/\allowbreak{}default/\allowbreak{}src/\allowbreak{}extraction/\allowbreak{}EInlining.ml:32-33}) does not rewrite \texttt{tCase} scrutinees, and \texttt{if n = 0} puts \texttt{instDecidableEqNat n 0} in one. D11 (R-25).
\end{itemize}
\item \textbf{Effect on emitted .ast (corpus):} of 284 files, 214 are byte-identical and 70 \texttt{.ast} files differ only in hygienic binder names; all 116 \texttt{.ast.inlinings} and 52 \texttt{.mli} files are byte-identical. \texttt{scripts/\allowbreak{}corpus-diff.sh -{}-normalize} reports 214 identical, 70 identical after normalization, 0 differing. Every difference is a suffix \texttt{\_\allowbreak{}hyg.N} of a binder name (the \texttt{\_\allowbreak{}alt} binders that \texttt{inlineMatchers} creates during erasure, and binders of definitions), whose number grows by the number of \texttt{config} clauses elaborated in the same file before the name was created: each \texttt{config} clause now takes one more macro scope, as \texttt{ErasureConfig} has one more field.

\begin{itemize}
\item Benchmarks (the Makefile writes \texttt{config \{remove\_\allowbreak{}irrel\_\allowbreak{}constr\_\allowbreak{}args := true,\}} for \texttt{prune} and \texttt{config \{\}} for \texttt{noprune}): in each variant 19 of the 20 \texttt{.ast} files differ, in 300 suffixes, all by +1; \texttt{iflazy.ast} has no hygienic name and is identical.
\item \texttt{examples/\allowbreak{}Defects} (7 \texttt{.ast} files) and \texttt{examples/\allowbreak{}PortProbe} (25): in a file written by the k-th \texttt{\#erase} with a \texttt{config} clause, or by a later \texttt{\#erase} without one, the suffixes grow by 0 to k, and by k for the \texttt{\_\allowbreak{}alt} binders created by that \texttt{\#erase} (k is at most 22).
\item \texttt{examples/\allowbreak{}Scope}: identical.
\end{itemize}

Without a \texttt{config} clause the output is unchanged: the 20 natio benchmarks erased with a bare \texttt{\#erase \textless{}test\textgreater{} to \textquotedbl{}\textless{}test\textgreater{}.ast\textquotedbl{} mli \textquotedbl{}\textless{}test\textgreater{}.mli\textquotedbl{}} give byte-identical \texttt{.ast}, \texttt{.ast.inlinings} and \texttt{.mli} files (60) before and after the merge. Binder names carry no meaning in \ensuremath{\lambda}\ensuremath{\Box}, whose terms use de Bruijn indices. In the logs, the signature logged for \texttt{\#erase \textquotedbl{}abc\textquotedbl{}} (no \texttt{mli} path) becomes \texttt{val main: string}, without the warning.
\item \textbf{Regression test:}

\begin{itemize}
\item \leanfile{tests/\allowbreak{}regress/\allowbreak{}mli\_\allowbreak{}types.lean} fails before (8 of its 9 \texttt{.mli} files differ) and passes after. It holds the eight signatures checked natively above, and \texttt{(Nat \ensuremath{\to} Nat) \ensuremath{\to} Nat}, printed \texttt{(Z.t -\textgreater{} Z.t) -\textgreater{} Z.t} before and after (checked natively too), which guards arrows in argument position. The cases of D1 to D3 are left to their fixes.
\item \leanfile{tests/\allowbreak{}regress/\allowbreak{}auto\_\allowbreak{}inline.lean} fails before (the option is not a field) and passes after. Its outputs \texttt{default.ast} and \texttt{default.ast.inlinings} are byte-identical to those of the code before the merge, \texttt{off.*} equals \texttt{default.*}, \texttt{on.ast} equals \texttt{default.ast}, and \texttt{on.ast.inlinings} adds six constants. With \texttt{PEREGRINE} set, \texttt{on3.ast} validates, and \texttt{default3.ast} and \texttt{on3.ast} evaluate to 3 with their inlinings.
\item \leanfile{tests/\allowbreak{}regress/\allowbreak{}runtime\_\allowbreak{}inlined\_\allowbreak{}hint.lean} checks the runtime files instead of calling \texttt{\#erase}. It fails before (\texttt{nat-inline.ml: 15 bare [@inlined] attribute(s)}) and, after, writes \texttt{nat-inline.ml: [@inlined hint] 15, [@inlined] 0} and \texttt{int-inline.ml: [@inlined hint] 9, [@inlined] 0}.
\item C4 has no test, as the co-target alone changes no observable behaviour; the fix of D4 adds one. C5 concerns peregrine and a file outside the repository.
\item \leanfile{tests/\allowbreak{}regress/\allowbreak{}smoke.lean} passes with \texttt{fact3.ast} re-baselined (hygienic suffixes only); its other outputs and its peregrine outputs are unchanged.
\end{itemize}
\end{itemize}

\section{\texorpdfstring{S-3: Nested products are parenthesized in the \texttt{.mli} signature}{S-3: Nested products are parenthesized in the .mli signature}}\label{reg:S-3}

\begin{itemize}
\item \textbf{Commit:} the commit whose subject starts with \texttt{shipping(S-3):} (\texttt{git log -{}-grep='\textasciicircum{}shipping(S-3):'}).
\item \textbf{Files and functions:}

\begin{itemize}
\item \leanfile{LeanToLambdaBox/\allowbreak{}Erasure.lean}: \texttt{MLType.toString}, case \texttt{prod}: both components are printed with \texttt{protCtor}, which parenthesizes arrows and products, instead of \texttt{protArrow}, which parenthesizes only arrows.
\item \leanfile{tests/\allowbreak{}regress/\allowbreak{}mli\_\allowbreak{}types.lean}: docstring; new cases \texttt{rNestL}, \texttt{rNestR}, \texttt{rNestArg}; their 9 expected outputs in \texttt{tests/\allowbreak{}regress/\allowbreak{}expected/\allowbreak{}mli\_\allowbreak{}types/}.
\item \texttt{doc/\allowbreak{}SHIPPING-CHANGES.md}: this entry.
\end{itemize}
\item \textbf{Why necessary:} defect D1 of S-2, reproduced on this branch. A product nested in a product is printed as an OCaml triple, while the erased value is a pair containing a pair, so an OCaml harness written against the \texttt{.mli} type-checks and then crashes or reads garbage without any warning. Before S-2 the same programs gave a warning and \texttt{unit}.
\item \textbf{Behaviour before:} the programs \texttt{rNestL n := ((n, n + 1), n + 2)}, \texttt{rNestR n := (n, (n + 1, n + 2))}, \texttt{rNestArg (p : (Nat \ensuremath{\times} Nat) \ensuremath{\times} Nat) := p.1.1 + p.1.2 + p.2} and \texttt{rNestFun n : (Nat \ensuremath{\times} Nat) \ensuremath{\times} (Nat \ensuremath{\to} Nat) := ((n, n + 1), (\ensuremath{\cdot} + n))} are erased with \texttt{\{remove\_\allowbreak{}irrel\_\allowbreak{}constr\_\allowbreak{}args := true\}} and compiled natively as for S-2 (peregrine with \texttt{unbox.config} after the R-1 rewrite, malfunction, OCaml 4.14.2 without flambda, the \texttt{via\_\allowbreak{}malfunction} runtime objects), each linked with a harness written against its \texttt{.mli}:

\begin{itemize}
\item \texttt{rNestL.mli} is \texttt{val main: Z.t -\textgreater{} Z.t * Z.t * Z.t}; the harness \texttt{let (a, b, c) = RNestL.main (Z.of\_\allowbreak{}int 5)} prints a garbage number of about 235 digits (it varies between runs), then \texttt{7 3195};
\item \texttt{rNestR.mli} is the same; its harness exits with status 139 (segmentation fault);
\item \texttt{rNestArg.mli} is \texttt{val main: Z.t * Z.t * Z.t -\textgreater{} Z.t}; the harness passing \texttt{(5, 6, 7)} exits with status 139;
\item \texttt{rNestFun.mli} is \texttt{val main: Z.t -\textgreater{} Z.t * Z.t * (Z.t -\textgreater{} Z.t)}; its harness exits with status 139.
\end{itemize}
\item \textbf{Behaviour after:} the signatures are \texttt{val main: Z.t -\textgreater{} (Z.t * Z.t) * Z.t}, \texttt{val main: Z.t -\textgreater{} Z.t * (Z.t * Z.t)}, \texttt{val main: (Z.t * Z.t) * Z.t -\textgreater{} Z.t} and \texttt{val main: Z.t -\textgreater{} (Z.t * Z.t) * (Z.t -\textgreater{} Z.t)}. Harnesses written against them print \texttt{5 6 7}, \texttt{5 6 7}, \texttt{18} and \texttt{5 6 15} (the Lean values; the function is applied to 10). The harnesses of before no longer compile (\texttt{This expression has type (Z.t * Z.t) * Z.t but an expression was expected of type 'a * 'b * 'c}). The \texttt{.ast} and \texttt{.ast.inlinings} of these programs are byte-identical before and after, and products that are not nested print as before (\texttt{rPair}, \texttt{rFunPair} and \texttt{rListPair} of \leanfile{tests/\allowbreak{}regress/\allowbreak{}mli\_\allowbreak{}types.lean} are unchanged).
\item \textbf{Effect on emitted .ast (corpus):} byte-identical for all 284 files (116 \texttt{.ast}, 116 \texttt{.ast.inlinings}, 52 \texttt{.mli}); \texttt{scripts/\allowbreak{}corpus-diff.sh} against the corpus of S-2 reports 284 identical. No corpus \texttt{.mli} contains a product: 51 are \texttt{val main: Z.t -\textgreater{} Z.t}, one is \texttt{val main: Z.t -\textgreater{} unit}.
\item \textbf{Regression test:} \leanfile{tests/\allowbreak{}regress/\allowbreak{}mli\_\allowbreak{}types.lean}, cases \texttt{rNestL}, \texttt{rNestR} and \texttt{rNestArg}. It fails before (their 3 \texttt{.mli} files are \texttt{Z.t * Z.t * Z.t} forms) and passes after; the 27 outputs of its other cases are unchanged.
\end{itemize}

\section{\texorpdfstring{S-4: The benchmark Makefile compiles a \texttt{.mli} that names \texttt{LeanArray}}{S-4: The benchmark Makefile compiles a .mli that names LeanArray}}\label{reg:S-4}

\begin{itemize}
\item \textbf{Commit:} the commit whose subject starts with \texttt{shipping(S-4):} (\texttt{git log -{}-grep='\textasciicircum{}shipping(S-4):'}).
\item \textbf{Files and functions:}

\begin{itemize}
\item \texttt{benchmarks/\allowbreak{}via\_\allowbreak{}malfunction/\allowbreak{}Makefile}: pattern rule \texttt{\$(build)/\allowbreak{}\%.cmi} (new prerequisite \texttt{\$(build)/\allowbreak{}LeanArray.cmi}; the command gains \texttt{-I \$(build)}; one comment line added); new explicit rules \texttt{\$(build)/\allowbreak{}decidable.cmi}, \texttt{\$(build)/\allowbreak{}eq.cmi} and \texttt{\$(build)/\allowbreak{}LeanArray.cmi}, each with the command the pattern rule gave it before (\texttt{\$(OCAMLOPT) -c \$\textless{}}), in place of the comments \texttt{\# \textless{}module\textgreater{}.cmi handled by generic rule (no prerequisites)}.
\item new \leanfile{tests/\allowbreak{}regress/\allowbreak{}makefile\_\allowbreak{}cmi.lean} and \texttt{tests/\allowbreak{}regress/\allowbreak{}expected/\allowbreak{}makefile\_\allowbreak{}cmi/\allowbreak{}cmi\_\allowbreak{}commands.txt}.
\item \texttt{doc/\allowbreak{}SHIPPING-CHANGES.md}: this entry.
\end{itemize}
\item \textbf{Why necessary:} defect D2 of S-2, reproduced on this branch. Since S-2 the eraser writes \texttt{Z.t LeanArray.array} for an \texttt{Array} in the \texttt{.mli}, but the Makefile compiles a generated \texttt{.mli} without \texttt{-I \$(build)}, the directory that holds \texttt{LeanArray.cmi}, and without building \texttt{LeanArray.cmi} first; the build of a program with an \texttt{Array} in its signature stops there. The three explicit rules are needed because the pattern rule now depends on \texttt{LeanArray.cmi}, which would be a circular dependency for \texttt{LeanArray.cmi} itself; they keep the commands for these three runtime interfaces as they were.
\item \textbf{Behaviour before:} reproduction: a fresh build directory \texttt{B} holds \texttt{rArr.ast}, \texttt{rArr.ast.inlinings} and \texttt{rArr.mli} written by \texttt{\#erase rArr config \{remove\_\allowbreak{}irrel\_\allowbreak{}constr\_\allowbreak{}args := true\} to ... mli ...} for \texttt{rArr (n : Nat) : Array Nat := Array.mk [n, n + 1]}; \texttt{rArr.mli} is \texttt{val main: Z.t -\textgreater{} Z.t LeanArray.array}. In \texttt{benchmarks/\allowbreak{}via\_\allowbreak{}malfunction}, \texttt{make build=B FLAMBDA=0 MALFUNCTION\_\allowbreak{}NO\_\allowbreak{}FLAMBDA\_\allowbreak{}SWITCH=peregrine ARRAYML=JCFArrayOCaml4.ml PEREGRINE=\textless{}wrapper applying the R-1 rewrite\textgreater{} -o B/\allowbreak{}rArr.mli -o B/\allowbreak{}rArr.ast -o B/\allowbreak{}rArr.ast.inlinings B/\allowbreak{}rArr.cmi} runs \texttt{ocamlfind ocamlopt -package zarith -c B/\allowbreak{}rArr.mli}, which fails with \texttt{Error: Unbound module LeanArray} (make exits with status 2). The target \texttt{B/\allowbreak{}rArr.cmx} fails the same way after running peregrine. There, \texttt{rArr.cmi} comes before \texttt{axioms.cmx}, the only target that leads to \texttt{LeanArray.cmi}, so adding \texttt{-I} alone would not be enough in a fresh directory.
\item \textbf{Behaviour after:} the same \texttt{make ... B/\allowbreak{}rArr.cmi} copies \texttt{LeanArray.mli}, compiles \texttt{LeanArray.cmi}, then runs \texttt{ocamlfind ocamlopt -package zarith -I B -c B/\allowbreak{}rArr.mli}, and succeeds. \texttt{make ... B/\allowbreak{}rArr.cmx} succeeds too. After \texttt{make ... B/\allowbreak{}LeanArray.cmx}, the objects link with a harness written against \texttt{rArr.mli} (\texttt{LeanArray.def\_\allowbreak{}\_\allowbreak{}Array\_\allowbreak{}size (Obj.repr ()) (RArr.main (Z.of\_\allowbreak{}int 5))}), which prints \texttt{2}, the size of \texttt{rArr 5 = \#[5, 6]}. The benchmarks are unaffected: in a fresh build directory, \texttt{make -n} for the binary of \texttt{even} lists the same commands before and after, except \texttt{-I \$(build)} in the compilation of \texttt{even.mli}; a real build of \texttt{even} with the options above gives byte-identical \texttt{even.mli} and \texttt{even.mlf}, and the binary prints 1, 0 and 1 for the inputs 0, 7 and 1000, before and after.
\item \textbf{Effect on emitted .ast (corpus):} byte-identical for all 284 files, since the Makefile's erasure rules are unchanged; \texttt{scripts/\allowbreak{}corpus-diff.sh} against the corpus of S-3 reports 284 identical.
\item \textbf{Regression test:} \leanfile{tests/\allowbreak{}regress/\allowbreak{}makefile\_\allowbreak{}cmi.lean} erases \texttt{rArr} into a fresh build directory, asks \texttt{make -n} for the commands that build \texttt{rArr.cmi} there (with \texttt{OCAMLOPT=ocamlopt}, so it needs \texttt{make} but not OCaml), and writes them to \texttt{cmi\_\allowbreak{}commands.txt}. It fails before (\texttt{LeanArray.cmi is not built before rArr.cmi}: the only command is \texttt{ocamlopt -c \$(build)/\allowbreak{}rArr.mli}) and passes after, with the commands \texttt{cp LeanArray.mli \$(build)/\allowbreak{}LeanArray.mli}, \texttt{ocamlopt -c \$(build)/\allowbreak{}LeanArray.mli} and \texttt{ocamlopt -I \$(build) -c \$(build)/\allowbreak{}rArr.mli}.
\end{itemize}

\section{\texorpdfstring{S-5: \texttt{String} is no longer printed as \texttt{string} in the \texttt{.mli} signature}{S-5: String is no longer printed as string in the .mli signature}}\label{reg:S-5}

\begin{itemize}
\item \textbf{Commit:} the commit whose subject starts with \texttt{shipping(S-5):} (\texttt{git log -{}-grep='\textasciicircum{}shipping(S-5):'}).
\item \textbf{Files and functions:}

\begin{itemize}
\item \leanfile{LeanToLambdaBox/\allowbreak{}Erasure.lean}: \texttt{MLType} (constructor \texttt{string} removed), \texttt{MLType.toString} (case \texttt{string} removed), \texttt{to\_\allowbreak{}ml\_\allowbreak{}type} (case \texttt{String} removed). These are the \texttt{String} parts of \texttt{94a429c}; the code for \texttt{String} is again that of \texttt{main}.
\item \leanfile{tests/\allowbreak{}regress/\allowbreak{}mli\_\allowbreak{}types.lean}: docstring; new case \texttt{rStr}; its 3 expected outputs in \texttt{tests/\allowbreak{}regress/\allowbreak{}expected/\allowbreak{}mli\_\allowbreak{}types/}.
\item \texttt{doc/\allowbreak{}SHIPPING-CHANGES.md}: this entry; R-17 no longer lists \texttt{String} among the handled types; new R-27, a defect of the benchmark Makefile found while checking S-4.
\end{itemize}
\item \textbf{Why necessary:} defect D3 of S-2, reproduced on this branch. The \texttt{.mli} states the OCaml representation of the value, and since S-2 it gave \texttt{string} for a Lean \texttt{String}, without a warning. Nothing in the pipeline represents a Lean \texttt{String} as an OCaml \texttt{string}: a string literal is erased to \texttt{\ensuremath{\Box}} after a panic (R-4), and the constructor \texttt{String.mk} and every \texttt{String} operation are \texttt{@[extern]} constants, erased to axioms that the benchmark runtime does not implement (R-21). A program that builds or inspects a string therefore cannot be linked, and one that does neither links whatever type the \texttt{.mli} gives. The signature \texttt{string} made an unfounded claim and removed the warning that says the type is not supported.

Two fixes were possible. Implementing \texttt{String} would mean choosing an OCaml representation of Lean strings (UTF-8 bytes, or a list of code points as in \texttt{String.mk}), writing in the runtime \texttt{String.mk} and each \texttt{String}, \texttt{Char} and \texttt{UInt32} operation that a program uses, and erasing string literals (R-4). That is a new feature, which neither the verification goal nor any benchmark requires. The fix taken withdraws the claim: \texttt{String} goes back to the fallback path (warning and \texttt{unit}) that every type without a runtime representation takes (R-17).
\item \textbf{Behaviour before:} reproduction: \texttt{rStr n := String.mk (List.replicate n (Char.ofNat 97))}, \texttt{strLen (s : String) : Nat := s.length}, \texttt{strId (s : String) : String := s} and the literal \texttt{\textquotedbl{}abc\textquotedbl{}}, erased with \texttt{\{remove\_\allowbreak{}irrel\_\allowbreak{}constr\_\allowbreak{}args := true\}}, give the signatures \texttt{val main: Z.t -\textgreater{} string}, \texttt{val main: string -\textgreater{} Z.t}, \texttt{val main: string -\textgreater{} string} and \texttt{val main: string}, with no warning. \texttt{rStr.ast} has the axioms \texttt{String.mk}, \texttt{Char.ofNatAux}, \texttt{UInt32.ofBitVec}, \texttt{Nat.pow}, \texttt{Nat.decLt}, \texttt{Nat.beq} and \texttt{Nat.sub}, and \texttt{strLen.ast} has \texttt{String.length}. The literal becomes \texttt{(Untyped () (Some tBox))}, after \texttt{PANIC ... String literals not supported.} Compiled natively as for S-3, \texttt{rStr} fails in \texttt{malfunction cmx} with \texttt{Unbound value Axioms.def\_\allowbreak{}\_\allowbreak{}String\_\allowbreak{}mk}, and \texttt{strLen} with \texttt{Unbound value Axioms.def\_\allowbreak{}\_\allowbreak{}String\_\allowbreak{}length}. Only \texttt{strId}, which does not touch its argument, links: its harness passes \texttt{\textquotedbl{}abc\textquotedbl{}} and prints \texttt{abc}.
\item \textbf{Behaviour after:} the four signatures are \texttt{val main: Z.t -\textgreater{} unit}, \texttt{val main: unit -\textgreater{} Z.t}, \texttt{val main: unit -\textgreater{} unit} and \texttt{val main: unit}, byte-identical to those of \texttt{main} (\texttt{58701f8}), and each \texttt{String} in them is reported with \texttt{warning: failed to translate String into ML type, emitting unit instead.} The \texttt{.ast} and \texttt{.ast.inlinings} files are byte-identical before and after. Linking still fails for \texttt{rStr} and \texttt{strLen}, since the program itself needs the missing axioms (R-21).
\item \textbf{Effect on emitted .ast (corpus):} byte-identical for all 284 files; \texttt{scripts/\allowbreak{}corpus-diff.sh} against the corpus of S-4 reports 284 identical. No corpus \texttt{.mli} involves \texttt{String}. In the logs (\texttt{\_\allowbreak{}meta/\allowbreak{}logs}, not part of the corpus), the signature logged for \texttt{\#erase \textquotedbl{}abc\textquotedbl{}} in \leanfile{tests/\allowbreak{}corpus/\allowbreak{}Defects.lean} becomes \texttt{val main: unit}, with the warning above, as on \texttt{main}.
\item \textbf{Regression test:} \leanfile{tests/\allowbreak{}regress/\allowbreak{}mli\_\allowbreak{}types.lean}, case \texttt{rStr}, whose \texttt{\#guard\_\allowbreak{}msgs (warning)} requires the warning. It fails before (Lean reports that the \texttt{\#guard\_\allowbreak{}msgs} docstring does not match: no warning is produced) and passes after, with \texttt{rStr.mli} equal to \texttt{val main: Z.t -\textgreater{} unit}; the outputs of the other cases are unchanged.
\end{itemize}

\section{\texorpdfstring{S-6: The benchmark Makefile regenerates a missing \texttt{.ast.inlinings} before running peregrine}{S-6: The benchmark Makefile regenerates a missing .ast.inlinings before running peregrine}}\label{reg:S-6}

\begin{itemize}
\item \textbf{Commit:} the commit whose subject starts with \texttt{shipping(S-6):} (\texttt{git log -{}-grep='\textasciicircum{}shipping(S-6):'}).
\item \textbf{Files and functions:}

\begin{itemize}
\item \texttt{benchmarks/\allowbreak{}via\_\allowbreak{}malfunction/\allowbreak{}Makefile}: the \texttt{INLINING=1} pattern rule \texttt{\$(build)/\allowbreak{}\%.mlf} gains the prerequisite \texttt{\$(build)/\allowbreak{}\%.ast.inlinings}, and a comment line above the \texttt{ifeq} says why; its command is unchanged. The \texttt{INLINING=0} rule is unchanged.
\item new \leanfile{tests/\allowbreak{}regress/\allowbreak{}makefile\_\allowbreak{}inlinings.lean} and \texttt{tests/\allowbreak{}regress/\allowbreak{}expected/\allowbreak{}makefile\_\allowbreak{}inlinings/\allowbreak{}mlf\_\allowbreak{}commands.txt}.
\item \texttt{doc/\allowbreak{}SHIPPING-CHANGES.md}: this entry; new R-28, found while running the new test with GNU make 4.3.
\end{itemize}
\item \textbf{Why necessary:} defect D4 of S-2 (claim C4), reproduced on this branch. With \texttt{INLINING=1} the rule for \texttt{\%.mlf} runs \texttt{peregrine compile \%.ast \textless{}config\textgreater{} -{}-attributes=\%.ast.inlinings}, but its only prerequisite is \texttt{\%.ast}. Since \texttt{94a429c}, \texttt{\%.ast.inlinings} is a target of the rule that runs \texttt{lake lean}, so make knows how to rebuild it, but make rebuilds a file only when a target depends on it. When the \texttt{.ast} is present and up to date and the \texttt{.ast.inlinings} is missing (for instance in a build directory filled by a frontend that did not write the file yet, the situation of the Zulip report of Feb 17), make runs peregrine on a missing file. The target added by \texttt{94a429c} is what makes the new prerequisite possible: without it, as at \texttt{a104486}, the prerequisite gives \texttt{No rule to make target} (Zulip, Feb 13).
\item \textbf{Behaviour before:} reproduction: in \texttt{benchmarks/\allowbreak{}via\_\allowbreak{}malfunction}, with the options \texttt{FLAMBDA=0 MALFUNCTION\_\allowbreak{}NO\_\allowbreak{}FLAMBDA\_\allowbreak{}SWITCH=peregrine ARRAYML=JCFArrayOCaml4.ml PEREGRINE=\textless{}wrapper applying the R-1 rewrite\textgreater{}} (build directory \texttt{build/\allowbreak{}7ee730d2}), \texttt{make build/\allowbreak{}7ee730d2/\allowbreak{}even.mlf} succeeds from scratch. After deleting \texttt{even.ast.inlinings} and \texttt{even.mlf}, the same command runs peregrine alone, which fails with \texttt{peregrine: option '-{}-attributes': invalid element in list (build/\allowbreak{}7ee730d2/\allowbreak{}even.ast.inlinings): no build/\allowbreak{}7ee730d2/\allowbreak{}even.ast.inlinings file or directory}, and make exits with status 2: the error reported on Zulip on Feb 17. Every later run fails the same way, and so does \texttt{make bin/\allowbreak{}even}, until the file is restored by hand.
\item \textbf{Behaviour after:} the same command first runs \texttt{lake lean build/\allowbreak{}7ee730d2/\allowbreak{}even.lean}, which rewrites \texttt{even.ast}, \texttt{even.ast.inlinings} and \texttt{even.mli}, then runs peregrine, and succeeds; the next run prints \texttt{make: 'build/\allowbreak{}7ee730d2/\allowbreak{}even.mlf' is up to date.} With \texttt{even.ast.inlinings} deleted, \texttt{make bin/\allowbreak{}even} erases, compiles and links, and the binary prints 1, 0 and 1 for the inputs 0, 7 and 1000. The regenerated \texttt{even.ast} and \texttt{even.mli} are byte-identical to those of a build from scratch. In a fresh build directory, \texttt{make -n bin/\allowbreak{}even} lists the same commands before and after, with \texttt{INLINING=1} and with \texttt{INLINING=0}.
\item \textbf{Effect on emitted .ast (corpus):} byte-identical for all 284 files, since the erasure rule is unchanged; \texttt{scripts/\allowbreak{}corpus-diff.sh} against the corpus of S-5 reports 284 identical.
\item \textbf{Regression test:} \leanfile{tests/\allowbreak{}regress/\allowbreak{}makefile\_\allowbreak{}inlinings.lean} erases \texttt{rInl} into a fresh build directory and asks \texttt{make -n} for the commands that build \texttt{rInl.mlf} there, with \texttt{PEREGRINE=peregrine} and the generated \texttt{rInl.lean} taken as old (\texttt{-o}), so it needs \texttt{make} but neither peregrine nor OCaml. With every output present, the only command is peregrine's. With \texttt{rInl.ast.inlinings} deleted, \texttt{lake lean \$(build)/\allowbreak{}rInl.lean} comes before it. With the file deleted and \texttt{INLINING=0}, the only command is peregrine's, without \texttt{-{}-attributes}. The three lists are written to \texttt{mlf\_\allowbreak{}commands.txt}. Before, the second list is the peregrine command alone and the test fails with \texttt{rInl.ast.inlinings missing: expected the lake lean command before the peregrine command}; after, it passes with GNU make 4.4.1 and with make 4.3.
\end{itemize}

\section{\texorpdfstring{S-7: Auto-inlining marks a constant only if its inlined body is small}{S-7: Auto-inlining marks a constant only if its inlined body is small}}\label{reg:S-7}

\begin{itemize}
\item \textbf{Commit:} the commit whose subject starts with \texttt{shipping(S-7):} (\texttt{git log -{}-grep='\textasciicircum{}shipping(S-7):'}).
\item \textbf{Files and functions:}

\begin{itemize}
\item \leanfile{LeanToLambdaBox/\allowbreak{}Basic.lean}: \texttt{ModPath} and \texttt{Kername} also derive \texttt{BEq} and \texttt{Hashable}, so that kernames can key a hash map.
\item \leanfile{LeanToLambdaBox/\allowbreak{}Erasure.lean}: \texttt{ErasureState} (new field \texttt{inlinedSizes}); \texttt{ErasureConfig} (docstring of \texttt{auto\_\allowbreak{}inline\_\allowbreak{}typeclass\_\allowbreak{}dispatch}); new \texttt{autoInlineMaxSize} (40) and \texttt{LBTerm.inlinedSize}; \texttt{erase.visitMutual}: in the non-recursive branch, the inlined size of the erased body is computed and recorded for an \texttt{@[inline]} constant, and a constant that the option would mark is marked only if its inlined size is at most \texttt{autoInlineMaxSize}, with a log message that gives the size or the refusal; in the recursive branch, the inlined size of the fixpoint of an \texttt{@[inline]} constant is recorded.
\item new \leanfile{tests/\allowbreak{}regress/\allowbreak{}auto\_\allowbreak{}inline\_\allowbreak{}size.lean}, its expected outputs \texttt{tests/\allowbreak{}regress/\allowbreak{}expected/\allowbreak{}auto\_\allowbreak{}inline\_\allowbreak{}size/} (8 files) and \texttt{tests/\allowbreak{}regress/\allowbreak{}expected-peregrine/\allowbreak{}auto\_\allowbreak{}inline\_\allowbreak{}size/} (4 files).
\item \texttt{doc/\allowbreak{}SHIPPING-CHANGES.md}: this entry.
\end{itemize}
\item \textbf{Why necessary:} defect D5 of S-2, reproduced on this branch. With the option on, every instance and every trivial alias is marked, whatever its size. Peregrine's inlining pass (MetaRocq \texttt{EInlining.inline\_\allowbreak{}env}) replaces each use of a marked constant by its body, into which the marked constants were already inlined, so the code grows without bound. \texttt{e46224f} limited the erased body of an instance to 40 nodes; \texttt{8db2dd9} removed the limit. A limit on the erased body alone does not bound the growth either: in a chain of instances where each one calls the method of the previous one twice, every body has at most 12 nodes, and the inlined code doubles at each level (see below). The limit is therefore put on what Peregrine substitutes: the body with the marked constants, including \texttt{@[inline]} ones, inlined into it (\texttt{LBTerm.inlinedSize}). The bound 40 is the value of \texttt{e46224f}. On the 20 benchmarks, every constant still marked measures at most 38 and every refused one at least 140, so any bound from 38 to 139 marks the same constants.
\item \textbf{Behaviour before:} reproduction: the 20 \texttt{natio} benchmarks are erased with \texttt{config \{remove\_\allowbreak{}irrel\_\allowbreak{}constr\_\allowbreak{}args := true, auto\_\allowbreak{}inline\_\allowbreak{}typeclass\_\allowbreak{}dispatch := true\}} (source file elaborated with \texttt{lake env lean} in \texttt{benchmarks/\allowbreak{}via\_\allowbreak{}malfunction}), and each \texttt{.ast} is compiled with \texttt{peregrine compile \textless{}t\textgreater{}.ast unbox.config -{}-attributes=\textless{}t\textgreater{}.ast.inlinings} after the R-1 rewrite.

\begin{itemize}
\item unionfind: 37 constants are marked, among them \texttt{UnionFind.StateT'.instMonad} and \texttt{UnionFind.ExceptT'.instMonad}. The \texttt{.mlf} has 14 510 809 bytes (126 374 with the option off) and peregrine takes 3.34 s (0.05 s). The benchmark Makefile (\texttt{FLAMBDA=0}, switch \texttt{peregrine}, OCaml 4.14.2, \texttt{ARRAYML=JCFArrayOCaml4.ml}, the given \texttt{.ast} files) builds the binary in 15.2 s (0.9 s), and the binary has 36 297 424 bytes (2 972 176). unionfind\_noinline: 5 818 655 bytes (89 106).
\item \leanfile{tests/\allowbreak{}regress/\allowbreak{}auto\_\allowbreak{}inline\_\allowbreak{}size.lean} (option on): \texttt{instMonadSt}, a \texttt{Monad} instance, and the four instances \texttt{i0} to \texttt{i3} of the chain are marked; \texttt{chain.mlf} has 3 208 bytes and \texttt{twice.mlf} 20 223. With the chain extended to eleven levels (\texttt{j0} to \texttt{j10}, each calling the previous method twice), all eleven are marked, and the \texttt{.mlf} of a program calling level 3, 6 or 10 has 3 208, 26 332 or 425 747 bytes.
\end{itemize}
\item \textbf{Behaviour after:} same reproduction.

\begin{itemize}
\item unionfind: 35 constants are marked; the two monad instances are refused (inlined sizes 881 and 1 098). The \texttt{.mlf} has 129 332 bytes, peregrine takes 0.05 s, the Makefile builds the binary in 0.8 s, and the binary has 2 900 376 bytes. unionfind\_noinline: 88 726 bytes, without its two monad instances. In qsort, qsort\_fin and qsort\_single, \texttt{Array.instGetElem?NatLtSize} (140) and \texttt{Vector.instGetElemNatLt} (146) are no longer marked (qsort \texttt{.mlf}: 34 933 to 34 583 bytes). The other 15 benchmarks mark the same constants as before, and their \texttt{.mlf} files are identical.
\item The binaries print the same results. Running times, minimum of 5 runs, with the option off / on before / on after: binarytrees 17: 0.75 / 0.38 / 0.39 s; qsort 1000: 0.51 / 0.36 / 0.37 s; triangle\_rec 10000000 (with \texttt{ulimit -s unlimited}): 7.51 / 0.109 / 0.107 s; unionfind 100000: 0.75 / 0.37 / 0.72 s. The speedup of unionfind came from inlining its two monad instances, at the price of the code growth above; it is lost. The other speedups are kept.
\item \leanfile{tests/\allowbreak{}regress/\allowbreak{}auto\_\allowbreak{}inline\_\allowbreak{}size.lean}: \texttt{instMonadSt} is refused (inlined size 453); \texttt{i0}, \texttt{i1} and \texttt{i3} are marked (8, 30 and 16) and \texttt{i2} is refused (74), after which \texttt{i3} refers to \texttt{i2} without inlining it. \texttt{chain.mlf} has 1 416 bytes and \texttt{twice.mlf} 7 892. In the eleven-level chain, \texttt{j0} and the odd levels are marked, and the three programs give 1 416, 2 292 and 3 468 bytes.
\item \texttt{peregrine eval} of \texttt{chain 2} and \texttt{(twice 3).1} under Peano naturals, with their inlinings, gives 10 and 7, before and after.
\item Each decision is logged, for example \texttt{Auto-inlining typeclass instance i1 (inlined size 30).} and \texttt{Not auto-inlining typeclass instance i2: inlined size 74 exceeds 40.}
\item The bound concerns Peregrine's inlining pass. Its beta-reduction pass (MetaRocq \texttt{EBeta.betared}) then substitutes the arguments of an inlined \ensuremath{\lambda} into its body, which copies an argument once per occurrence of the bound variable.
\end{itemize}
\item \textbf{Effect on emitted .ast (corpus):} byte-identical for all 284 files; \texttt{scripts/\allowbreak{}corpus-diff.sh} against the corpus of S-6 reports 284 identical. The corpus does not turn the option on, and the \texttt{.ast.inlinings} of other configurations lists only \texttt{@[inline]} constants, which this change does not affect. With the option on (outside the corpus), the 20 benchmark \texttt{.ast} files are byte-identical before and after, and only the \texttt{.ast.inlinings} files of the five benchmarks named above change, each losing the constants named above.
\item \textbf{Regression test:} \leanfile{tests/\allowbreak{}regress/\allowbreak{}auto\_\allowbreak{}inline\_\allowbreak{}size.lean} erases \texttt{twice} and \texttt{chain} with the option on, and \texttt{(twice 3).1} and \texttt{chain 2} under Peano naturals. It fails before (\texttt{chain.ast.inlinings} and \texttt{chain2.ast.inlinings} list \texttt{i2}, and \texttt{twice.ast.inlinings} and \texttt{twice3.ast.inlinings} list \texttt{instMonadSt}) and passes after. With \texttt{PEREGRINE} set, \texttt{chain2.ast} and \texttt{twice3.ast} validate and evaluate to 10 and 7.
\end{itemize}

\section{\texorpdfstring{S-8: Auto-inlining marks only constants whose erased body is a value}{S-8: Auto-inlining marks only constants whose erased body is a value}}\label{reg:S-8}

\begin{itemize}
\item \textbf{Commit:} the commit whose subject starts with \texttt{shipping(S-8):} (\texttt{git log -{}-grep='\textasciicircum{}shipping(S-8):'}).
\item \textbf{Files and functions:}

\begin{itemize}
\item \leanfile{LeanToLambdaBox/\allowbreak{}Erasure.lean}: new \texttt{LBTerm.isValue} (with its auxiliary \texttt{LBTerm.isValue.isConstructorApp}); \texttt{ErasureConfig} (docstring of \texttt{auto\_\allowbreak{}inline\_\allowbreak{}typeclass\_\allowbreak{}dispatch}); \texttt{erase.visitMutual}: in the non-recursive branch, a constant that the option would mark is not marked if its erased body is not a value, with a log message.
\item new \leanfile{tests/\allowbreak{}regress/\allowbreak{}auto\_\allowbreak{}inline\_\allowbreak{}values.lean}, its expected outputs \texttt{tests/\allowbreak{}regress/\allowbreak{}expected/\allowbreak{}auto\_\allowbreak{}inline\_\allowbreak{}values/} (4 files) and \texttt{tests/\allowbreak{}regress/\allowbreak{}expected-peregrine/\allowbreak{}auto\_\allowbreak{}inline\_\allowbreak{}values/} (2 files).
\item \texttt{doc/\allowbreak{}SHIPPING-CHANGES.md}: this entry.
\end{itemize}
\item \textbf{Why necessary:} defect D6 of S-2, reproduced on this branch. The option marks an instance whatever its body. The compiled program evaluates a top-level constant once; once the constant is inlined, its body is evaluated at every use, so a body that computes is computed again each time. The bound of S-7 does not prevent this, since a small body can start a long computation. The example of S-2, \texttt{instTbl} (\texttt{let s := slowSum 100000; \ensuremath{\langle}fun i =\textgreater{} s + i\ensuremath{\rangle}}), is refused by that bound (inlined size 75), but smaller bodies of the same kind are marked. The fix marks a constant only if its erased body is a value (\texttt{LBTerm.isValue}): a \ensuremath{\lambda}, \ensuremath{\Box}, a primitive, a constant, or a constructor applied to such values. Evaluating such a body at each use repeats no computation other than building it. A constant counts as a value because it refers to a top-level definition, which the compiled program evaluates once.
\item \textbf{Behaviour before:} reproduction: \texttt{slowSum : Nat \ensuremath{\to} Nat} adds \texttt{n + (n-1) + \ldots{} + 1} by recursion, \texttt{instSlow : Inhabited Nat := \ensuremath{\langle}slowSum 100000\ensuremath{\rangle}} with \texttt{useSlow n := (List.range n).foldl (fun acc \_\allowbreak{} =\textgreater{} acc + (default : Nat)) 0}, and \texttt{instLet : Tbl := let s := slowSum 100000; \ensuremath{\langle}Nat.add s\ensuremath{\rangle}} (with \texttt{class Tbl where get : Nat \ensuremath{\to} Nat}) with \texttt{useLet n := (List.range n).foldl (fun acc i =\textgreater{} acc + Tbl.get (self := instLet) i) 0}. Both programs are erased with \texttt{config \{remove\_\allowbreak{}irrel\_\allowbreak{}constr\_\allowbreak{}args := true, auto\_\allowbreak{}inline\_\allowbreak{}typeclass\_\allowbreak{}dispatch := true\}} and built natively as for S-7 (benchmark Makefile, \texttt{FLAMBDA=0}, OCaml 4.14.2, peregrine after the R-1 rewrite).

\begin{itemize}
\item \texttt{instSlow} and \texttt{instLet} are marked (inlined sizes 26 and 28).
\item \texttt{useSlow 1000} takes 0.64 s and \texttt{useSlow 3000} 1.85 s; \texttt{useLet 1000} takes 0.68 s and \texttt{useLet 3000} 2.05 s. \texttt{slowSum 100000} is computed once per list element.
\item Benchmarks with the option on (as for S-7): qsort, qsort\_fin and qsort\_single mark \texttt{instMinNat} and \texttt{Nat.instMax} (inlined size 25), whose bodies are the applications \texttt{minOfLe \ldots{}} and \texttt{maxOfLe \ldots{}} that build a dictionary.
\end{itemize}
\item \textbf{Behaviour after:} same reproduction.

\begin{itemize}
\item \texttt{instSlow} and \texttt{instLet} are not marked, with the log messages \texttt{Not auto-inlining typeclass instance instSlow: its body is not a value.} and the same for \texttt{instLet}.
\item \texttt{useSlow 1000} and \texttt{useSlow 3000} take 0.006 s each, and \texttt{useLet} 0.004 s and 0.005 s. The printed results are unchanged: 5000050000000 and 15000150000000 for \texttt{useSlow}, 5000050499500 and 15000154498500 for \texttt{useLet}.
\item Benchmarks: \texttt{instMinNat} and \texttt{Nat.instMax} are no longer marked, in the three qsort benchmarks only (47 to 45, 52 to 50 and 46 to 44 marked constants; qsort \texttt{.mlf} 34 583 to 34 009 bytes). Running times, minimum of 5 runs, before / after: qsort 1000: 0.371 / 0.369 s; qsort\_fin 1000: 0.382 / 0.385 s; qsort\_single 100000: 0.179 / 0.177 s; the results are unchanged. The other 17 benchmarks mark the same constants, and their \texttt{.mlf} files are identical.
\end{itemize}
\item \textbf{Effect on emitted .ast (corpus):} byte-identical for all 284 files; \texttt{scripts/\allowbreak{}corpus-diff.sh} against the corpus of S-7 reports 284 identical. The corpus does not turn the option on. With the option on (outside the corpus), the 20 benchmark \texttt{.ast} files are byte-identical before and after, and only the \texttt{.ast.inlinings} files of qsort, qsort\_fin and qsort\_single change, each losing \texttt{instMinNat} and \texttt{Nat.instMax}.
\item \textbf{Regression test:} \leanfile{tests/\allowbreak{}regress/\allowbreak{}auto\_\allowbreak{}inline\_\allowbreak{}values.lean} erases \texttt{useAll}, which uses \texttt{instSlow}, \texttt{instLet} and \texttt{instLam} (\texttt{\ensuremath{\langle}fun i =\textgreater{} slowSum i\ensuremath{\rangle}}, a value), with the option on, and \texttt{useAll 2} under Peano naturals. It fails before (\texttt{useAll.ast.inlinings} lists \texttt{instSlow} and \texttt{instLet}) and passes after, where only \texttt{instLam} of the three is listed. With \texttt{PEREGRINE} set, \texttt{useAll2.ast} validates and evaluates to 25.
\end{itemize}

\section{\texorpdfstring{S-9: The dead \texttt{containsFix} guard of auto-inlining is removed}{S-9: The dead containsFix guard of auto-inlining is removed}}\label{reg:S-9}

\begin{itemize}
\item \textbf{Commit:} the commit whose subject starts with \texttt{shipping(S-9):} (\texttt{git log -{}-grep='\textasciicircum{}shipping(S-9):'}).
\item \textbf{Files and functions:}

\begin{itemize}
\item \leanfile{LeanToLambdaBox/\allowbreak{}Erasure.lean}: \texttt{LBTerm.containsFix} removed; \texttt{erase.visitMutual}: the condition \texttt{!t.containsFix} of the auto-inline step and its comment removed, and the comment now says that recursive definitions are never marked; \texttt{ErasureConfig}: the paragraph of the docstring of \texttt{auto\_\allowbreak{}inline\_\allowbreak{}typeclass\_\allowbreak{}dispatch} about \texttt{LBTerm.fix} is replaced by one that says that only non-recursive definitions are considered.
\item new \leanfile{tests/\allowbreak{}regress/\allowbreak{}auto\_\allowbreak{}inline\_\allowbreak{}recursive.lean}, its expected outputs \texttt{tests/\allowbreak{}regress/\allowbreak{}expected/\allowbreak{}auto\_\allowbreak{}inline\_\allowbreak{}recursive/} (4 files) and \texttt{tests/\allowbreak{}regress/\allowbreak{}expected-peregrine/\allowbreak{}auto\_\allowbreak{}inline\_\allowbreak{}recursive/} (2 files).
\item \texttt{doc/\allowbreak{}SHIPPING-CHANGES.md}: this entry.
\end{itemize}
\item \textbf{Why necessary:} defect D7 of S-2, confirmed on this branch. \texttt{8db2dd9} added the guard \texttt{!t.containsFix}, documented as refusing to inline recursive bodies, in the non-recursive branch of \texttt{erase.visitMutual}. The term \texttt{t} there is the result of \texttt{visitExpr}, and \texttt{LBTerm.fix} is built only in the recursive branch of \texttt{erase.visitMutual} (the declaration \texttt{.fix defs i}), so \texttt{t} never contains one and the guard never applies. Recursive definitions were never candidates, since the auto-inline step exists only in the non-recursive branch. The guard and its documentation described a protection that the code gets elsewhere, and they suggested that a body calling a recursive function is refused, which it is not (it contains a constant, not a fixpoint). The change removes the dead code and states where the protection comes from.
\item \textbf{Behaviour before:} reproduction: \leanfile{tests/\allowbreak{}regress/\allowbreak{}auto\_\allowbreak{}inline\_\allowbreak{}recursive.lean} with the option on. \texttt{depthInst : (n : Nat) \ensuremath{\to} Depth n}, recursive and declared an instance with \texttt{attribute [instance]}, is erased to a fixpoint and is not marked; \texttt{instSum : Tbl := \ensuremath{\langle}fun i =\textgreater{} sumTo i\ensuremath{\rangle}}, which calls the recursive \texttt{sumTo}, is marked (inlined size 6). A copy of the eraser at the commit of S-8 that logs every non-recursive body for which \texttt{containsFix} holds logs nothing on the corpus examples, the nine regression tests and the 20 benchmarks erased with the option on (164 \texttt{.ast} files, 88 of them with a fixpoint).
\item \textbf{Behaviour after:} the same: outputs and log messages of the test above are byte-identical, and so are the \texttt{.ast}, \texttt{.ast.inlinings} and \texttt{.mlf} files of the 20 benchmarks erased with the option on. With \texttt{PEREGRINE} set, \texttt{useBoth 2} evaluates to 6 under Peano naturals, before and after.
\item \textbf{Effect on emitted .ast (corpus):} byte-identical for all 284 files; \texttt{scripts/\allowbreak{}corpus-diff.sh} against the corpus of S-8 reports 284 identical.
\item \textbf{Regression test:} the change does not alter behaviour, so no test fails before it. \leanfile{tests/\allowbreak{}regress/\allowbreak{}auto\_\allowbreak{}inline\_\allowbreak{}recursive.lean} guards what the removed guard claimed and what holds without it: with the option on, a recursive instance (\texttt{depthInst}) is not marked, and a non-recursive instance that calls a recursive function (\texttt{instSum}) is marked. It passes before and after.
\end{itemize}

\section{\texorpdfstring{S-10: Auto-inlining does not mark constants tagged \texttt{@[noinline]}}{S-10: Auto-inlining does not mark constants tagged @[noinline]}}\label{reg:S-10}

\begin{itemize}
\item \textbf{Commit:} the commit whose subject starts with \texttt{shipping(S-10):} (\texttt{git log -{}-grep='\textasciicircum{}shipping(S-10):'}).
\item \textbf{Files and functions:}

\begin{itemize}
\item \leanfile{LeanToLambdaBox/\allowbreak{}Erasure.lean}: \texttt{erase.visitMutual} (new local \texttt{leanNoinline}, read from \texttt{Compiler.getInlineAttribute?}; in the non-recursive branch, a constant that the option would mark is not marked if it is tagged \texttt{@[noinline]}, with a log message); \texttt{ErasureConfig} (docstring of \texttt{auto\_\allowbreak{}inline\_\allowbreak{}typeclass\_\allowbreak{}dispatch}).
\item new \leanfile{tests/\allowbreak{}regress/\allowbreak{}auto\_\allowbreak{}inline\_\allowbreak{}noinline.lean}, its expected outputs \texttt{tests/\allowbreak{}regress/\allowbreak{}expected/\allowbreak{}auto\_\allowbreak{}inline\_\allowbreak{}noinline/} (4 files) and \texttt{tests/\allowbreak{}regress/\allowbreak{}expected-peregrine/\allowbreak{}auto\_\allowbreak{}inline\_\allowbreak{}noinline/} (2 files).
\item \texttt{doc/\allowbreak{}SHIPPING-CHANGES.md}: this entry.
\end{itemize}
\item \textbf{Why necessary:} defect D8 of S-2, reproduced on this branch. \texttt{@[noinline]} is the user's request that a definition not be inlined. \texttt{erase.visitMutual} consulted only \texttt{@[inline]} and \texttt{@[always\_\allowbreak{}inline]}, so the option marked a \texttt{@[noinline]} instance or alias like any other, and Peregrine inlined it. The example of S-2, \texttt{@[noinline] instance instBar : Inhabited Nat := \ensuremath{\langle}42\ensuremath{\rangle}}, is no longer marked since S-8, because the literal \texttt{42} is erased to an application of \texttt{OfNat.ofNat}, which is not a value; the reproduction below uses bodies that are values.
\item \textbf{Behaviour before:} reproduction: \leanfile{tests/\allowbreak{}regress/\allowbreak{}auto\_\allowbreak{}inline\_\allowbreak{}noinline.lean}, which declares \texttt{@[noinline] instance instNo : OpNo := \ensuremath{\langle}fun n =\textgreater{} n.succ\ensuremath{\rangle}} and \texttt{@[noinline] def addNo : Nat \ensuremath{\to} Nat \ensuremath{\to} Nat := Nat.add}, and the same without the attribute (\texttt{instYes}, \texttt{addYes}), and erases \texttt{useAll n := addNo (OpNo.op n) (addYes (OpYes.op n) n)} with the option on. \texttt{useAll.ast.inlinings} lists \texttt{instNo} and \texttt{addNo}, and the log says \texttt{Auto-inlining typeclass instance instNo (inlined size 6).} and \texttt{Auto-inlining trivial alias addNo (inlined size 1).}
\item \textbf{Behaviour after:} \texttt{instNo} and \texttt{addNo} are not listed, and the log says \texttt{Not auto-inlining typeclass instance instNo: it is tagged @[noinline].} and the same for \texttt{addNo}; \texttt{instYes} and \texttt{addYes} are still listed. The \texttt{.ast} files are byte-identical. With \texttt{PEREGRINE} set, \texttt{useAll 2} evaluates to 8 under Peano naturals, before and after. The 20 benchmarks erased with the option on give byte-identical \texttt{.ast} and \texttt{.ast.inlinings} files: none of their candidates is tagged \texttt{@[noinline]}.
\item \textbf{Effect on emitted .ast (corpus):} byte-identical for all 284 files; \texttt{scripts/\allowbreak{}corpus-diff.sh} against the corpus of S-9 reports 284 identical. The corpus does not turn the option on.
\item \textbf{Regression test:} \leanfile{tests/\allowbreak{}regress/\allowbreak{}auto\_\allowbreak{}inline\_\allowbreak{}noinline.lean} fails before (\texttt{useAll.ast.inlinings} and \texttt{useAll2.ast.inlinings} list \texttt{instNo} and \texttt{addNo}) and passes after.
\end{itemize}

\section{\texorpdfstring{S-11: The documentation of auto-inlining describes the shapes the code accepts}{S-11: The documentation of auto-inlining describes the shapes the code accepts}}\label{reg:S-11}

\begin{itemize}
\item \textbf{Commit:} the commit whose subject starts with \texttt{shipping(S-11):} (\texttt{git log -{}-grep='\textasciicircum{}shipping(S-11):'}).
\item \textbf{Files and functions:}

\begin{itemize}
\item \leanfile{LeanToLambdaBox/\allowbreak{}Erasure.lean}: docstrings of \texttt{ErasureConfig.auto\_\allowbreak{}inline\_\allowbreak{}typeclass\_\allowbreak{}dispatch} (the list of candidates), \texttt{LBTerm.isTrivialAlias} and \texttt{LBTerm.stripLambdas}. No code changes.
\item new \leanfile{tests/\allowbreak{}regress/\allowbreak{}auto\_\allowbreak{}inline\_\allowbreak{}shapes.lean}, its expected outputs \texttt{tests/\allowbreak{}regress/\allowbreak{}expected/\allowbreak{}auto\_\allowbreak{}inline\_\allowbreak{}shapes/} (4 files) and \texttt{tests/\allowbreak{}regress/\allowbreak{}expected-peregrine/\allowbreak{}auto\_\allowbreak{}inline\_\allowbreak{}shapes/} (2 files).
\item \texttt{doc/\allowbreak{}SHIPPING-CHANGES.md}: this entry.
\end{itemize}
\item \textbf{Why necessary:} defect D9 of S-2, confirmed on this branch. The docstring of the option said that a non-instance is a candidate when its erased body is "a bare \texttt{const}/\texttt{proj}, or a single-ctor structure literal whose fields are shallow", and that of \texttt{LBTerm.isTrivialAlias} named "a single-ctor structure literal (the usual \texttt{Foo.mk arg\ensuremath{_1} \ldots{} arg\ensuremath{_n}} shape \ldots{})". The code tests, after the leading \ensuremath{\lambda}s, for \texttt{.const}, \texttt{.proj} or \texttt{.construct \_\allowbreak{} 0 \_\allowbreak{}}. The eraser emits a constructor as a \texttt{construct} node with an empty argument list, applied to its parameters and fields with \texttt{LBTerm.app} (\texttt{erase.visitConstructor}), so the third case matches only a constructor of index 0 without parameters or fields: never a structure literal, whatever its fields, and \texttt{false} but not \texttt{true}. The first case also matches constant functions, which the documentation called aliases. The docstring of \texttt{LBTerm.stripLambdas} said that the stripped \ensuremath{\lambda}s are typeclass-instance parameters; they are any parameters, and in a projection function they include the structure argument. The docstrings now state what the code accepts. The code is kept: a definition of any of these shapes is marked only after the checks of S-7 to S-10 (at most 40 nodes once inlined, a value, not recursive, not \texttt{@[noinline]}), and matching structure literals would be a new feature.

The commit messages of \texttt{94a429c} and \texttt{8db2dd9}, which D9 also concerns, cannot be changed; S-2 records what they claim.
\item \textbf{Behaviour before:} reproduction: \leanfile{tests/\allowbreak{}regress/\allowbreak{}auto\_\allowbreak{}inline\_\allowbreak{}shapes.lean} with the option on. Marked: \texttt{addAlias : Nat \ensuremath{\to} Nat \ensuremath{\to} Nat := Nat.add} (a constant), \texttt{kfun (\_\allowbreak{} : Nat) : Nat := seven} (a constant function), the projection function \texttt{P.b}, and \texttt{falseDef : Bool := false}. Not marked: \texttt{trueDef : Bool := true}, \texttt{noneDef : Option Nat := none} (a constructor applied to its type parameter) and the structure literal \texttt{pVal : P := \ensuremath{\langle}Nat.zero, Nat.succ Nat.zero\ensuremath{\rangle}}. This contradicts the old documentation for \texttt{pVal} (a single-constructor structure literal with shallow fields) and for \texttt{kfun} (not an alias).
\item \textbf{Behaviour after:} the same outputs and log messages, byte for byte; the documentation now describes them. The 20 benchmarks erased with the option on give byte-identical \texttt{.ast} and \texttt{.ast.inlinings} files. With \texttt{PEREGRINE} set, \texttt{useAll 2} evaluates to 10 under Peano naturals.
\item \textbf{Effect on emitted .ast (corpus):} byte-identical for all 284 files; \texttt{scripts/\allowbreak{}corpus-diff.sh} against the corpus of S-10 reports 284 identical.
\item \textbf{Regression test:} the change does not alter behaviour, so no test fails before it. \leanfile{tests/\allowbreak{}regress/\allowbreak{}auto\_\allowbreak{}inline\_\allowbreak{}shapes.lean} guards the documented shapes: it pins the four marked and the three unmarked definitions above. It passes before and after.
\end{itemize}

\section{\texorpdfstring{S-12: Port to Lean v4.33.0-rc2, with a Lake dependency on \texttt{lean4lean}}{S-12: Port to Lean v4.33.0-rc2, with a Lake dependency on lean4lean}}\label{reg:S-12}

\begin{itemize}
\item \textbf{Commit:} the commit whose subject starts with \texttt{shipping(S-12):} (\texttt{git log -{}-grep='\textasciicircum{}shipping(S-12):'}).
\item \textbf{Files and functions:}

\begin{itemize}
\item \texttt{lean-toolchain}: \texttt{leanprover/\allowbreak{}lean4:v4.22.0} becomes \texttt{leanprover/\allowbreak{}lean4:v4.33.0-rc2}.
\item \texttt{lakefile.toml}: new \texttt{[[require]]} of \texttt{lean4lean}, git \texttt{https:/\allowbreak{}/\allowbreak{}github.com/\allowbreak{}barabbs/\allowbreak{}lean4lean}, rev \texttt{8223d223ed98661882e95d9d6a7126df7097cd76} (the fork's \texttt{master}, whose toolchain is \texttt{v4.33.0-rc2}). \texttt{defaultTargets} is unchanged (\texttt{LeanToLambdaBox}).
\item \texttt{lake-manifest.json}: regenerated by \texttt{lake update}. The format goes from 1.1.0 to 1.2.0 (new key \texttt{fixedToolchain: false}), and it lists \texttt{lean4lean} at \texttt{8223d223\ldots{}} and, inherited from it, \texttt{batteries} at \texttt{76e1c118b0700b4ceafe99532e887d6431625e1a} (input rev \texttt{v4.33.0-rc2}).
\item \leanfile{LeanToLambdaBox/\allowbreak{}Erasure.lean}:

\begin{itemize}
\item \texttt{register\_\allowbreak{}inductive}: the type ascription \texttt{List projection\_\allowbreak{}body} becomes \texttt{List ProjectionBody} (E1);
\item \texttt{fvar\_\allowbreak{}to\_\allowbreak{}name}: the predicate passed to \texttt{String.all} becomes \texttt{fun (c : Char) =\textgreater{} decide (33 \textless{}= c.toNat /\allowbreak{}\textbackslash{} c.toNat \textless{} 127)} (E2);
\item new \texttt{csimpReplaceConstants}; \texttt{prepare\_\allowbreak{}erasure} calls it in place of \texttt{Compiler.CSimp.replaceConstants} (E3);
\item \texttt{erase.visitCases}: \texttt{casesInfo.altsRange.start} becomes \texttt{casesInfo.altsRange.lower} (four places, in the machine-\texttt{Nat} and machine-\texttt{Int} paths); in the generic path, the loop reads each entry of \texttt{casesInfo.altNumParams} as \texttt{.ctor \_\allowbreak{} numFields} and throws an error on a \texttt{.default} entry (E4).
\end{itemize}
\item \texttt{benchmarks/\allowbreak{}via\_\allowbreak{}malfunction/\allowbreak{}lean-toolchain} (to \texttt{v4.33.0-rc2}) and \texttt{benchmarks/\allowbreak{}via\_\allowbreak{}malfunction/\allowbreak{}lake-manifest.json} (regenerated by \texttt{lake update}: the same format change, and \texttt{lean4lean} and \texttt{batteries} as inherited git packages); \texttt{benchmarks/\allowbreak{}lean-toolchain} and \texttt{benchmarks/\allowbreak{}via\_\allowbreak{}lean/\allowbreak{}lean-toolchain.template} (to \texttt{v4.33.0-rc2}).
\item \texttt{.github/\allowbreak{}workflows/\allowbreak{}build.yml}: the lean-action step gets \texttt{build-args: \textquotedbl{}LeanToLambdaBox Lean4Lean Lean4Lean.Theory Lean4Lean.Verify\textquotedbl{}} and a comment.
\item new \leanfile{tests/\allowbreak{}regress/\allowbreak{}compiler\_\allowbreak{}api.lean}, its expected outputs \texttt{tests/\allowbreak{}regress/\allowbreak{}expected/\allowbreak{}compiler\_\allowbreak{}api/} (19 files) and \texttt{tests/\allowbreak{}regress/\allowbreak{}expected-peregrine/\allowbreak{}compiler\_\allowbreak{}api/} (2 files).
\item \texttt{tests/\allowbreak{}regress/\allowbreak{}expected/}: 29 \texttt{.ast} files re-baselined: \texttt{auto\_\allowbreak{}inline/\allowbreak{}\{default,default3,off,on, on3\}.ast}, \texttt{auto\_\allowbreak{}inline\_\allowbreak{}noinline/\allowbreak{}useAll2.ast}, \texttt{auto\_\allowbreak{}inline\_\allowbreak{}recursive/\allowbreak{}\{useBoth,useBoth2\}.ast}, \texttt{auto\_\allowbreak{}inline\_\allowbreak{}shapes/\allowbreak{}\{useAll,useAll2\}.ast}, \texttt{auto\_\allowbreak{}inline\_\allowbreak{}size/\allowbreak{}\{chain2,twice,twice3\}.ast}, \texttt{auto\_\allowbreak{}inline\_\allowbreak{}values/\allowbreak{}\{useAll,useAll2\}.ast}, \texttt{mli\_\allowbreak{}types/\allowbreak{}\{rArr,rFunPair,rHof,rInt,rListFun, rListPair,rNestArg,rNestL,rNestR,rOpt,rOptFun,rPair,rStr\}.ast}, \texttt{smoke/\allowbreak{}fact3.ast}.
\item \texttt{doc/\allowbreak{}SHIPPING-CHANGES.md}: this entry; the reproductions of R-6, R-12 and R-19, which named the behaviour of Lean v4.22, now give that of v4.33.
\end{itemize}

No file of the shipping code imports \texttt{lean4lean}.
\item \textbf{Why necessary:} \S{}4 of the specification: the \texttt{shipping} branch is the shipping code ported to the \texttt{lean4lean} base, with a Lake dependency on \texttt{lean4lean} \texttt{master} at \texttt{8223d223} and the toolchain bumped to its \texttt{v4.33.0-rc2}, so that the verification code can build on \texttt{lean4lean} and the eraser in one Lake workspace. With only \texttt{lean-toolchain}, \texttt{lakefile.toml} and \texttt{lake-manifest.json} changed, \texttt{lake build} fails with 11 errors in \texttt{Erasure.lean}; the four code edits are the smallest ones that remove them:

\begin{itemize}
\item E1: \texttt{Unknown identifier 'projection\_\allowbreak{}body'}. No such name exists. Lean v4.22 accepted it because its \texttt{do} elaborator did not elaborate the type ascription of \texttt{let x : T \ensuremath{\leftarrow} if \ldots{} then \ldots{} else \ldots{}} (a probe \texttt{let xs : List noSuchType \ensuremath{\leftarrow} if b then pure [1] else pure []} runs on v4.22 and is rejected with \texttt{Unknown identifier} on v4.33). The value already had type \texttt{List ProjectionBody}.
\item E2: three errors, all caused by the predicate of line 264 (\texttt{Invalid field notation: Type of c is not known} twice, \texttt{stuck at solving universe constraint}). In v4.33, \texttt{String.all (s : String) (pat : \ensuremath{\rho}) [ForwardPattern pat]} takes a pattern instead of a \texttt{Char \ensuremath{\to} Bool} (\leanloc{Init/\allowbreak{}Data/\allowbreak{}String/\allowbreak{}TakeDrop.lean}{241}), so the binder type of the lambda is no longer known. The edit states the binder type and the \texttt{decide} that v4.22 inserted; the \texttt{ForwardPattern} instance of a \texttt{Char \ensuremath{\to} Bool} accepts a string iff every character satisfies the predicate.
\item E3: \texttt{Unknown identifier 'Compiler.CSimp.replaceConstants'}. v4.33 removed it; only the single-constant \texttt{CSimp.replaceConstant?} remains, and the map of the \texttt{csimp} extension now holds \texttt{CSimp.Entry} records (\texttt{fromDeclName}, \texttt{toDeclName}, \texttt{thmName}) instead of names (\leanfile{Lean/\allowbreak{}Compiler/\allowbreak{}CSimpAttr.lean}). \texttt{csimpReplaceConstants} is the v4.22 body with \texttt{entry.toDeclName} in place of the name.
\item E4: five errors. \texttt{CasesInfo} moved from \texttt{Lean.Compiler.LCNF} to \texttt{Lean} (\leanfile{Lean/\allowbreak{}Meta/\allowbreak{}CasesInfo.lean}); \texttt{altsRange} is a \texttt{Std.Rco Nat} (field \texttt{lower}, no \texttt{start}) and \texttt{altNumParams} an \texttt{Array CasesAltInfo} (\texttt{.ctor ctorName numFields} or \texttt{.default numHyps}) instead of an \texttt{Array Nat}. For a \texttt{T.casesOn}, every entry is \texttt{.ctor} and \texttt{lower} is the old \texttt{start}.
\end{itemize}

Benchmarks: \texttt{benchmarks/\allowbreak{}via\_\allowbreak{}malfunction} is the Lake workspace in which the benchmark Makefile runs \texttt{lake lean}, and it requires the root package by path; with its toolchain left at v4.22, \texttt{lake build LeanToLambdaBox} there fails with 6 errors (\texttt{Invalid field 'toDeclName'}, \texttt{Invalid field 'lower'} four times, \texttt{Unknown identifier 'Nat.ctor'}). \texttt{from\_\allowbreak{}lean\_\allowbreak{}common} (\texttt{benchmarks/}) is compiled by \texttt{via\_\allowbreak{}malfunction} with v4.33 and by the per-test packages of \texttt{via\_\allowbreak{}lean} with the toolchain of \texttt{lean-toolchain.template}, both into \texttt{benchmarks/\allowbreak{}.lake/\allowbreak{}build}; with two toolchains each build recompiles it for its own (observed: 15 jobs rebuilt at every alternation), and the Lean-native baseline of the benchmarks would be compiled by another Lean than the one whose terms the eraser reads. \texttt{FromLeanCommon} compiles on v4.33 without change, and \texttt{make -C benchmarks/\allowbreak{}via\_\allowbreak{}lean bin/\allowbreak{}even} builds a binary that prints 1, 0 and 1 for 0, 7 and 1000.

CI: lean-action installs the toolchain that \texttt{lean-toolchain} names, and saves \texttt{.lake} to the cache right after its build step, so the \texttt{lean4lean} libraries are built in that step (\texttt{build-args}) to be cached with it. \texttt{Lean4Lean}, \texttt{Lean4Lean.Theory} and \texttt{Lean4Lean.Verify} take 130 jobs (1 min 42 s here) and give 22 \texttt{declaration uses 'sorry'} warnings and no error. The regression step still runs \texttt{lake build} without arguments, which builds \texttt{LeanToLambdaBox} only.
\item \textbf{Behaviour before:} at the parent commit (Lean v4.22), the eraser as in S-11. With only \texttt{lean-toolchain}, \texttt{lakefile.toml} and \texttt{lake-manifest.json} changed, \texttt{lake build} fails: \texttt{Erasure.lean:242:22} (E1), \texttt{:264:27}, \texttt{:264:38}, \texttt{:260:48} (E2), \texttt{:443:9} (E3), \texttt{:602:47}, \texttt{:604:47}, \texttt{:622:48}, \texttt{:623:50}, \texttt{:641:27} (E4), line numbers of the parent commit.
\item \textbf{Behaviour after:} \texttt{lake build} succeeds (6 jobs), with one warning that no edit requires: \texttt{Printing.lean:162}: \texttt{List.asString} has been deprecated. \texttt{lake build LeanToLambdaBox Lean4Lean Lean4Lean.Theory Lean4Lean.Verify} succeeds (136 jobs); a plain \texttt{lake build} does not build \texttt{lean4lean}. \texttt{scripts/\allowbreak{}regress.sh} passes its 12 tests, with and without \texttt{PEREGRINE}.

\begin{itemize}
\item On the inputs of \leanfile{tests/\allowbreak{}regress/\allowbreak{}compiler\_\allowbreak{}api.lean}, which exercise E1 to E4, the 19 outputs are those of the parent commit on v4.22 up to hygienic names (8 \texttt{.ast} files differ only there; the \texttt{.ast.inlinings} and \texttt{.mli} files are byte-identical), and peregrine prints the same result (\texttt{all} evaluates to 21 under Peano naturals).
\item Re-baselined goldens: 28 of the 29 \texttt{.ast} files differ only in hygienic binder names. Lean v4.33 encodes macro scopes differently: \texttt{x.\_\allowbreak{}@.M.\_\allowbreak{}hyg.N} is now \texttt{x.\_\allowbreak{}@.M.\textless{}hash\textgreater{}.\_\allowbreak{}hygCtx.\_\allowbreak{}hyg.N'} (for example \texttt{x.\_\allowbreak{}@.\_\allowbreak{}stdin.\_\allowbreak{}hyg.18} becomes \texttt{x.\_\allowbreak{}@.\_\allowbreak{}stdin.303384890.\_\allowbreak{}hygCtx.\_\allowbreak{}hyg.5}); the number of such names per file is unchanged. \texttt{mli\_\allowbreak{}types/\allowbreak{}rStr.ast} differs in two more ways, both from the standard library: the two hygienic binders of \texttt{List.replicateTR.loop} are based on \texttt{a} instead of \texttt{x}, and the declaration of the inductive \texttt{String} is no longer emitted. In v4.33, \texttt{String} is \texttt{structure String where ofByteArray :: bytes : ByteArray \ldots{}} and \texttt{String.mk} is a deprecated \texttt{@[extern \textquotedbl{}lean\_\allowbreak{}string\_\allowbreak{}mk\textquotedbl{}]} definition (\leanloc{Init/\allowbreak{}Data/\allowbreak{}String/\allowbreak{}Bootstrap.lean}{149}), so the program refers to the axiom \texttt{String.mk} without registering \texttt{String} through a constructor. Its seven axioms and its \texttt{.mli} are unchanged. Every peregrine output of the regression tests is unchanged.
\item \textbf{Known defect left at this commit, fixed by the next entry: sparse \texttt{casesOn}.} In v4.33, \texttt{getCasesInfo?} also recognizes two new kinds of declarations as cases-like (\texttt{Lean.isCasesOnLike}, \leanloc{Lean/\allowbreak{}AuxRecursor.lean}{63}): the sparse \texttt{casesOn} that the match compiler creates for a match with a wildcard or inaccessible pattern (\texttt{F.\_\allowbreak{}sparseCasesOn\_\allowbreak{}\textless{}i\textgreater{}}: alternatives for some constructors, and a catch-all alternative whose \texttt{CasesAltInfo} is \texttt{.default}), and the per-constructor eliminators \texttt{T.c.elim} of the derived \texttt{BEq} and \texttt{DecidableEq} of an inductive with at least 10 constructors. \texttt{erase.visitCases} takes the type name as \texttt{casesInfo.declName.getPrefix}, which for these is a definition or a constructor, so \texttt{let .inductInfo indVal \ensuremath{\leftarrow} getConstInfo typeName \textbar{} unreachable!} panics (\texttt{PANIC at Erasure.erase.visitCases LeanToLambdaBox.Erasure:652:55: unreachable code has been reached}); the panic's default value makes the whole match \texttt{\ensuremath{\Box}}. \texttt{\#erase} still succeeds, exits with status 0, and peregrine validates the output. Reproductions:

\begin{itemize}
\item peregrine eval of corpus \texttt{examples/\allowbreak{}PortProbe/\allowbreak{}predOr0.peano.ast} (\texttt{predOr0 3}, with \texttt{predOr0 : \textbar{} n+1 =\textgreater{} n \textbar{} \_\allowbreak{} =\textgreater{} 0}) prints \texttt{constr con\_\allowbreak{}15} / \texttt{constr con\_\allowbreak{}105} instead of \texttt{Nat.succ (Nat.succ Nat.zero)}; \texttt{secondN.peano.ast} (\texttt{secondN 3}) prints \texttt{constr con\_\allowbreak{}15} / \texttt{constr con\_\allowbreak{}108} instead of \texttt{Nat.succ Nat.zero};
\item the 20 natio benchmarks built natively (benchmark Makefile, peregrine with \texttt{unbox.config} after the R-1 rewrite, malfunction and OCaml 4.14.2 without flambda, \texttt{JCFArrayOCaml4.ml}) and run on 0, 1, 2, 5, 10, 50 and 1000 (0, 1, 2, 5 and 8 for binarytrees, const\_fold and deriv): 15 programs print the same values as at the parent commit; \texttt{rbmap\_\allowbreak{}beans}, \texttt{rbmap\_\allowbreak{}std}, \texttt{rbmap\_\allowbreak{}mono} and \texttt{rbmap\_\allowbreak{}raw} print 1 for 50 and 1000 instead of 5 and 100; \texttt{const\_\allowbreak{}fold} crashes with a segmentation fault on every input; \texttt{deriv} crashes on 0 and runs past a 60 s timeout without output on 1, 2, 5 and 8 (the parent commit prints 2, 6, 22, 2202 and 598592). The affected declarations of the corpus are listed below, under "sparse \texttt{casesOn}".
\end{itemize}
\end{itemize}
\item \textbf{Effect on emitted .ast (corpus):} of 284 files, 155 are byte-identical to the corpus of S-11 and 129 differ (none added or removed). All 129 contain hygienic names, which Lean v4.33 prints in the new encoding; \texttt{scripts/\allowbreak{}corpus-diff.sh -{}-normalize} reports 155 identical, 55 identical after normalization and 74 differing (53 \texttt{.ast}, 21 \texttt{.ast.inlinings}; every \texttt{.mli} is byte-identical). The 55 are \texttt{benchmarks/\allowbreak{}\{prune,noprune\}/\allowbreak{}\{even,iflazy,triangle\_\allowbreak{}acc, triangle\_\allowbreak{}rec\}.ast}, 14 of the 16 \texttt{examples/\allowbreak{}Defects/\allowbreak{}*.ast} (\texttt{strlit.ast} is byte-identical and \texttt{wf\_\allowbreak{}rec.ast} differs), 20 of the 27 \texttt{examples/\allowbreak{}Scope/\allowbreak{}*.ast} (the other 7 are byte-identical) and 13 of the 33 \texttt{examples/\allowbreak{}PortProbe/\allowbreak{}*.ast} (\texttt{double}, \texttt{fact} and \texttt{width} in the three configurations, \texttt{useHalf} and \texttt{usePred} default and pruned). The axioms of every file are unchanged, except in predOr0 default and pruned, which lose \texttt{Nat.beq} and \texttt{Nat.sub} (sparse \texttt{casesOn}, below). Every difference of the 74 files falls in one of these classes (declarations compared after parsing, per kername):

\begin{itemize}
\item \textbf{Printing (standard library):} names that Lean v4.33 gives differently, the terms being the same. The hygienic binders of the core definitions \texttt{List.mapTR.loop}, \texttt{List.replicateTR.loop} and \texttt{List.range.loop} are based on \texttt{a} instead of \texttt{x} (binarytrees, list\_sum\_foldl, list\_sum\_foldr, list\_sum\_rev, triangle\_foldl, qsort, qsort\_fin); the core made \texttt{Array.qsort.sort} and \texttt{Array.qpartition.loop} private, so their kernames become \texttt{\_\allowbreak{}private.Init.Data.Array.QSort.Basic.0.Array.qsort.sort} and \texttt{\_\allowbreak{}private.Init.Data.Array.QSort.Basic.0.Array.qpartition.loop} (qsort, qsort\_fin, qsort\_single); and \texttt{Nat.recCompiled} is no longer private (\texttt{colorCode.peano}, \texttt{shapeEq.peano}, \texttt{bigEq.peano}). 16 benchmark \texttt{.ast} files (binarytrees, list\_sum\_foldl, list\_sum\_foldr, list\_sum\_rev, triangle\_foldl, qsort, qsort\_fin, qsort\_single, both variants) differ only in these names.
\item \textbf{\texttt{@[inline]} attributes (standard library):} v4.33 tags \texttt{Array.empty} and \texttt{Array.qpartition} \texttt{@[inline]}, so the \texttt{.ast.inlinings} of qsort, qsort\_fin and qsort\_single list both, and those of unionfind and unionfind\_noinline list \texttt{Array.empty} (both variants, 10 files).
\item \textbf{Value patterns (standard library):} a match on literals now goes through the new \texttt{abbrev Eq.ndrec\_\allowbreak{}symm} (\leanloc{Init/\allowbreak{}Prelude.lean}{385}) instead of two nested \texttt{Eq.ndrec}. It changes \texttt{Deriv.Expr.pown}, \texttt{Tiny.colorCode}, \texttt{Tiny.shapeOf} and \texttt{Tiny.bigOf}, adds the declaration \texttt{Eq.ndrec\_\allowbreak{}symm} (an ordinary definition whose body applies \texttt{Eq.ndrec}, so no axiom is added), and adds \texttt{Eq.ndrec\_\allowbreak{}symm} to the \texttt{.ast.inlinings} of deriv (both variants) and of colorCode, shapeEq and bigEq (three configurations each), 11 files.
\item \textbf{\texttt{do} notation (elaborator):} in \texttt{mkNodes} (\texttt{\textbar{} n+1 =\textgreater{} do \_\allowbreak{} \ensuremath{\leftarrow} mk; mkNodes n}) the ignored result is no longer bound by a trivial \texttt{let}, and the binders that \texttt{do} introduces in \texttt{test}, \texttt{findEntryAux}, \texttt{mkNodes}, \texttt{mk} and \texttt{StateT'.bind} are named \texttt{\_\allowbreak{}\_\allowbreak{}r} and \texttt{\_\allowbreak{}\_\allowbreak{}x} instead of \texttt{x} and \texttt{\_\allowbreak{}\_\allowbreak{}discr} (unionfind, unionfind\_noinline, both variants).
\item \textbf{Derived instances (standard library):} the auxiliary definitions of derived \texttt{BEq} and \texttt{DecidableEq} get stable names (\texttt{Tiny.instBEqPt.beq}, \texttt{Tiny.instBEqShape.beq}, \texttt{Tiny.instDecidableEqShape.decEq}, \ldots{}) instead of hygienic ones (\texttt{Tiny.beqPt.\_\allowbreak{}@.\_\allowbreak{}stdin.\_\allowbreak{}hyg.N}, \ldots{}), with bodies identical for ptEqN (three configurations) and for the \texttt{DecidableEq} of Shape. For the 10 constructors of \texttt{Tiny.Big}, v4.33 compares \texttt{Tiny.Big.ctorIdx} (a new declaration) with the root-level \texttt{decEq} (also new) and uses per-constructor eliminators.
\item \textbf{Core \texttt{Nat} functions (standard library)}, erased from their definitions under Peano naturals: \texttt{Nat.div} and \texttt{Nat.modCore} pass \texttt{Nat.succ x} instead of \texttt{x + 1} (so \texttt{Nat.add}, \texttt{instAddNat}, \texttt{instHAdd}, \texttt{Add}, \texttt{Add.add}, \texttt{HAdd} and \texttt{HAdd.hAdd} are no longer emitted in \texttt{colorCode.peano}); \texttt{Nat.div.go} and \texttt{Nat.modCore.go} are compiled differently (new \texttt{Nat.modCore.go.\_\allowbreak{}f} and \texttt{Nat.brecOn.go}); \texttt{Nat.ble} no longer splits on its second argument when the first is 0. Files: \texttt{Defects/\allowbreak{}wf\_\allowbreak{}rec.ast}, \texttt{PortProbe/\allowbreak{}\{useHalf,usePred\}.peano.ast}, and the peano variants of colorCode, shapeEq and bigEq. The declaration order changes in wf\_rec, useHalf.peano and bigEq.peano, whose dependencies change.
\item \textbf{Sparse \texttt{casesOn} (the known defect above):} a match compiled to a sparse \texttt{casesOn} or a per-constructor eliminator becomes \texttt{\ensuremath{\Box}}: \texttt{Expr.reassoc}, \texttt{Expr.appendAdd}, \texttt{Expr.appendMul}, \texttt{Expr.constFolding} (const\_fold, 6 matches); \texttt{Deriv.Expr.ln}, \texttt{add}, \texttt{mul}, \texttt{pow} (deriv); \texttt{setBlack}, \texttt{isRed}, \texttt{balance1}, \texttt{balance2} (each rbmap benchmark); \texttt{Tiny.isRed} (colorCode), \texttt{Tiny.second} (secondN), \texttt{Tiny.predOr0} (predOr0), the three branches of \texttt{Tiny.instBEqShape.beq} (shapeEq), and the ten eliminator branches of each of \texttt{Tiny.instBEqBig.beq} and \texttt{Tiny.instDecidableEqBig.decEq} (bigEq). These are all the \texttt{PANIC} messages of the corpus logs apart from the four of R-4 and R-5: 26 per benchmark variant, 78 in \texttt{PortProbe}. Declarations reachable only from a lost match are no longer emitted (\texttt{Unit.unit} and \texttt{PUnit} in deriv, which also loses \texttt{Unit.unit} from its \texttt{.ast.inlinings}; \texttt{Nat}, \texttt{Bool}, \texttt{Nat.beq} and \texttt{Nat.sub} in predOr0 default and pruned), and the declaration order changes in the rbmap benchmarks and secondN.
\end{itemize}

No difference comes from a change of the eraser's behaviour on an unchanged input term: E1 to E4 change no output (see \texttt{compiler\_\allowbreak{}api} above), and the sparse \texttt{casesOn} defect comes from declarations that the v4.22 match compiler did not produce.
\item \textbf{Regression test:} \leanfile{tests/\allowbreak{}regress/\allowbreak{}compiler\_\allowbreak{}api.lean} guards E1 to E4: the projections of a structure (\texttt{ptSum.ast} declares \texttt{Pt} with two \texttt{projection\_\allowbreak{}body} entries), an ASCII and a non-ASCII binder (\texttt{names.ast}: \texttt{a} is named, \texttt{\ensuremath{\alpha}\ensuremath{_1}} is \texttt{nAnon}), \texttt{csimp} on and off (\texttt{sumFoldr.ast} uses \texttt{List.foldrTR}, \texttt{sumFoldr.nocsimp.ast} uses \texttt{List.foldr}), and the alternatives of the machine \texttt{Nat} and \texttt{Int} paths and of a user inductive (\texttt{natCase.ast}, \texttt{intCase.ast}, \texttt{triVal.ast}); with \texttt{PEREGRINE} set, \texttt{all.ast} validates and evaluates to 21. The parent commit cannot run it on v4.33, since the package does not compile there; on v4.22 it passes with the same outputs up to hygienic names. The 29 re-baselined goldens keep guarding what they guarded, and \texttt{smoke}, \texttt{mli\_\allowbreak{}types} and the auto-inline tests pass unchanged apart from them.
\end{itemize}

\section{\texorpdfstring{S-13: Sparse \texttt{casesOn} and per-constructor eliminators are erased by constructor name}{S-13: Sparse casesOn and per-constructor eliminators are erased by constructor name}}\label{reg:S-13}

\begin{itemize}
\item \textbf{Commit:} the commit whose subject starts with \texttt{shipping(S-13):} (\texttt{git log -{}-grep='\textasciicircum{}shipping(S-13):'}).
\item \textbf{Files and functions:}

\begin{itemize}
\item \leanfile{LeanToLambdaBox/\allowbreak{}Erasure.lean}: \texttt{erase.visitCases} (new docstring). The inductive type is \texttt{casesInfo.indName} instead of the prefix of \texttt{casesInfo.declName}. The alternatives are collected by constructor name from the \texttt{.ctor} entries of \texttt{casesInfo.altNumParams}, and the catch-all from its \texttt{.default} entry. The branch of a constructor without an alternative is the catch-all applied to one \texttt{\ensuremath{\Box}} per hypothesis, or \texttt{\ensuremath{\Box}} when there is no catch-all; the catch-all is erased once, and only if some constructor lacks an alternative. The machine-\texttt{Nat} path (arms of \texttt{Nat.zero} and \texttt{Nat.succ}), the machine-\texttt{Int} path (arms of \texttt{Int.ofNat} and \texttt{Int.negSucc}) and the generic path (a loop over the constructors of the type) look up each arm by name. The generic path no longer throws an error on a \texttt{.default} entry (the E4 edit of S-12).
\item new \leanfile{tests/\allowbreak{}regress/\allowbreak{}sparse\_\allowbreak{}cases.lean}, its expected outputs \texttt{tests/\allowbreak{}regress/\allowbreak{}expected/\allowbreak{}sparse\_\allowbreak{}cases/} (32 files) and \texttt{tests/\allowbreak{}regress/\allowbreak{}expected-peregrine/\allowbreak{}sparse\_\allowbreak{}cases/} (4 files).
\item \texttt{doc/\allowbreak{}SHIPPING-CHANGES.md}: this entry, and a report of the defect that S-14 fixes.
\end{itemize}
\item \textbf{Why necessary:} the known defect of S-12 ("sparse \texttt{casesOn}"), a silent miscompilation: in Lean v4.33, \texttt{getCasesInfo?} (\leanloc{Lean/\allowbreak{}Meta/\allowbreak{}CasesInfo.lean}{56}) also describes two kinds of declarations that are not a \texttt{T.casesOn}, and \texttt{erase.visitCases} turned each application of them into \texttt{\ensuremath{\Box}} after a panic, with exit status 0 and output that \texttt{peregrine validate} accepts.

\begin{itemize}
\item The sparse \texttt{casesOn} \texttt{F.\_\allowbreak{}sparseCasesOn\_\allowbreak{}\textless{}i\textgreater{}} that the match compiler creates for a match with a wildcard or inaccessible pattern. Its alternatives are those of the constructors that the match names, in the order in which they first appear in the match (\texttt{collectCtors}, \leanloc{Lean/\allowbreak{}Meta/\allowbreak{}Match/\allowbreak{}Match.lean}{543}), so not necessarily the constructor order; a last argument, the catch-all, takes a proof of \texttt{Nat.hasNotBit mask x.ctorIdx}. Its definition (\texttt{mkSparseCasesOn}, \leanloc{Lean/\allowbreak{}Meta/\allowbreak{}Constructions/\allowbreak{}SparseCasesOn.lean}{59}) is a \texttt{T.rec} whose minor premise for a constructor without alternative applies the catch-all to such a proof; the erasure of that application is the erased catch-all applied to \texttt{\ensuremath{\Box}}.
\item The per-constructor eliminator \texttt{T.c.elim motive x h alt} (\texttt{mkConstructorElim}, \leanloc{Lean/\allowbreak{}Meta/\allowbreak{}Constructions/\allowbreak{}CtorElim.lean}{165}), whose side condition \texttt{h : x.ctorIdx = i} makes every constructor other than \texttt{c} unreachable; it has one alternative and no catch-all. The derived \texttt{BEq} and \texttt{DecidableEq} of an inductive with at least 10 constructors use it.
\end{itemize}

Lean's own compiler treats both the same way (\texttt{ToLCNF.visitAlt} and \texttt{ToLCNF.visitCases}, \leanloc{Lean/\allowbreak{}Compiler/\allowbreak{}LCNF/\allowbreak{}ToLCNF.lean}{584,621}): it applies the catch-all to erased arguments and leaves constructors without an alternative out of the \texttt{cases}. A \ensuremath{\lambda}\ensuremath{\Box} \texttt{case} needs one branch per constructor, so an unreachable branch is \texttt{\ensuremath{\Box}}.
\item \textbf{Behaviour before:} at the parent commit, with \leanfile{tests/\allowbreak{}regress/\allowbreak{}sparse\_\allowbreak{}cases.lean} copied into a checkout: \texttt{lake env lean} exits with status 0 after 47 messages \texttt{PANIC at Erasure.erase.visitCases LeanToLambdaBox.Erasure:652:55: unreachable code has been reached}, and every match that Lean compiles to a sparse \texttt{casesOn} or an eliminator is \texttt{\ensuremath{\Box}}: for example \texttt{Sparse.isRed} is \texttt{\ensuremath{\lambda}x. let \_\allowbreak{}alt := \ldots{} in let \_\allowbreak{}alt := \ldots{} in \ensuremath{\Box}} and \texttt{Sparse.viaElim} is \texttt{\ensuremath{\lambda}x h k. \ensuremath{\Box} k}. \texttt{peregrine validate} accepts \texttt{count.ast} and \texttt{sum.ast}, and \texttt{peregrine eval} of either fails with \texttt{Could not evaluate program: Case: \textless{}15\textgreater{} branch not found}. The corpus reproductions of S-12 hold: \texttt{peregrine eval} of \texttt{examples/\allowbreak{}PortProbe/\allowbreak{}predOr0.peano.ast} and \texttt{secondN.peano.ast} prints \texttt{constr con\_\allowbreak{}15} / \texttt{constr con\_\allowbreak{}105} and \texttt{constr con\_\allowbreak{}15} / \texttt{constr con\_\allowbreak{}108}; natively, \texttt{rbmap\_\allowbreak{}beans}, \texttt{rbmap\_\allowbreak{}std}, \texttt{rbmap\_\allowbreak{}mono} and \texttt{rbmap\_\allowbreak{}raw} print 1 for 50 and 1000, \texttt{const\_\allowbreak{}fold} crashes on every input, and \texttt{deriv} crashes or runs past a 60 s timeout.
\item \textbf{Behaviour after:} no panic. The same reproduction gives, for example, with \texttt{r} and \texttt{w} the let-bound alternatives of \texttt{.red} and of the wildcard, \texttt{Sparse.isRed = \ensuremath{\lambda}x. let r := \ldots{} in let w := \ldots{} in case x of red =\textgreater{} r Unit.unit \textbar{} green =\textgreater{} (\ensuremath{\lambda}h. w x) \ensuremath{\Box} \textbar{} blue =\textgreater{} (\ensuremath{\lambda}h. w x) \ensuremath{\Box} \textbar{} black =\textgreater{} (\ensuremath{\lambda}h. w x) \ensuremath{\Box}}; \texttt{Sparse.pick}, whose match lists \texttt{.c} before \texttt{.a}, gets each alternative in its constructor's branch; the machine-\texttt{Nat} path gives, with \texttt{s} and \texttt{w} the alternatives of \texttt{n + 1} and of the wildcard, \texttt{Sparse.predOr0 = \ensuremath{\lambda}x. let s := \ldots{} in let w := \ldots{} in let n := x in case (Nat.beq n 0) of false =\textgreater{} (\ensuremath{\lambda}n. s n) (Nat.sub n 1) \textbar{} true =\textgreater{} (\ensuremath{\lambda}h. w x) \ensuremath{\Box}}; and \texttt{Sparse.viaElim = \ensuremath{\lambda}x h k. (case x of c0 \_\allowbreak{} =\textgreater{} \ensuremath{\Box} \textbar{} c1 a b =\textgreater{} \ensuremath{\lambda}k. a + b * k \textbar{} c2 =\textgreater{} \ensuremath{\Box} \textbar{} \ldots{} \textbar{} c9 \_\allowbreak{} =\textgreater{} \ensuremath{\Box}) k}. With \texttt{PEREGRINE} set, \texttt{count.ast} evaluates to 25 (all 25 results of the test equal Lean's values, which \texttt{\#guard} checks) and \texttt{sum.ast} to 90 (their sum). Further checks, with the tools of the benchmark pipeline (peregrine with \texttt{unbox.config} after the R-1 rewrite, malfunction, OCaml 4.14.2 without flambda, the runtime of \texttt{benchmarks/\allowbreak{}via\_\allowbreak{}malfunction} with \texttt{decidable-pruned.ml} and \texttt{JCFArrayOCaml4.ml}):

\begin{itemize}
\item \texttt{peregrine eval} of the corpus files \texttt{predOr0.peano.ast} and \texttt{secondN.peano.ast} prints \texttt{Nat.succ (Nat.succ Nat.zero)} and \texttt{Nat.succ Nat.zero}, Lean's values of \texttt{predOr0 3} and \texttt{secondN 3};
\item the 20 natio benchmarks, built natively and run on 0, 1, 2, 5, 10, 50 and 1000 (0, 1, 2, 5 and 8 for binarytrees, const\_fold and deriv), print in all 134 runs the values of S-11, which are those of Lean's \texttt{\#eval};
\item the 11 functions of \leanfile{tests/\allowbreak{}corpus/\allowbreak{}PortProbe.lean}, erased with constructor pruning and run natively on 16 inputs from 0 to 101, print Lean's \texttt{\#eval} values in all 176 runs; so do the 13 functions of \leanfile{tests/\allowbreak{}regress/\allowbreak{}sparse\_\allowbreak{}cases.lean}, each wrapped as a function \texttt{Nat \ensuremath{\to} Nat} and erased with constructor pruning, in all 208 runs; these cover the machine-\texttt{Nat} and machine-\texttt{Int} paths, which peregrine cannot evaluate (R-1).
\end{itemize}

On a plain \texttt{T.casesOn} every constructor has an alternative, in constructor order, and there is no catch-all, so \texttt{visitCases} makes the same calls in the same order as before: the 12 other regression tests pass unchanged, \texttt{compiler\_\allowbreak{}api} among them (the machine-\texttt{Nat} and machine-\texttt{Int} paths and a user inductive), and 252 of the 284 corpus files are byte-identical.
\item \textbf{Effect on emitted .ast (corpus):} 252 of the 284 files are byte-identical to the corpus of S-12; 32 differ (27 \texttt{.ast} and 5 \texttt{.ast.inlinings} files; every \texttt{.mli} is byte-identical), all listed by S-12 under "sparse \texttt{casesOn}". In each file, the declarations that change are exactly those that contain such a match, whose \texttt{\ensuremath{\Box}} is now the match:

\begin{itemize}
\item \texttt{benchmarks/\allowbreak{}\{prune,noprune\}/\allowbreak{}const\_\allowbreak{}fold.ast}: \texttt{Expr.reassoc}, \texttt{Expr.appendAdd}, \texttt{Expr.appendMul}, \texttt{Expr.constFolding}.
\item \texttt{benchmarks/\allowbreak{}\{prune,noprune\}/\allowbreak{}deriv.ast}: \texttt{Deriv.Expr.ln}, \texttt{add}, \texttt{mul}, \texttt{pow}; the declarations \texttt{Unit.unit} and \texttt{PUnit}, reachable only from these matches, are emitted again, and \texttt{deriv.ast.inlinings} (both variants) lists \texttt{Unit.unit} again.
\item \texttt{benchmarks/\allowbreak{}\{prune,noprune\}/\allowbreak{}rbmap\_\allowbreak{}\{beans,mono,raw,std\}.ast}: \texttt{setBlack}, \texttt{isRed}, \texttt{balance1}, \texttt{balance2}; the declaration order is that of S-11 again.
\item \texttt{examples/\allowbreak{}PortProbe/}, each in its three configurations (\texttt{default}, \texttt{prune}, \texttt{peano}): \texttt{colorCode} (\texttt{Tiny.isRed}), \texttt{secondN} (\texttt{Tiny.second}; the declaration order is that of S-11 again), \texttt{predOr0} (\texttt{Tiny.predOr0}; in \texttt{default} and \texttt{prune}, \texttt{Bool} and the axioms \texttt{Nat.beq} and \texttt{Nat.sub} are emitted again), \texttt{shapeEq} (\texttt{Tiny.instBEqShape.beq}), \texttt{bigEq} (\texttt{Tiny.instBEqBig.beq}, \texttt{Tiny.instDecidableEqBig.decEq}; the declarations and the entries of \texttt{bigEq.*.ast.inlinings} change order).
\end{itemize}

Against the corpus of S-11 (Lean v4.22), with binder names and \texttt{\_\allowbreak{}private} prefixes ignored, the axioms of all 32 files are the same, the declarations are the same apart from the standard-library differences listed in S-12, and the recovered matches differ from v4.22 only where v4.33's match compiler does. A constructor that a wildcard covers gets \texttt{(\ensuremath{\lambda}h. alt d) \ensuremath{\Box}}, where \texttt{d} is the discriminant, in place of v4.22's \texttt{alt (C args)} with the constructor rebuilt from the matched fields (for \texttt{RBNode.isRed}: \texttt{leaf =\textgreater{} (\ensuremath{\lambda}h. alt t) \ensuremath{\Box}} for \texttt{alt (leaf \ensuremath{\Box} \ensuremath{\Box})}); and the case trees of \texttt{Deriv.Expr.mul}, \texttt{Deriv.Expr.pow} and \texttt{Expr.constFolding} have 33, 22 and 13 \texttt{case} nodes where v4.22 had 35, 23 and 15.
\item \textbf{Regression test:} \leanfile{tests/\allowbreak{}regress/\allowbreak{}sparse\_\allowbreak{}cases.lean} erases a match of each kind on each path of \texttt{visitCases}: the generic path with a catch-all for three constructors (\texttt{isRed}), alternatives out of constructor order (\texttt{pick}), a catch-all that receives the scrutinee (\texttt{leftOr}), nested matches (\texttt{second}), a discriminant that is not a variable (\texttt{redCode}, see S-14), a catch-all for a constructor with a proof field with and without pruning (\texttt{optVal}) and a match applied to an extra argument (\texttt{applyTo}); the machine \texttt{Nat} path with a catch-all for \texttt{Nat.zero} (\texttt{predOr0}) and for \texttt{Nat.succ} (\texttt{isZero}); the machine \texttt{Int} path with a catch-all for \texttt{Int.ofNat} (\texttt{negPart}) and for \texttt{Int.negSucc} (\texttt{natPart}); and eliminators, called directly with an extra argument (\texttt{viaElim}) and through a derived \texttt{BEq} of 10 constructors (\texttt{bigEq}). With \texttt{PEREGRINE} set, \texttt{count.ast} and \texttt{sum.ast} validate and evaluate to 25 and 90. It fails before (47 panics, 23 of its 32 outputs differ, \texttt{peregrine eval} fails) and passes after.
\end{itemize}

\section{\texorpdfstring{S-14: A catch-all no longer re-evaluates a discriminant that is not a variable}{S-14: A catch-all no longer re-evaluates a discriminant that is not a variable}}\label{reg:S-14}

\begin{itemize}
\item \textbf{Commit:} the commit whose subject starts with \texttt{shipping(S-14):} (\texttt{git log -{}-grep='\textasciicircum{}shipping(S-14):'}).
\item \textbf{Files and functions:}

\begin{itemize}
\item \leanfile{LeanToLambdaBox/\allowbreak{}Erasure.lean}: \texttt{erase.visitCases} (docstring extended) and new \texttt{erase.visitCasesOn}. \texttt{visitCases} erases the discriminant and collects the alternatives and the catch-all as in S-13, then drops the catch-all if every constructor has an alternative. If the catch-all is used and the erased discriminant is neither a variable nor \ensuremath{\Box}, it adds a let-declaration \texttt{discr} whose value is the discriminant to the local context, replaces every occurrence of the discriminant in the catch-all by \texttt{discr} (the catch-all stays well-typed, since \texttt{discr} unfolds to the discriminant), builds the \texttt{case} with \texttt{discr} as its discriminant, and returns \texttt{let discr := \textless{}erased discriminant\textgreater{} in \textless{}case\textgreater{}}. The extra arguments of an over-applied \texttt{casesOn} are applied to the result, outside the \texttt{let}, as before. \texttt{visitCasesOn} is the rest of S-13's \texttt{visitCases}, unchanged: the erasure of the catch-all, then the machine-\texttt{Nat}, machine-\texttt{Int} and generic paths.
\item new \leanfile{tests/\allowbreak{}regress/\allowbreak{}sparse\_\allowbreak{}discr.lean}, its expected outputs \texttt{tests/\allowbreak{}regress/\allowbreak{}expected/\allowbreak{}sparse\_\allowbreak{}discr/} (12 files) and \texttt{tests/\allowbreak{}regress/\allowbreak{}expected-peregrine/\allowbreak{}sparse\_\allowbreak{}discr/} (2 files).
\item \texttt{tests/\allowbreak{}regress/\allowbreak{}expected/\allowbreak{}sparse\_\allowbreak{}cases/\allowbreak{}\{redCode,count,sum\}.ast}: the declaration \texttt{Sparse.redCode}, the only match of that test whose discriminant is not a variable.
\item \texttt{doc/\allowbreak{}SHIPPING-CHANGES.md}: this entry.
\end{itemize}
\item \textbf{Why necessary:} a defect of S-13 and a performance regression against Lean v4.22 (\texttt{main}, and S-11 on this branch). In Lean v4.33, the matcher of a match with a wildcard passes its discriminant to the wildcard's alternative inside the catch-all of its sparse \texttt{casesOn} (\texttt{fun motive x h\_\allowbreak{}1 h\_\allowbreak{}2 =\textgreater{} F.\_\allowbreak{}sparseCasesOn\_\allowbreak{}1 x (h\_\allowbreak{}1 ()) fun h =\textgreater{} h\_\allowbreak{}2 x}), and \texttt{inlineMatchers} substitutes the discriminant for \texttt{x}. S-13 erases the catch-all as given, so a discriminant that is a call is computed once as the scrutinee and once more in every branch of the catch-all. The results are right, but a function that matches on its own recursive call computes that call twice at every level, and takes time exponential in the depth of the recursion. Under v4.22 the same match was a plain \texttt{casesOn} whose wildcard alternative received the constructor rebuilt from the matched fields, so the discriminant was computed once, as it is in Lean's own compiled code.
\item \textbf{Behaviour before:} at the parent commit, with \leanfile{tests/\allowbreak{}regress/\allowbreak{}sparse\_\allowbreak{}discr.lean} copied in, \texttt{lean} exits with status 1 and the error \texttt{expected one reference each: [stepN.ast refers to Discr.step 3 times, predSub2.ast refers to Discr.sub2 2 times, negOr.ast refers to Discr.neg 2 times]}, and 4 of the 12 outputs differ from the expected ones (\texttt{stepN}, \texttt{predSub2}, \texttt{negOr} and \texttt{all}). With \texttt{g} and \texttt{w} the let-bound alternatives of \texttt{.green} and of the wildcard, \texttt{Discr.stepN = \ensuremath{\lambda}n. let g := \ldots{} in let w := \ldots{} in case (Discr.step n) of red =\textgreater{} (\ensuremath{\lambda}h. w (Discr.step n)) \ensuremath{\Box} \textbar{} green =\textgreater{} g Unit.unit \textbar{} blue =\textgreater{} (\ensuremath{\lambda}h. w (Discr.step n)) \ensuremath{\Box}}, and \texttt{Discr.step} puts its recursive call \texttt{step n} in the same three places. On the machine paths, \texttt{Discr.predSub2 = \ensuremath{\lambda}n. \ldots{} let n := Discr.sub2 n in case (Nat.beq n 0) of false =\textgreater{} \ldots{} \textbar{} true =\textgreater{} (\ensuremath{\lambda}h. w (Discr.sub2 n)) \ensuremath{\Box}}, and \texttt{Discr.negOr} likewise with \texttt{Discr.neg}. Built natively with the tools of the benchmark pipeline (as in S-13), \texttt{stepN} erased with constructor pruning prints 1 (Lean's value) for 20, 22, 24, 26 and 28 after 0.012, 0.034, 0.120, 0.460 and 1.830 s; under Peano naturals, \texttt{peregrine eval} of \texttt{stepN 16} and \texttt{stepN 20} prints 1 after 0.105 and 1.476 s.
\item \textbf{Behaviour after:} the test passes. For example, \texttt{Discr.stepN = \ensuremath{\lambda}n. let g := \ldots{} in let w := \ldots{} in let discr := Discr.step n in case discr of red =\textgreater{} (\ensuremath{\lambda}h. w discr) \ensuremath{\Box} \textbar{} green =\textgreater{} g Unit.unit \textbar{} blue =\textgreater{} (\ensuremath{\lambda}h. w discr) \ensuremath{\Box}}; \texttt{Discr.step} binds its recursive call in the same way; \texttt{Discr.predSub2 = \ensuremath{\lambda}n. \ldots{} let discr := Discr.sub2 n in let n := discr in case (Nat.beq n 0) of false =\textgreater{} \ldots{} \textbar{} true =\textgreater{} (\ensuremath{\lambda}h. w discr) \ensuremath{\Box}}, and \texttt{Discr.negOr} likewise. Natively, \texttt{stepN} prints 1 for each of 20 to 28 after 0.004 s, and for 1000 and 100000 after 0.004 and 0.006 s (for 1000000 the stack overflows, as the recursion of \texttt{step} is not a tail call); \texttt{peregrine eval} of \texttt{stepN 20} and \texttt{stepN 200} takes 0.011 and 0.029 s. For comparison, S-11 (Lean v4.22) takes 0.004 to 0.007 s natively for 20 to 100000, and Lean v4.33's own compiled code 0.006 to 0.008 s.

A match gets no \texttt{let} if its \texttt{casesOn} has no catch-all, if its catch-all is not used, or if its discriminant is a variable. So the 12 tests other than \texttt{sparse\_\allowbreak{}cases} and \texttt{sparse\_\allowbreak{}discr} pass unchanged, \texttt{compiler\_\allowbreak{}api} (plain \texttt{casesOn} on the three paths) among them; in \texttt{sparse\_\allowbreak{}cases} only the declaration \texttt{Sparse.redCode} changes, and \texttt{count.ast} and \texttt{sum.ast} still evaluate to 25 and 90; and in the new test, \texttt{codeOf} (a plain \texttt{casesOn} on a call) and \texttt{shiftBy} (a match applied to an extra argument, which \texttt{inlineMatchers} leaves as a \ensuremath{\beta}-redex whose parameter is the discriminant) are byte-identical to the parent's output.
\item \textbf{Effect on emitted .ast (corpus):} byte-identical for all 284 files. In the corpus, every \texttt{case} with a catch-all branch (\texttt{(\ensuremath{\lambda}h. \ldots{}) \ensuremath{\Box}}) has a variable discriminant: 267 on the generic path, whose scrutinee is a variable, and 2 on the machine-\texttt{Nat} path (\texttt{Tiny.predOr0} in \texttt{examples/\allowbreak{}PortProbe/\allowbreak{}predOr0.\{default,prune\}.ast}), whose \texttt{n} is bound to a variable.
\item \textbf{Regression test:} \leanfile{tests/\allowbreak{}regress/\allowbreak{}sparse\_\allowbreak{}discr.lean} erases a match with a used catch-all and a discriminant that is a call on each path of \texttt{visitCases}: the generic path (\texttt{stepN}, and \texttt{step}, which matches on its own recursive call), the machine \texttt{Nat} path (\texttt{predSub2}) and the machine \texttt{Int} path (\texttt{negOr}); and, as controls, a plain \texttt{casesOn} on a call (\texttt{codeOf}) and a match applied to an extra argument (\texttt{shiftBy}). The expected outputs show one \texttt{let discr} per such match. An \texttt{\#eval} counts, in each file, the references to the function that computes the discriminant, and fails unless there is exactly one: a check of the emitted term, independent of timing. With \texttt{PEREGRINE} set, \texttt{all.ast} (every function applied to arguments, under Peano naturals) validates and evaluates to 26, Lean's value, which \texttt{\#guard} checks. It fails before (the \texttt{\#eval} reports 3, 2 and 2 references; 4 of 12 outputs differ) and passes after. The expected outputs of \texttt{sparse\_\allowbreak{}cases} pin \texttt{Sparse.redCode} with its \texttt{let}.
\end{itemize}

\section{\texorpdfstring{S-15: The corpus gains examples inside and near the verification scope, irreducible-alias programs and unsafe recursion}{S-15: The corpus gains examples inside and near the verification scope, irreducible-alias programs and unsafe recursion}}\label{reg:S-15}

\begin{itemize}
\item \textbf{Commit:} the commit whose subject starts with \texttt{shipping(S-15):} (\texttt{git log -{}-grep='\textasciicircum{}shipping(S-15):'}).
\item \textbf{Files and functions:}

\begin{itemize}
\item new \leanfile{tests/\allowbreak{}corpus/\allowbreak{}Examples.lean}, \leanfile{tests/\allowbreak{}corpus/\allowbreak{}Names.lean}, \leanfile{tests/\allowbreak{}corpus/\allowbreak{}NearMiss.lean}: example programs inside and near the verification scope. Each program \texttt{t} named \texttt{\textless{}n\textgreater{}} is erased with \texttt{\#erase t config \{nat := .peano\} to \textquotedbl{}\textless{}n\textgreater{}.peano.ast\textquotedbl{}} and \texttt{\#erase t to \textquotedbl{}\textless{}n\textgreater{}.default.ast\textquotedbl{}}, each readout with the first of these two. In \texttt{NearMiss.lean}, the two \texttt{\#erase}s of \texttt{(\_\allowbreak{} : NM.CNat)} fail and are wrapped in \texttt{\#guard\_\allowbreak{}msgs (error, substring := true)} on \texttt{unknown metavariable} (the metavariable's id depends on what the file elaborates before it).
\item new \leanfile{tests/\allowbreak{}corpus/\allowbreak{}IrrAlias.lean}: \texttt{Irr.useFI := fI two} with \texttt{fI : EndoC} and \texttt{@[irreducible] def EndoC : Type 1 := CNat \ensuremath{\to} CNat}; \texttt{Irr.useLamHR := guardR (lamHR two) six} with \texttt{theorem lamHR : CNat \ensuremath{\to} R := fun \_\allowbreak{} =\textgreater{} hR}; \texttt{Irr.pidHR := pid hR}; where \texttt{axiom R : IProp}, \texttt{axiom hR : R} and \texttt{@[irreducible] def IProp : Type := Prop}. The four \texttt{\#erase}s of \texttt{useFI} and \texttt{useLamHR} fail and are wrapped in \texttt{\#guard\_\allowbreak{}msgs (error)} with their exact errors.
\item new \leanfile{tests/\allowbreak{}corpus/\allowbreak{}UnsafeRec.lean}, over \texttt{CN.\{u\} := (\ensuremath{\alpha} : Type u) \ensuremath{\to} (\ensuremath{\alpha} \ensuremath{\to} \ensuremath{\alpha}) \ensuremath{\to} \ensuremath{\alpha} \ensuremath{\to} \ensuremath{\alpha}}: \texttt{URec.uf one.\{0\}} with \texttt{unsafe def uf (n : CN.\{0\}) : CN.\{0\} := (fun \_\allowbreak{} =\textgreater{} n) (fun x =\textgreater{} uf x)}; \texttt{URec.ua bfalse one.\{0\}} with the mutual block \texttt{unsafe def ua b n := b \_\allowbreak{} (fun m =\textgreater{} m) (fun m =\textgreater{} ub btrue (csucc m)) n} and \texttt{ub} symmetric, over Church booleans in \texttt{Type 2}; and \texttt{URec.ug one.\{0\}} with \texttt{unsafe def ug : CN.\{0\} \ensuremath{\to} CN.\{0\} := (fun \_\allowbreak{} n =\textgreater{} n) (fun x =\textgreater{} ug x)}, whose value is an application.
\item \texttt{doc/\allowbreak{}SHIPPING-CHANGES.md}: this entry; the corpus description above; the reproductions of R-11 and R-14.
\end{itemize}

No file under \texttt{LeanToLambdaBox/}, no package file, no file under \texttt{benchmarks/} or \texttt{scripts/} and no regression test changes.
\item \textbf{Why necessary:} the verification plans a second erasure path for programs without inductive types, to be checked against today's path, and with \texttt{peregrine validate} and \texttt{eval}, on every corpus program of that fragment. Before this change the corpus had 25 such programs (\texttt{Scope.lean}). None of them has an \texttt{@[irreducible]} alias, a \texttt{@[macro\_\allowbreak{}inline]} constant, an opaque, a name that collides or needs escaping, a universe-polymorphic type that is a proposition at universe 0, or recursion. Without the new files, those checks would cover neither the behaviours in which the two paths are expected to differ (an irreducible alias, \texttt{@[macro\_\allowbreak{}inline]}) nor recursion (\texttt{tFix}). The new files add 74 programs of the fragment: 61 in \texttt{Examples.lean}, 7 in \texttt{Names.lean}, 3 in \texttt{IrrAlias.lean} and 3 in \texttt{UnsafeRec.lean}.
\item \textbf{Behaviour before:} the eraser as at S-14. \texttt{scripts/\allowbreak{}corpus.sh} writes 284 files.
\item \textbf{Behaviour after:} the eraser is unchanged. \texttt{scripts/\allowbreak{}corpus.sh} writes 704 files, with no \texttt{PANIC} in the logs of the new files. On the new programs:

\begin{itemize}
\item \texttt{\#erase Irr.useFI} fails with \texttt{function expected} on \texttt{Irr.fI Irr.two}, and \texttt{\#erase Irr.useLamHR} with \texttt{type expected} on \texttt{Irr.R} (R-14). \texttt{Irr.pidHR} is emitted with the axioms \texttt{Irr.R} and \texttt{Irr.hR} (R-14): it validates, and \texttt{peregrine eval} stops with \texttt{Axioms found \ldots{} .Irr.R, .Irr.hR}.
\item \texttt{URec.uf} is a one-member \texttt{tFix}; the program validates and evaluates to \texttt{one}. \texttt{URec.ua} and \texttt{URec.ub} are two declarations, each the two-member \texttt{tFix} of the block (at index 0 and 1); \texttt{ua bfalse one} validates and evaluates to \texttt{csucc one}, a \ensuremath{\lambda}. \texttt{URec.ug} is a \texttt{tFix} whose body is an application; \texttt{peregrine validate} and \texttt{eval} reject the program with \texttt{Fixpoint body is not a lambda} (R-11).
\item The emitted files of \texttt{Examples.lean}, \texttt{Names.lean} and \texttt{NearMiss.lean} are those the scope note reports (erased at S-13), up to hygienic name suffixes and, in \texttt{N\_\allowbreak{}private.*}, the main module's name in \texttt{\_\allowbreak{}private.\textless{}module\textgreater{}.0}.
\end{itemize}
\item \textbf{Effect on emitted .ast (corpus):} the 284 files of S-14 are byte-identical. 420 files are new (210 \texttt{.ast}, 210 \texttt{.ast.inlinings}, no \texttt{.mli}): \texttt{examples/\allowbreak{}Examples/} 340, \texttt{examples/\allowbreak{}Names/} 32, \texttt{examples/\allowbreak{}NearMiss/} 32, \texttt{examples/\allowbreak{}IrrAlias/} 4 (\texttt{pidHR}), \texttt{examples/\allowbreak{}UnsafeRec/} 12.
\item \textbf{Regression test:} the eraser is unchanged, so no test fails before. The new corpus files guard the fragment's behaviours listed above, byte for byte through the corpus, and the failing \texttt{\#erase}s through \texttt{\#guard\_\allowbreak{}msgs}: \texttt{scripts/\allowbreak{}corpus.sh} stops if a corpus file reports an error, so a change of these errors, or an \texttt{\#erase} that starts to succeed, shows. The regression tests are unchanged and pass.
\end{itemize}

\section{\texorpdfstring{S-16: The erasure traversal is total and generic over its backend}{S-16: The erasure traversal is total and generic over its backend}}\label{reg:S-16}

\begin{itemize}
\item \textbf{Commit:} the commit whose subject starts with \texttt{shipping(S-16):} (\texttt{git log -{}-grep='\textasciicircum{}shipping(S-16):'}).
\item \textbf{Files and functions:}

\begin{itemize}
\item \leanfile{LeanToLambdaBox/\allowbreak{}Basic.lean}: \texttt{toBvar} is no longer \texttt{partial}: it is defined by structural recursion, together with the new \texttt{toBvarList}, \texttt{toBvarAlts} and \texttt{toBvarDefs}, which apply it to the arguments of a \texttt{construct}, the alternatives of a \texttt{case} and the definitions of a \texttt{fix} (in place of \texttt{List.map}).
\item \leanfile{LeanToLambdaBox/\allowbreak{}Erasure.lean}:

\begin{itemize}
\item new class \texttt{Backend m}: the environment and \texttt{Meta} operations of the traversal (\texttt{findConst?}, \texttt{unknownConstant}, \texttt{declInfo?}, \texttt{unsafeRecBase?}, \texttt{freshFVarId}, \texttt{instantiate1}, \texttt{isErasable}, \texttt{inferType}, \texttt{casesInfo?}, \texttt{ctorArity?}, \texttt{argMask}, \texttt{isExtern}, \texttt{inlineAttr?}, \texttt{isInstance}, \texttt{prepare}, \texttt{log}, \texttt{outOfFuel}); \texttt{isErasable}, \texttt{inferType} and \texttt{argMask} receive the traversal's \texttt{LocalContext};
\item new \texttt{instance : Backend CoreM}, whose operations are the calls the traversal made before: \texttt{Environment.find?}, \texttt{throwUnknownConstant}, \texttt{Compiler.LCNF.getDeclInfo?}, \texttt{Compiler.isUnsafeRecName?}, \texttt{mkFreshFVarId}, \texttt{Expr.instantiate1}, \texttt{isErasable} and \texttt{Meta.inferType} run by the new \texttt{runMetaM} (the former \texttt{liftMetaM}, in \texttt{CoreM}), \texttt{getCasesInfo?}, \texttt{getCtorArity?}, the new \texttt{argMaskCore} (the argument-mask code of \texttt{register\_\allowbreak{}inductive}, moved unchanged), \texttt{isExtern}, \texttt{Compiler.getInlineAttribute?}, \texttt{Meta.isInstance}, \texttt{prepareErasure} (the former \texttt{prepare\_\allowbreak{}erasure}, which now receives the configuration), \texttt{logInfo}, and \texttt{throwError} for \texttt{outOfFuel};
\item new \texttt{EraseT m := StateT ErasureState (ReaderT ErasureContext m)} in place of \texttt{EraseM} (which was \texttt{EraseT CoreM}); \texttt{run} is generic;
\item new \texttt{getConst}: \texttt{getConstInfo} through the backend (\texttt{findConst?}, then \texttt{unknownConstant});
\item generic over the backend, with bodies otherwise unchanged: \texttt{addAxiom}, \texttt{register\_\allowbreak{}inductive}, \texttt{fvar\_\allowbreak{}to\_\allowbreak{}name}, \texttt{mkLambda}, \texttt{mkLetIn}, \texttt{mkAlt}, \texttt{mkDef}, \texttt{withLocalDecl}, \texttt{withLocalDef}, \texttt{lambdaMonocular}, \texttt{letMonocular}, \texttt{forallMonocular}, \texttt{lambdaMonocularOrIntro}, \texttt{lambdaOrIntroToArity}, \texttt{remove\_\allowbreak{}unsafe\_\allowbreak{}rec}, \texttt{name\_\allowbreak{}occurs};
\item \texttt{withAppEtaToMinArity} is no longer \texttt{partial}: it takes a \texttt{fuel} argument, and its \texttt{go} recurses structurally on it;
\item the functions of the \texttt{where} block of \texttt{erase} (\texttt{visitExpr}, \texttt{visitLiteral}, \texttt{visitLambda}, \texttt{visitLet}, \texttt{visitProj}, \texttt{visitApp}, \texttt{visitConst}, \texttt{visitConstApp}, \texttt{visitConstructor}, \texttt{visitAppArgs}, \texttt{visitCases}, \texttt{visitCasesOn}, \texttt{visitAlt}, \texttt{get\_\allowbreak{}constant\_\allowbreak{}kername}, \texttt{visitMutual}) become the top-level \texttt{mutual} block \leandecl{Erasure.visitExpr}, \ldots{}, generic over the backend. Each takes a first argument \texttt{fuel}, fails with \texttt{Backend.outOfFuel} at 0 and calls the functions of the block with \texttt{fuel - 1}; they are defined by structural recursion on it. Their bodies are otherwise unchanged;
\item new \texttt{travFuel := 2 \textasciicircum{} 32}; \texttt{erase} is no longer \texttt{partial}: it runs \texttt{visitExpr travFuel} after \texttt{prepare} with the \texttt{CoreM} backend, as before.
\end{itemize}
\item new \leanfile{tests/\allowbreak{}regress/\allowbreak{}traversal\_\allowbreak{}total.lean}, its expected outputs \texttt{tests/\allowbreak{}regress/\allowbreak{}expected/\allowbreak{}traversal\_\allowbreak{}total/} (4 files) and \texttt{tests/\allowbreak{}regress/\allowbreak{}expected-peregrine/\allowbreak{}traversal\_\allowbreak{}total/} (2 files).
\item \leanfile{tests/\allowbreak{}regress/\allowbreak{}compiler\_\allowbreak{}api.lean}, \leanfile{tests/\allowbreak{}regress/\allowbreak{}sparse\_\allowbreak{}cases.lean}: the module docstrings name \leandecl{Erasure.visitCases} in place of \texttt{erase.visitCases}.
\item \texttt{doc/\allowbreak{}SHIPPING-CHANGES.md}: this entry; the "Where" fields of R-4, R-5, R-6, R-8, R-9, R-10, R-11, R-12, R-13, R-21 and R-24, and the reproduction of R-4, name the new functions.
\end{itemize}
\item \textbf{Why necessary:} the verification proves a theorem about the traversal that \texttt{\#erase} runs. A \texttt{partial} definition is an opaque constant: it has no equations, so nothing about \texttt{erase}, its \texttt{where} block, \texttt{withAppEtaToMinArity} or \texttt{toBvar} can be proved. The traversal cannot be structural on the term, since \texttt{visitMutual} erases the body of another constant and \texttt{lambdaMonocular} instantiates the body it visits; a fuel argument that bounds the depth of the recursion makes it total without changing the calls it makes. The traversal also ran in \texttt{CoreM}, whose state and \texttt{Meta} operations are opaque to proofs; being generic over the backend lets the verified path run this same traversal with a pure backend (S-19, S-20) instead of a copy of it, while \texttt{\#erase} keeps the \texttt{CoreM} backend.
\item \textbf{Behaviour before:} \texttt{erase}, its \texttt{where} functions, \texttt{withAppEtaToMinArity} and \texttt{toBvar} are \texttt{partial}. Reproduction: \leanfile{tests/\allowbreak{}regress/\allowbreak{}traversal\_\allowbreak{}total.lean} at the parent commit fails with 21 errors (\texttt{Unknown identifier visitExpr}, \texttt{Unknown identifier Backend.outOfFuel}, \ldots{}), and a file stating \texttt{toBvar x 0 (.letIn .anon (.fvar x) (.fvar x)) = .letIn .anon (.bvar 0) (.bvar 1) := rfl} fails with \texttt{Type mismatch}, because \texttt{toBvar} does not unfold.
\item \textbf{Behaviour after:} \texttt{\#erase} gives the same outputs and the same log messages; a \texttt{PANIC} message names the new function (\texttt{PANIC at Erasure.visitLiteral} in place of \texttt{PANIC at Erasure.erase.visitLiteral}). The traversal functions and \texttt{toBvar} have equation lemmas (\leandecl{Erasure.visitExpr.eq_1}: \texttt{visitExpr 0 e = liftM (Backend.outOfFuel \textquotedbl{}visitExpr\textquotedbl{})}; \texttt{visitExpr.eq\_\allowbreak{}2}: at \texttt{fuel + 1} on \texttt{.app}, the body of \texttt{visitExpr}; likewise for the other functions and for \texttt{withAppEtaToMinArity.go}), and \texttt{rfl} computes \texttt{toBvar}. With fuel 2, \texttt{visitExpr} on \texttt{fun (x : Nat) =\textgreater{} x} fails with \texttt{erasure: recursion bound reached in visitExpr}; with fuel 3 it gives \texttt{\ensuremath{\lambda}x. x}. \texttt{\#erase} runs with fuel \texttt{2 \textasciicircum{} 32}, which no program reaches: the recursion is as deep as before and exhausts the stack first (R-6). All axioms of \texttt{visitExpr} and \texttt{visitMutual} are \texttt{propext}, \texttt{Classical.choice} and \texttt{Quot.sound}; \texttt{toBvar} has none.
\item \textbf{Effect on emitted .ast (corpus):} byte-identical for all 704 files; \texttt{scripts/\allowbreak{}corpus-diff.sh} against the corpus of S-15 reports 704 identical. The corpus logs are identical except for the location lines of the 4 \texttt{PANIC} messages of \texttt{Defects.lean} (the function names and line numbers) and their backtraces.
\item \textbf{Regression test:} \leanfile{tests/\allowbreak{}regress/\allowbreak{}traversal\_\allowbreak{}total.lean} states, for every backend, the equations of \texttt{visitExpr} at fuel 0 and at \texttt{fuel + 1} on a free variable, of \texttt{visitAppArgs} and of \texttt{visitMutual} at fuel 0 (proved by \texttt{rw} with the equation lemmas); the value of \texttt{toBvar} under a \ensuremath{\lambda}, a let, a case with two alternatives and a fixpoint with two definitions (proved by \texttt{rfl}); the fuel error and the result at sufficient fuel of \texttt{visitExpr} with the \texttt{CoreM} backend (\texttt{\#guard\_\allowbreak{}msgs}); and it pins the output of a program that goes through every function of the traversal: a let of a literal, a projection, a match with its alternatives, a constructor \ensuremath{\eta}-expanded to its arity (\texttt{List.map T.b}), a mutual recursion (\texttt{ev}/\texttt{od}), under Peano naturals (\texttt{all.ast}, which peregrine evaluates to 10, Lean's value) and under the default configuration (\texttt{go.ast}). It fails before (21 errors; the outputs are the same) and passes after.
\end{itemize}

\section{\texorpdfstring{S-17: The traversal's context lists its locals, and binder names come from that list}{S-17: The traversal's context lists its locals, and binder names come from that list}}\label{reg:S-17}

\begin{itemize}
\item \textbf{Commit:} the commit whose subject starts with \texttt{shipping(S-17):} (\texttt{git log -{}-grep='\textasciicircum{}shipping(S-17):'}).
\item \textbf{Files and functions:}

\begin{itemize}
\item \leanfile{LeanToLambdaBox/\allowbreak{}Erasure.lean}:

\begin{itemize}
\item \texttt{ErasureContext} becomes \texttt{TravCtx}, with the fields \texttt{lctx}, \texttt{locals}, \texttt{fixvars} and \texttt{config}: the new field \texttt{locals : List Local} holds the same binders as \texttt{lctx}, innermost first;
\item new structure \texttt{Local} (\texttt{fvarId}, \texttt{userName}, \texttt{type}, \texttt{value?});
\item new \texttt{binderNameOf}: the name test of \texttt{fvar\_\allowbreak{}to\_\allowbreak{}name}, as a function of the user name (an ASCII graphic name is kept, any other becomes anonymous); new \texttt{fixDefName}: the name that \texttt{mkDef} gave a fixpoint definition;
\item \texttt{withLocalDecl} and \texttt{withLocalDef} push the new binder onto \texttt{locals} (with its value for a \texttt{let}) as well as into \texttt{lctx};
\item \texttt{fvar\_\allowbreak{}to\_\allowbreak{}name} is \texttt{binderNameOf} of the user name of the variable's entry in \texttt{locals}, in place of a lookup in \texttt{lctx}; \texttt{mkDef} names the definition with \texttt{fixDefName};
\item \texttt{Backend.isErasable} and \texttt{Backend.inferType} also receive the list of locals; the \texttt{CoreM} backend ignores it and runs \texttt{Meta} in \texttt{lctx}, as before; \texttt{EraseT m} reads a \texttt{TravCtx}.
\end{itemize}
\item new \leanfile{tests/\allowbreak{}regress/\allowbreak{}traversal\_\allowbreak{}locals.lean}, its expected outputs \texttt{tests/\allowbreak{}regress/\allowbreak{}expected/\allowbreak{}traversal\_\allowbreak{}locals/} (6 files) and \texttt{tests/\allowbreak{}regress/\allowbreak{}expected-peregrine/\allowbreak{}traversal\_\allowbreak{}locals/} (2 files).
\item \leanfile{tests/\allowbreak{}regress/\allowbreak{}traversal\_\allowbreak{}total.lean}: its equation of \texttt{visitExpr} on a free variable passes the list of locals to \texttt{Backend.isErasable}.
\item \texttt{doc/\allowbreak{}SHIPPING-CHANGES.md}: this entry.
\end{itemize}
\item \textbf{Why necessary:} the verified backend's oracle types free variables from a list of locals: a \texttt{LocalContext} lookup goes through a \texttt{PersistentHashMap}, whose operations are opaque to proofs (lean4lean states them as axioms). The erasure relation of the verification fixes the \ensuremath{\lambda}\ensuremath{\Box} name of a binder as a function of its Lean user name (\texttt{binderNameOf}) and the name of a fixpoint definition as a function of the constant's name (\texttt{fixDefName}), so the traversal computes them from these. Keeping \texttt{lctx} leaves the \texttt{CoreM} backend unchanged.
\item \textbf{Behaviour before:} \texttt{fvar\_\allowbreak{}to\_\allowbreak{}name} reads the binder's user name from \texttt{lctx} (\texttt{lctx.fvarIdToDecl.find!}), and the context has no list of locals. Reproduction: \leanfile{tests/\allowbreak{}regress/\allowbreak{}traversal\_\allowbreak{}locals.lean} at the parent commit fails with 9 errors (\texttt{Unknown identifier binderNameOf}, \texttt{Invalid field locals}, \ldots{}); its six output files are those expected.
\item \textbf{Behaviour after:} \texttt{\#erase} gives the same outputs and the same log messages. After \texttt{withLocalDecl a} and \texttt{withLocalDef b}, \texttt{locals} is \texttt{[b (with its value), a]} and \texttt{lctx} holds both. \texttt{binderNameOf} gives \texttt{x \ensuremath{\mapsto} \textquotedbl{}x\textquotedbl{}}, \texttt{\ensuremath{\alpha}\ensuremath{_1} \ensuremath{\mapsto} anonymous}, \texttt{a.b \ensuremath{\mapsto} \textquotedbl{}a.b\textquotedbl{}}, \texttt{\ensuremath{\langle\!\langle}a b\ensuremath{\rangle\!\rangle} \ensuremath{\mapsto} anonymous}, as \texttt{fvar\_\allowbreak{}to\_\allowbreak{}name} did. A variable without an entry in \texttt{locals} makes \texttt{fvar\_\allowbreak{}to\_\allowbreak{}name} panic (\texttt{Option.get!}), as a variable without an entry in \texttt{lctx} did (\texttt{PersistentHashMap.find!}); the traversal introduces every variable it names with \texttt{withLocalDecl} or \texttt{withLocalDef}, and no corpus or regression log has a new \texttt{PANIC}.
\item \textbf{Effect on emitted .ast (corpus):} byte-identical for all 704 files; \texttt{scripts/\allowbreak{}corpus-diff.sh} against the corpus of S-16 reports 704 identical. The corpus logs are identical except for the line numbers in the locations of the 4 \texttt{PANIC} messages of \texttt{Defects.lean}, and their backtraces.
\item \textbf{Regression test:} \leanfile{tests/\allowbreak{}regress/\allowbreak{}traversal\_\allowbreak{}locals.lean} checks \texttt{binderNameOf} and \texttt{fixDefName} on ASCII, non-ASCII and hierarchical names, and the contents of \texttt{locals} and \texttt{lctx} after a \texttt{withLocalDecl} and a \texttt{withLocalDef} (\texttt{\#guard\_\allowbreak{}msgs}); and it pins the output of programs with a binder at every place where the traversal introduces one: a \ensuremath{\lambda} with an ASCII and a non-ASCII binder, a \texttt{let}, a constructor \ensuremath{\eta}-expanded to its arity (\texttt{List.map T.b}), a \texttt{casesOn} with an alternative that is not a \ensuremath{\lambda} (\texttt{g}, \ensuremath{\eta}-expanded through its type) and one that is, a sparse \texttt{casesOn} whose discriminant is let-bound (\texttt{discr}, S-14), a \texttt{Nat} match in machine mode (\texttt{n}), and a recursive definition (the fixpoint definition \texttt{Loc.step}): \texttt{names.ast}, \texttt{step.ast} (default configuration) and \texttt{all.ast} (Peano naturals, which peregrine evaluates to 11, Lean's value). A binder missing from \texttt{locals} would make the test fail on the panic. It fails before (9 errors; the outputs are the same) and passes after.
\end{itemize}

\section{\texorpdfstring{S-18: The dependency closure of a program is computed by \texttt{collectDeps}, which \texttt{\#erase} does not call yet}{S-18: The dependency closure of a program is computed by collectDeps, which \#erase does not call yet}}\label{reg:S-18}

\begin{itemize}
\item \textbf{Commit:} the commit whose subject starts with \texttt{shipping(S-18):} (\texttt{git log -{}-grep='\textasciicircum{}shipping(S-18):'}).
\item \textbf{Files and functions:}

\begin{itemize}
\item new \leanfile{LeanToLambdaBox/\allowbreak{}Erasure/\allowbreak{}Collect.lean} (module \texttt{LeanToLambdaBox.Erasure.Collect}; it imports \texttt{LeanToLambdaBox.Basic} and lean4lean's \texttt{Lean4Lean.Verify.Axioms}):

\begin{itemize}
\item derived \texttt{DecidableEq} instances for \texttt{ModPath} and \texttt{Kername} (\texttt{instDecidableEqModPath}, \texttt{instDecidableEqKername});
\item \leandecl{Erasure.EraseError}, with the constructors \texttt{outOfFragment}, \texttt{nameCollision}, \texttt{fuel}, \texttt{failed};
\item \leandecl{Erasure.EnvView}, with the fields \texttt{find?}, \texttt{isExtern}, \texttt{inlineAttr?}, and \leandecl{Erasure.EnvView.ofEnvironment}, which reads them from an \texttt{Environment} (\texttt{Environment.find?}, \texttt{isExtern}, \texttt{Compiler.getInlineAttribute?});
\item \leandecl{Erasure.findConst}: lookup by name in a list of declarations;
\item \leandecl{Erasure.exprConsts}: the constants of a term, or \texttt{outOfFragment} for a free variable, a metavariable, a universe metavariable (tested with lean4lean's \texttt{Level.hasMVar'}), a literal or a projection;
\item \leandecl{Erasure.declDeps}: the names that an axiom, a definition, a theorem or an opaque depends on (its type, its value, the members \texttt{all} of its block), or \texttt{outOfFragment} for a quotient, an inductive type, a constructor or a recursor;
\item \leandecl{Erasure.closure}: the depth-first closure of a list of names, each name expanded once, with a bound on its steps (\texttt{fuel} when reached); a name the view does not have is \texttt{outOfFragment}, and a view that answers with a declaration of another name is \texttt{failed};
\item \leandecl{Erasure.findCollision}: two declarations of a list with the same kername (\texttt{toKername});
\item \texttt{Erasure.collectFuel := 2 \textasciicircum{} 32} and \texttt{Erasure.collectDeps view e}: the closure of the constants of \texttt{e}, and \texttt{nameCollision} if two of its declarations have the same kername.
\end{itemize}
\item \leanfile{LeanToLambdaBox.lean}: imports the new module.
\item new \leanfile{tests/\allowbreak{}regress/\allowbreak{}collect\_\allowbreak{}deps.lean}, its expected outputs \texttt{tests/\allowbreak{}regress/\allowbreak{}expected/\allowbreak{}collect\_\allowbreak{}deps/} (2 files) and \texttt{tests/\allowbreak{}regress/\allowbreak{}expected-peregrine/\allowbreak{}collect\_\allowbreak{}deps/} (2 files).
\item \texttt{doc/\allowbreak{}SHIPPING-CHANGES.md}: this entry.
\end{itemize}
\item \textbf{Why necessary:} the verified path erases a program over its own environment, the dependency closure of the term, not over the whole Lean environment, and \texttt{collectDeps} computes that closure from a view of the environment. S-21 routes on its result: a program whose closure meets an inductive type, a literal, a projection, a metavariable or an unknown constant (\texttt{outOfFragment}) is to keep the unchanged path. The verification's statements \texttt{collectDeps\_\allowbreak{}spec}, \texttt{collectDeps\_\allowbreak{}sub} and \texttt{collectDeps\_\allowbreak{}not\_\allowbreak{}outOfFragment} are about this function, so it is total and works on lists, which proofs and kernel evaluation see through. It tests universe metavariables with lean4lean's structural \texttt{Level.hasMVar'}: Lean's \texttt{Level.hasMVar} reads a cached field, which lean4lean relates to \texttt{hasMVar'} only by an axiom (\texttt{Level.hasMVar\_\allowbreak{}eq}, \leanloc{Lean4Lean/\allowbreak{}Verify/\allowbreak{}Axioms.lean}{289}). The collision check makes the kernames of the closure distinct, since \texttt{toKername} is not injective (R-3). \texttt{findCollision} compares kernames with the derived \texttt{BEq}; the \texttt{DecidableEq} instances decide the equality of kernames that the verification states (\texttt{KernameInj}).
\item \textbf{Behaviour before:} none of these declarations exists. Reproduction: \leanfile{tests/\allowbreak{}regress/\allowbreak{}collect\_\allowbreak{}deps.lean} at the parent commit fails with 21 errors (\texttt{Unknown identifier collectDeps}, \texttt{Unknown identifier EnvView.ofEnvironment}, \texttt{The expected type EnvView is not an inductive type}, \ldots{}); its two output files are those expected.
\item \textbf{Behaviour after:} \texttt{\#erase} is unchanged: nothing calls the new functions. Over the elaboration environment (\texttt{EnvView.ofEnvironment}), \texttt{collectDeps} gives \texttt{[A, CN, one]} for \texttt{one (fun a : A =\textgreater{} a)} (\texttt{axiom A : Type}, \texttt{def CN := (A \ensuremath{\to} A) \ensuremath{\to} A \ensuremath{\to} A}, \texttt{def one : CN := fun s z =\textgreater{} s z}); \texttt{[one, ub, ua, A, CN, useU]} for a constant \texttt{useU} that uses a two-member \texttt{unsafe} block \texttt{ua}/\texttt{ub}, so the block is closed under its members; \texttt{outOfFragment \textquotedbl{}level metavariable\textquotedbl{}} for \texttt{@pid} at a universe metavariable; \texttt{outOfFragment \textquotedbl{}literal\textquotedbl{}} for \texttt{\ensuremath{\langle\!\langle}a b\ensuremath{\rangle\!\rangle}}, whose closure contains \texttt{a\_\allowbreak{}u32b} (the same kername) and a literal, so the scan of the fragment comes before the collision check; and \texttt{nameCollision c\_\allowbreak{}u32d \ensuremath{\langle\!\langle}c d\ensuremath{\rangle\!\rangle}} for an in-fragment closure with those two constants. The kernel evaluates \texttt{collectDeps} with its bound \texttt{2 \textasciicircum{} 32} (\texttt{decide}). No new declaration depends on an axiom of lean4lean: \texttt{exprConsts}, \texttt{declDeps}, \texttt{closure}, \texttt{findConst} and the two instances depend on no axiom, \texttt{findCollision}, \texttt{collectDeps} and \texttt{EnvView.ofEnvironment} on \texttt{propext}, \texttt{Classical.choice} and \texttt{Quot.sound}. Every file that imports \texttt{LeanToLambdaBox} now also imports \texttt{Lean4Lean.Verify.Axioms} and, through it, 25 modules of \texttt{batteries} and 2 more of \texttt{lean4lean} (28 modules): their declarations, including lean4lean's axioms and \texttt{@[simp]} axioms about \texttt{Expr} and \texttt{Level}, are in scope. For every constant declared outside these 28 modules, the \texttt{@[inline]}-family attribute, \texttt{@[extern]}, \texttt{@[implemented\_\allowbreak{}by]}, the reducibility, the instance status and the \texttt{@[csimp]} replacement are the same with and without them (64119 constants with a value other than the default, all equal); the 16 \texttt{@[csimp]} lemmas that the 28 modules add all replace functions declared in \texttt{batteries} (\texttt{Batteries.Data.List.Basic}, e.g. \texttt{List.sublists}). The first \texttt{lake lean} of the benchmark pipeline in \texttt{benchmarks/\allowbreak{}via\_\allowbreak{}malfunction} builds the 28 modules once.
\item \textbf{Effect on emitted .ast (corpus):} byte-identical for all 704 files; \texttt{scripts/\allowbreak{}corpus-diff.sh} against the corpus of S-17 reports 704 identical. The corpus logs are identical except for the job counts of \texttt{lake} (6 jobs become 37 for the build of the package, 21 become 52 for each \texttt{lake lean} of the benchmarks) and the addresses in the backtraces of the 4 \texttt{PANIC} messages of \texttt{Defects.lean}. In a benchmark workspace where the 28 modules are not built yet, the first benchmark log also lists their build.
\item \textbf{Regression test:} \leanfile{tests/\allowbreak{}regress/\allowbreak{}collect\_\allowbreak{}deps.lean} proves by \texttt{decide} that \texttt{collectDeps} gives \texttt{[A, CN, one]} for \texttt{one (fun a : A =\textgreater{} a)} over a view built by hand; checks with \texttt{\#guard\_\allowbreak{}msgs} its results over the elaboration environment on that program, on the \texttt{unsafe} block, on \texttt{@pid} at a universe metavariable, on the collision-first program and on the in-fragment collision (as above); proves by \texttt{decide}, over views built by hand, the closure \texttt{[A, B]} of \texttt{B := A}, that the collision-first program is \texttt{outOfFragment}, that the in-fragment collision is \texttt{nameCollision}, that a constant missing from the view is \texttt{outOfFragment}, and that a view answering \texttt{find? a} with a declaration named \texttt{b} gives \texttt{failed}; and pins the output of \texttt{\#erase} on \texttt{one (fun a : A =\textgreater{} a)} (\texttt{nv1.ast}, which peregrine validates and evaluates to \texttt{\ensuremath{\lambda}z. (\ensuremath{\lambda}a. a) z}). It fails before (21 errors; the outputs are the same) and passes after.
\end{itemize}

\section{\texorpdfstring{S-19: A pure erasability oracle and the monad of a pure backend, which \texttt{\#erase} does not use yet}{S-19: A pure erasability oracle and the monad of a pure backend, which \#erase does not use yet}}\label{reg:S-19}

\begin{itemize}
\item \textbf{Commit:} the commit whose subject starts with \texttt{shipping(S-19):} (\texttt{git log -{}-grep='\textasciicircum{}shipping(S-19):'}).
\item \textbf{Files and functions:}

\begin{itemize}
\item new \leanfile{LeanToLambdaBox/\allowbreak{}Erasure/\allowbreak{}Pure.lean} (module \texttt{LeanToLambdaBox.Erasure.Pure}; it imports \texttt{LeanToLambdaBox.Erasure} and \texttt{LeanToLambdaBox.Erasure.Collect}):

\begin{itemize}
\item the oracle, in namespace \texttt{Erasure.Pure}:

\begin{itemize}
\item \texttt{instLevel ps us}: the level with each parameter of \texttt{ps} replaced by the level at the same position of \texttt{us}, without normalization; \texttt{instLevels ps us}: the same on every sort and constant of a term;
\item \texttt{Ctx}, whose only field \texttt{decls : List ConstantInfo} holds the declarations the oracle reads;
\item \texttt{alwaysZero}: a level that is zero under every assignment of its parameters, read off its shape (\texttt{zero}, \texttt{max} of two such levels, \texttt{imax} whose right side is one);
\item \texttt{findLocal}: the local of a free variable in a list of \leandecl{Erasure.Local}s;
\item \texttt{whnf cx fuel ls e}: the weak-head normal form of \texttt{e} by head \ensuremath{\beta} (with lean4lean's \texttt{Expr.instantiate1'}), \ensuremath{\zeta} (a \texttt{let}, a local of \texttt{ls} with a value), \texttt{mdata} removal, and \ensuremath{\delta} of every definition (\texttt{defnInfo}) of \texttt{cx.decls} used at its number of universe levels, whether or not it is \texttt{@[irreducible]};
\item \texttt{inferType cx fuel ls \ensuremath{\Gamma} e}: the type of \texttt{e}, inferred without checking. Inside \texttt{e} bound variables stay de Bruijn indices: \texttt{\ensuremath{\Gamma}} holds the types of the binders entered, and \texttt{bvar i} has type \texttt{\ensuremath{\Gamma}[i]} lifted by \texttt{i + 1} (lean4lean's \texttt{Expr.liftLooseBVars'}). A free variable has its type in \texttt{ls}, a constant the type of its declaration at the given levels. The type of the function of an application is reduced by \texttt{whnf} to a \ensuremath{\Pi}, and the types of the domain and codomain of a \ensuremath{\Pi} to sorts;
\item \texttt{isArity cx fuel ls T}: \texttt{whnf} of \texttt{T} is a sort, or a \ensuremath{\Pi} whose codomain is an arity;
\item \texttt{isErasable cx fuel ls e}: the type \texttt{T} of \texttt{e} is an arity, or the type of \texttt{T} reduces to a sort that is \texttt{alwaysZero}.
\end{itemize}

Each of \texttt{whnf}, \texttt{inferType}, \texttt{isArity} recurses on its fuel and fails with \texttt{EraseError.fuel} when it runs out. A literal, a projection, a metavariable, a loose bound variable, or a constant or free variable the oracle does not know gives \texttt{outOfFragment}; a constant at the wrong number of levels, an application whose function type does not reduce to a \ensuremath{\Pi}, and a \ensuremath{\Pi} whose domain or codomain type does not reduce to a sort give \texttt{failed};
\item \texttt{Erasure.oracleFuel := 2 \textasciicircum{} 20}, the fuel of one oracle call;
\item \leandecl{Erasure.PureCtx} (fields \texttt{decls : List ConstantInfo} and \texttt{view : EnvView}), \leandecl{Erasure.PureState} (field \texttt{next : Nat}) and \texttt{Erasure.PureM := ReaderT PureCtx (StateT PureState (Except EraseError))}, the monad of a backend without \texttt{Meta}; \texttt{Erasure.EraseT.runPure x st tc pc ps} runs an action \texttt{x} of \texttt{EraseT PureM} from the traversal's state \texttt{st} and context \texttt{tc} and the backend's context \texttt{pc} and state \texttt{ps};
\item operations of \texttt{PureM}: \leandecl{Erasure.PureM.findConst?} (\texttt{findConst} on \texttt{decls}), \leandecl{Erasure.PureM.freshFVarId} (the free variable \texttt{\_\allowbreak{}pure.\textless{}next\textgreater{}}, then \texttt{next + 1}), \leandecl{Erasure.PureM.instantiate1} (\texttt{Expr.instantiate1'}), \texttt{Erasure.PureM.isErasable ls e} (\texttt{Pure.isErasable \ensuremath{\langle}decls\ensuremath{\rangle} oracleFuel ls e}, whose error it throws), \leandecl{Erasure.PureM.casesInfo?} and \leandecl{Erasure.PureM.ctorArity?} (always \texttt{none}).
\end{itemize}
\item \leanfile{LeanToLambdaBox.lean}: imports the new module.
\item new \leanfile{tests/\allowbreak{}regress/\allowbreak{}pure\_\allowbreak{}oracle.lean}, its expected outputs \texttt{tests/\allowbreak{}regress/\allowbreak{}expected/\allowbreak{}pure\_\allowbreak{}oracle/} (2 files) and \texttt{tests/\allowbreak{}regress/\allowbreak{}expected-peregrine/\allowbreak{}pure\_\allowbreak{}oracle/} (1 file).
\item \texttt{doc/\allowbreak{}SHIPPING-CHANGES.md}: this entry.
\end{itemize}
\item \textbf{Why necessary:} the erasability test of \texttt{\#erase}, \leandecl{Erasure.isErasable}, runs \texttt{Meta.inferType}, \texttt{Meta.isProp} and \texttt{Meta.isTypeFormerType} in \texttt{MetaM}, whose state and operations are opaque to proofs, so its answers cannot be proved sound. The verification proves that an "erasable" answer of the oracle is sound (\texttt{Pure.isErasable\_\allowbreak{}sound}), which needs the oracle to be a total function on data: lists of declarations and locals, fuel, and lean4lean's \texttt{Expr.instantiate1'} and \texttt{Expr.liftLooseBVars'}, which are definitions (Lean's \texttt{Expr.instantiate1} and \texttt{Expr.liftLooseBVars} are related to them only by lean4lean's axioms \texttt{Expr.instantiate1\_\allowbreak{}eq} and \texttt{Expr.liftLooseBVars\_\allowbreak{}eq}). The kernel can evaluate it (\texttt{decide}). Its reductions unfold every definition, as the reductions of MetaRocq's \texttt{is\_\allowbreak{}erasableb} (\texttt{erasure/\allowbreak{}theories/\allowbreak{}ErasureFunction.v:894}), at \texttt{RedFlags.default}, unfold every constant with a body. \texttt{PureM}, its operations and \texttt{EraseT.runPure} are the backend with which the verified path is to run the traversal of S-16 (S-20 completes the backend), and the verification's statements about that run are stated with \texttt{EraseT.runPure}. The \texttt{Meta} oracle \leandecl{Erasure.isErasable} is unchanged, and \texttt{\#erase} keeps using it.
\item \textbf{Behaviour before:} none of these declarations exists. Reproduction: \leanfile{tests/\allowbreak{}regress/\allowbreak{}pure\_\allowbreak{}oracle.lean} at the parent commit fails with 49 errors (\texttt{Unknown constant Pure.isErasable}, \texttt{Unknown identifier PureM.isErasable}, \texttt{Unknown constant Pure.Ctx}, \ldots{}); its two output files are those expected.
\item \textbf{Behaviour after:} \texttt{\#erase} is unchanged: nothing calls the new functions. The kernel evaluates the oracle at \texttt{oracleFuel} (\texttt{decide}), on environments built by hand: the ill-typed spine \texttt{hq A a} of a proof \texttt{hq : Q}, where \texttt{Q : Prop := \ensuremath{\forall} P : Prop, P \ensuremath{\to} P}, \texttt{a : A} and \texttt{A : Type}, is kept; the proof \texttt{hq.\{v\} : P.\{v\}} of \texttt{P.\{v\} : Sort v} is kept at the parameter \texttt{v} and erased at level \texttt{0}; with \texttt{IProp : Type := Prop}, \texttt{R : IProp}, \texttt{hR : R}, \texttt{Endo : Type := A \ensuremath{\to} A}, \texttt{fI : Endo} and \texttt{a : A}, the oracle types and keeps \texttt{fun (\_\allowbreak{} : R) (x : A) =\textgreater{} x} and \texttt{fI a}, and erases \texttt{R} and \texttt{hR}. With fuel 1 the oracle fails on \texttt{fI a} with \texttt{fuel \textquotedbl{}inferType\textquotedbl{}}, and with fuel 8 it keeps it. \texttt{PureM.freshFVarId} from the counter 3 gives \texttt{\_\allowbreak{}pure.3}, then \texttt{\_\allowbreak{}pure.4}, and leaves 5. Over the declarations that \texttt{collectDeps} collects from the elaboration environment for \texttt{pidHR : R := pid hR}, with \texttt{IProp} made \texttt{@[irreducible]} after \texttt{R} and \texttt{hR} are declared, the oracle erases \texttt{R}, \texttt{hR} and \texttt{pidHR} and keeps \texttt{pid.\{1\}}, while \texttt{\#erase Irr.pidHR} keeps \texttt{R} and \texttt{hR} (R-14). \texttt{instLevel}, \texttt{Ctx}, \texttt{findLocal}, \texttt{oracleFuel}, \texttt{PureCtx}, \texttt{PureState}, \texttt{PureM} and the operations of \texttt{PureM} other than \texttt{isErasable} depend on no axiom; \texttt{instLevels}, \texttt{alwaysZero}, \texttt{whnf}, \texttt{inferType}, \texttt{isArity}, \texttt{Pure.isErasable} and \texttt{PureM.isErasable} on \texttt{propext}; \texttt{EraseT.runPure} on \texttt{propext}, \texttt{Classical.choice} and \texttt{Quot.sound}, as \texttt{EraseT} does (the hash maps of \texttt{ErasureState} and \texttt{TravCtx}). None depends on an axiom of lean4lean. \texttt{whnf}, \texttt{inferType}, \texttt{isArity}, \texttt{alwaysZero}, \texttt{instLevel} and \texttt{instLevels} have equation lemmas.
\item \textbf{Effect on emitted .ast (corpus):} byte-identical for all 704 files; \texttt{scripts/\allowbreak{}corpus-diff.sh} against the corpus of S-18 reports 704 identical. The corpus logs are identical except for the job counts of \texttt{lake} (37 jobs become 38 for the build of the package, 52 become 53 for each \texttt{lake lean} of the benchmarks) and the addresses in the backtraces of the 4 \texttt{PANIC} messages of \texttt{Defects.lean}.
\item \textbf{Regression test:} \leanfile{tests/\allowbreak{}regress/\allowbreak{}pure\_\allowbreak{}oracle.lean} proves by \texttt{decide}, at \texttt{oracleFuel}, the oracle's answers above on the three environments built by hand (the environments and terms of the verification's register tests \texttt{defHead\_\allowbreak{}kept}, \texttt{levelDependent\_\allowbreak{}kept} and \texttt{irreducibleAlias}); proves by \texttt{decide} that fuel 1 and fuel 0 give \texttt{fuel} and fuel 8 an answer, and pins the messages with \texttt{\#guard\_\allowbreak{}msgs}; proves by \texttt{decide} the results of \texttt{PureM.findConst?}, \texttt{casesInfo?}, \texttt{ctorArity?} and \texttt{isErasable} (its answers, the error it throws on an unknown free variable, and the types it reads from a list of locals, one of them with a value), and by \texttt{rfl} one instance of \texttt{PureM.instantiate1}; pins with \texttt{\#guard\_\allowbreak{}msgs} \texttt{PureM.freshFVarId} and a run of \texttt{EraseT.runPure}; pins with \texttt{\#guard\_\allowbreak{}msgs} the oracle's answers over the elaboration environment on \texttt{Irr.R}, \texttt{Irr.hR}, \texttt{Irr.pidHR} and \texttt{Irr.pid.\{1\}}; and pins the output of \texttt{\#erase} on \texttt{Irr.pidHR} (\texttt{pidHR.ast}, which peregrine validates). It fails before (49 errors; the outputs are the same) and passes after.
\end{itemize}

\section{\texorpdfstring{S-20: The pure backend is complete, and \texttt{erasePure} runs the traversal with it, which \texttt{\#erase} does not call yet}{S-20: The pure backend is complete, and erasePure runs the traversal with it, which \#erase does not call yet}}\label{reg:S-20}

\begin{itemize}
\item \textbf{Commit:} the commit whose subject starts with \texttt{shipping(S-20):} (\texttt{git log -{}-grep='\textasciicircum{}shipping(S-20):'}).
\item \textbf{Files and functions:}

\begin{itemize}
\item \leanfile{LeanToLambdaBox/\allowbreak{}Erasure/\allowbreak{}Pure.lean}:

\begin{itemize}
\item new operations of \texttt{PureM}: \leandecl{Erasure.PureM.declInfo?} (\texttt{findConst?}: the declaration itself, where the \texttt{CoreM} backend's \texttt{Compiler.LCNF.getDeclInfo?} prefers an \texttt{\_\allowbreak{}unsafe\_\allowbreak{}rec} version), \leandecl{Erasure.PureM.unsafeRecBase?} (always \texttt{none}) and \leandecl{Erasure.PureM.prepare} (returns its term);
\item new instance \texttt{Erasure.instBackendPureM : Backend PureM}. \texttt{findConst?}, \texttt{declInfo?}, \texttt{freshFVarId}, \texttt{instantiate1}, \texttt{casesInfo?}, \texttt{ctorArity?} and \texttt{prepare} are the operations of \texttt{PureM}, and \texttt{isErasable} is \texttt{PureM.isErasable} on the traversal's locals; \texttt{unsafeRecBase?}, a field without the monad, is \texttt{none}; \texttt{inferType} is \texttt{Pure.inferType} at \texttt{oracleFuel} on the locals, whose error it throws; \texttt{argMask} keeps every field; \texttt{isExtern} and \texttt{inlineAttr?} read the view; \texttt{isInstance} is \texttt{false}; \texttt{log} does nothing; \texttt{unknownConstant n} throws \texttt{outOfFragment \textquotedbl{}unknown constant n\textquotedbl{}}, and \texttt{outOfFuel site} throws \texttt{fuel site};
\item new \texttt{Erasure.erasePure view cfg decls e}: runs \texttt{visitExpr travFuel e} with this backend (\texttt{EraseT.runPure} from the empty traversal state, the configuration \texttt{cfg}, the backend context \texttt{\ensuremath{\langle}decls, view\ensuremath{\rangle}} and the counter 0), and assembles the program \texttt{.untyped gdecls (some t)} and the inlinings as \texttt{erase} does, without \texttt{prepare};
\item new \leandecl{Erasure.nameOccurs} (the traversal's \texttt{name\_\allowbreak{}occurs} with exact names), \leandecl{Erasure.isRecursiveDecl} (the block \texttt{all} does not have exactly one member, or the value mentions the constant's own name) and \leandecl{Erasure.axiomatized} (a single declaration tagged \texttt{@[extern]}, under \texttt{extern := .preferAxiom});
\item the module docstring.
\end{itemize}
\item new \leanfile{tests/\allowbreak{}regress/\allowbreak{}pure\_\allowbreak{}backend.lean}, its expected outputs \texttt{tests/\allowbreak{}regress/\allowbreak{}expected/\allowbreak{}pure\_\allowbreak{}backend/} (20 files) and \texttt{tests/\allowbreak{}regress/\allowbreak{}expected-peregrine/\allowbreak{}pure\_\allowbreak{}backend/} (7 files).
\item new \leanfile{tests/\allowbreak{}harness/\allowbreak{}PureHarness.lean} (the harness), \texttt{tests/\allowbreak{}harness/\allowbreak{}expected/\allowbreak{}pure-harness.txt} (its expected summary) and \texttt{scripts/\allowbreak{}pure-harness.sh}.
\item \texttt{.github/\allowbreak{}workflows/\allowbreak{}build.yml}: new step "Pure path next to the Meta path on the corpus".
\item \texttt{doc/\allowbreak{}SHIPPING-CHANGES.md}: this entry, and the section "Regression tests".
\end{itemize}
\item \textbf{Why necessary:} the verification proves its theorem about \texttt{erasePure}: the traversal of \texttt{\#erase} run with a backend whose every operation is a function on data. With the operations of S-19 and those added here, \texttt{PureM} is an instance of \texttt{Backend}, so \texttt{erasePure} runs the traversal of S-16 itself, not a copy of it. On this backend the traversal erases kernel values, as MetaRocq's \texttt{erase} erases the kernel term (\texttt{erasure/\allowbreak{}theories/\allowbreak{}ErasureFunction.v:1309}, \texttt{erase\_\allowbreak{}constant\_\allowbreak{}body} on \texttt{cst\_\allowbreak{}body}): \texttt{prepare} is the identity, since \texttt{macro\_\allowbreak{}inline} and matcher inlining give a term that is not an erasure of the kernel value, and csimp rests on an \texttt{Eq} theorem; \texttt{declInfo?} finds the declaration itself and \texttt{unsafeRecBase?} is \texttt{none}, since a user constant \texttt{x.\_\allowbreak{}unsafe\_\allowbreak{}rec} is unrelated to \texttt{x}. \texttt{nameOccurs}, \texttt{isRecursiveDecl} and \texttt{axiomatized} are the tests of \texttt{visitMutual} at this backend, which the verification's specification reads (\texttt{RecursiveDecl}, and the \texttt{@[extern]} remapping of the evaluation environment). The harness is the gate of this entry: on the corpus programs in the fragment, the outputs and oracle answers of the pure path equal those of the \texttt{Meta} path, except for the differences listed under "Behaviour after".
\item \textbf{Behaviour before:} none of these declarations exists. Reproduction: \leanfile{tests/\allowbreak{}regress/\allowbreak{}pure\_\allowbreak{}backend.lean} at the parent commit fails with 38 errors (\texttt{Unknown identifier erasePure}, \texttt{Unknown constant nameOccurs}, failed synthesis of \texttt{Backend PureM}, \ldots{}); its eight \texttt{\#erase} outputs are those expected. The harness does not compile (\texttt{erasePure} is unknown).
\item \textbf{Behaviour after:} \texttt{\#erase} is unchanged: nothing it runs calls the new declarations.

\begin{itemize}
\item At the pure backend the traversal's \texttt{name\_\allowbreak{}occurs} is \texttt{nameOccurs} by definition: \texttt{\ensuremath{\forall} n e, name\_\allowbreak{}occurs (m := PureM) n e = nameOccurs n e} holds by \texttt{rfl}. The kernel unfolds both structural recursions; the elaborator needs \texttt{smartUnfolding} off, or \texttt{delta name\_\allowbreak{}occurs nameOccurs} first. So \texttt{visitMutual}'s test \texttt{single\_\allowbreak{}decl \&\& !name\_\allowbreak{}occurs name value} is \texttt{!isRecursiveDecl ci} for the declaration \texttt{ci} that \texttt{declInfo? name} finds.
\item Over the elaboration environment, \texttt{isRecursiveDecl} holds for the \texttt{unsafe} recursive \texttt{uf} and not for \texttt{one}, \texttt{oneU} and the \texttt{@[extern]} definition \texttt{ext}; \texttt{axiomatized} holds for \texttt{ext} under the default configuration only.
\item \texttt{erasePure}, after \texttt{collectDeps}, gives the output of \texttt{\#erase} on \texttt{one (fun a : A =\textgreater{} a)} (NV-1; peregrine evaluates it to \texttt{\ensuremath{\lambda}z. (\ensuremath{\lambda}a. a) z}) and on the unsafe recursion \texttt{uf one.\{0\}} (a fixpoint; peregrine evaluates it to the value of \texttt{oneU}). It keeps the call \texttt{mfirst one one} that \texttt{\#erase} replaces by \texttt{one} (\texttt{@[macro\_\allowbreak{}inline] mfirst}). It erases \texttt{pidHR : R := pid hR}, whose type \texttt{R : IProp} is a proposition only through the \texttt{@[irreducible]} alias \texttt{IProp := Prop}, to \texttt{(Untyped () (Some tBox))}, where \texttt{\#erase} keeps it. It makes the \texttt{@[extern]} definition \texttt{ext} an axiom under the default configuration and keeps its body under \texttt{extern := .preferLogical}. It logs nothing, where \texttt{\#erase} logs, e.g., \texttt{No value found for name PB.R, emitting axiom.} A constant missing from the declarations fails with \texttt{outOfFragment \textquotedbl{}unknown constant\textquotedbl{}} (from the oracle), an ill-typed application with \texttt{failed \textquotedbl{}not a function\textquotedbl{}}, and \texttt{visitExpr} at fuel 1 with \texttt{fuel \textquotedbl{}visitApp\textquotedbl{}}.
\item The harness on the corpus: of 292 \texttt{\#erase} commands, 119 are out of the fragment, 2 have a \texttt{collectDeps} error, and 171 are in the fragment (25 of \texttt{Scope.lean}, 122 of \texttt{Examples.lean}, 12 of \texttt{Names.lean}, 6 of \texttt{IrrAlias.lean}, 6 of \texttt{UnsafeRec.lean}; each program of the last four files is erased under two configurations). The output of the pure path equals the output of the \texttt{Meta} path byte for byte on 159 of them, and on every in-fragment program except these:

\begin{itemize}
\item \texttt{attrEx} (both configurations): \texttt{\#erase} inlines the \texttt{@[macro\_\allowbreak{}inline]} \texttt{mfirst}, so its \texttt{attrEx} is \texttt{inlTwo}; the pure path keeps \texttt{mfirst inlTwo three} and declares \texttt{mfirst}, \texttt{three} and \texttt{add} in addition;
\item \texttt{irrAxiomArg}, \texttt{irrThmArg}: the proof arguments \texttt{hR}, \texttt{hR2} become \texttt{\ensuremath{\Box}}, and their declarations disappear;
\item \texttt{pidHR}: \texttt{(Untyped () (Some tBox))};
\item \texttt{useFI}, \texttt{useLamHR}: \texttt{\#erase} fails (\texttt{function expected}, \texttt{type expected}), and the pure path gives an output;
\item \texttt{N\_\allowbreak{}collision} (both configurations): \texttt{collectDeps} reports \texttt{nameCollision Nm.two\_\allowbreak{}u39 Nm.two'}.
\end{itemize}

The 6 outputs of \texttt{UnsafeRec.lean} are byte-identical on both paths. The pure oracle never fails. Of its 5735 calls, 10 get an answer that differs from the \texttt{Meta} oracle's: \texttt{Irr.fI Irr.two} in \texttt{useFI} (the pure oracle keeps it; \texttt{Meta} fails with \texttt{function expected}), and \texttt{hR} in \texttt{irrAxiomArg}, \texttt{hR2} in \texttt{irrThmArg}, \texttt{Irr.pidHR} and \texttt{Irr.lamHR Irr.two} in \texttt{useLamHR} (the pure oracle erases them; \texttt{Meta} keeps them), each under both configurations: the pure oracle reduces with the kernel's \ensuremath{\delta}, which unfolds the \texttt{@[irreducible]} aliases \texttt{EndoC} and \texttt{IProp}, and the \texttt{Meta} oracle, at default transparency, does not. On the 67 in-fragment programs of \texttt{Examples.lean} and \texttt{Names.lean} the pure run makes 2351 oracle calls under each configuration. The files that \texttt{\#erase} writes during the harness run are byte-identical to the corpus.
\item Axioms: \texttt{PureM.declInfo?}, \texttt{PureM.unsafeRecBase?} and \texttt{PureM.prepare} depend on none; \texttt{nameOccurs} and \texttt{axiomatized} on \texttt{propext}; \texttt{isRecursiveDecl}, the instance and \texttt{erasePure} on \texttt{propext}, \texttt{Classical.choice} and \texttt{Quot.sound}. None depends on an axiom of lean4lean.
\end{itemize}
\item \textbf{Effect on emitted .ast (corpus):} byte-identical for all 704 files; \texttt{scripts/\allowbreak{}corpus-diff.sh} against the corpus of S-19 reports 704 identical. The corpus logs are identical except for the addresses in the backtraces of the 4 \texttt{PANIC} messages of \texttt{Defects.lean}.
\item \textbf{Regression test:} \leanfile{tests/\allowbreak{}regress/\allowbreak{}pure\_\allowbreak{}backend.lean} proves by \texttt{rfl} the operations of the instance (\texttt{declInfo?}, \texttt{findConst?}, \texttt{unsafeRecBase?}, \texttt{remove\_\allowbreak{}unsafe\_\allowbreak{}rec}, \texttt{prepare}, \texttt{isInstance}, \texttt{log}, \texttt{instantiate1}, \texttt{isErasable}) and that \texttt{name\_\allowbreak{}occurs} is \texttt{nameOccurs} at the pure backend (pointwise and as functions); proves that \texttt{visitMutual}'s test is \texttt{isRecursiveDecl}; pins with \texttt{\#guard\_\allowbreak{}msgs} \texttt{isRecursiveDecl} and \texttt{axiomatized} on declarations of the elaboration environment, the errors of \texttt{erasePure} above, and the results of \texttt{unknownConstant}, \texttt{outOfFuel} and \texttt{argMask}; and, with a command \texttt{\#erase\_\allowbreak{}pure} defined in the test (\texttt{collectDeps} and \texttt{erasePure} on a term elaborated as \texttt{\#erase} elaborates it), pins next to the output of \texttt{\#erase} the outputs of the pure path on the programs above (\texttt{nv1}, \texttt{ufOne}, \texttt{useM}, \texttt{pidHR}, \texttt{useExt}, \texttt{useExtL}), which peregrine validates, evaluating \texttt{nv1} and \texttt{ufOne}. It fails before (38 errors; the outputs of \texttt{\#erase} are the same) and passes after. \texttt{scripts/\allowbreak{}pure-harness.sh} pins the summary of the comparison above; before, the harness does not compile.
\end{itemize}

\section{\texorpdfstring{S-21: \texttt{\#erase} erases programs in the fragment with the pure path}{S-21: \#erase erases programs in the fragment with the pure path}}\label{reg:S-21}

\begin{itemize}
\item \textbf{Commit:} the commit whose subject starts with \texttt{shipping(S-21):} (\texttt{git log -{}-grep='\textasciicircum{}shipping(S-21):'}).
\item \textbf{Files and functions:}

\begin{itemize}
\item new \leanfile{LeanToLambdaBox/\allowbreak{}Erasure/\allowbreak{}Entry.lean}:

\begin{itemize}
\item \leandecl{Erasure.Route} (\texttt{viaPure r}, \texttt{viaMeta});
\item \texttt{Erasure.route view cfg e}: \texttt{collectDeps view e = ok decls} gives \texttt{viaPure (erasePure view cfg decls e)}, \texttt{outOfFragment} gives \texttt{viaMeta}, and every other error \texttt{err} gives \texttt{viaPure (error err)};
\item \texttt{Erasure.eraseEntry view cfg e : CoreM (Program \ensuremath{\times} List Kername)}: \texttt{viaPure (ok r)} returns \texttt{r}, \texttt{viaPure (error err)} throws \texttt{erasure failed: \textless{}err.describe\textgreater{}}, \texttt{viaMeta} runs \texttt{erase e cfg};
\item \leandecl{Erasure.EraseError.describe}, the text of an error in that message.
\end{itemize}
\item new \leanfile{LeanToLambdaBox/\allowbreak{}Erasure/\allowbreak{}Command.lean}: the \texttt{\#erase} syntax \leandecl{Erasure.erasestx} and its elaborator \leandecl{Erasure.eraseElab}, moved from \leanfile{LeanToLambdaBox/\allowbreak{}Erasure.lean}. The moved code is unchanged except for one step: \texttt{erase e cfg} becomes \texttt{eraseEntry view cfg e}, with \texttt{view := EnvView.ofEnvironment (\ensuremath{\leftarrow} getEnv)}.
\item \leanfile{LeanToLambdaBox/\allowbreak{}Erasure.lean}: the syntax and the elaborator are removed.
\item \leanfile{LeanToLambdaBox.lean}: imports the two new modules.
\item corpus sources:

\begin{itemize}
\item \leanfile{tests/\allowbreak{}corpus/\allowbreak{}Scope.lean} gains \texttt{\#erase @Scope.pid.\{1\}} (\texttt{pidU1.ast}) and \texttt{\#erase @Scope.pcomp.\{1,1,1\}} (\texttt{pcompU111.ast});
\item \leanfile{tests/\allowbreak{}corpus/\allowbreak{}IrrAlias.lean} loses the \texttt{\#guard\_\allowbreak{}msgs} around the four \texttt{\#erase} of \texttt{useFI} and \texttt{useLamHR};
\item \leanfile{tests/\allowbreak{}corpus/\allowbreak{}Names.lean} wraps the two \texttt{\#erase Nm.both} (\texttt{N\_\allowbreak{}collision}) in \texttt{\#guard\_\allowbreak{}msgs}.
\end{itemize}
\item new \leanfile{tests/\allowbreak{}regress/\allowbreak{}entry\_\allowbreak{}route.lean}, with its expected outputs \texttt{tests/\allowbreak{}regress/\allowbreak{}expected/\allowbreak{}entry\_\allowbreak{}route/} (10 files) and \texttt{tests/\allowbreak{}regress/\allowbreak{}expected-peregrine/\allowbreak{}entry\_\allowbreak{}route/} (7 files).
\item \leanfile{tests/\allowbreak{}regress/\allowbreak{}pure\_\allowbreak{}backend.lean} and \leanfile{tests/\allowbreak{}regress/\allowbreak{}pure\_\allowbreak{}oracle.lean}: docstrings, and the messages pinned for \texttt{\#erase} of \texttt{pidHR}; the expected \texttt{pure\_\allowbreak{}backend/\allowbreak{}useM.ast}, \texttt{pure\_\allowbreak{}backend/\allowbreak{}pidHR.ast} and \texttt{pure\_\allowbreak{}oracle/\allowbreak{}pidHR.ast}.
\item \texttt{tests/\allowbreak{}harness/\allowbreak{}expected/\allowbreak{}pure-harness.txt}.
\item \texttt{doc/\allowbreak{}SHIPPING-CHANGES.md}: this entry; the section "Regression tests"; R-3 and R-14 (programs in the fragment); R-15 and R-17 (where \texttt{eraseElab} is).
\end{itemize}
\item \textbf{Why necessary:} the verification's final theorem is about the result of \texttt{eraseEntry} (its hypothesis \texttt{eraseEntry view cfg e = pure (p, inl)}). It is a theorem about \texttt{\#erase} only if \texttt{\#erase} runs \texttt{eraseEntry}. On an input in scope, \texttt{collectDeps} does not answer \texttt{outOfFragment}, so \texttt{eraseEntry} returns only a result of \texttt{erasePure}. A program outside the fragment keeps the \texttt{Meta} path unchanged. An error of the pure path is an error of \texttt{\#erase}, not a fallback to \texttt{Meta}: a fallback would give an input in scope a result that is not verified. The elaborator moves because \texttt{Entry.lean} imports \texttt{Pure.lean}, which imports \texttt{Erasure.lean}: the elaborator cannot call \texttt{eraseEntry} from \texttt{Erasure.lean}.
\item \textbf{Behaviour before:} \texttt{\#erase} runs \texttt{erase e cfg}, the \texttt{Meta} path, on every program. Reproduction: \leanfile{tests/\allowbreak{}regress/\allowbreak{}entry\_\allowbreak{}route.lean} at the parent commit fails with 25 errors (\texttt{route} and \texttt{eraseEntry} unknown, \texttt{\#route} not elaborated). Its \texttt{\#erase} outputs:

\begin{itemize}
\item \texttt{useM.ast}: \texttt{useM := one}, since \texttt{@[macro\_\allowbreak{}inline] mfirst} is inlined;
\item \texttt{pidHR.ast}: \texttt{pid R hR}, which declares \texttt{pid} and the axioms \texttt{R} and \texttt{hR}. The command logs \texttt{No value found for name ER.R, emitting axiom.};
\item \texttt{both.ast}: a program with two declarations named \texttt{ER.two\_\allowbreak{}u39}.
\end{itemize}
\item \textbf{Behaviour after:} \texttt{\#erase} runs \texttt{eraseEntry} over the elaboration environment.

\begin{itemize}
\item \textbf{Outside the fragment:} a program on which \texttt{collectDeps} answers \texttt{outOfFragment} takes the unchanged \texttt{Meta} path: its closure has an inductive type, a constructor, a recursor, a quotient, a literal, a projection, a free variable, a metavariable, a universe metavariable or an unknown constant. Its output and messages are those of the parent commit.
\item \textbf{In the fragment:} a program takes the pure path, and \texttt{\#erase} writes the output of \texttt{erasePure}. The \texttt{Meta} path differs in the following ways:

\begin{itemize}
\item \texttt{@[macro\_\allowbreak{}inline]} and matcher inlining do not happen;
\item a proof whose type is a proposition only through an \texttt{@[irreducible]} alias is erased;
\item a program on which \texttt{Meta} type inference fails through such an alias gets an output;
\item the pure backend logs nothing, so the messages \texttt{No value found for name \ldots{}, emitting axiom.} and \texttt{Name \ldots{} is marked as inline.} do not appear. The messages of the \texttt{.mli} step (\texttt{val main: \ldots{}}, \texttt{failed to translate \ldots{}}) are unchanged.
\end{itemize}
\item \textbf{Errors:} a name collision in the closure is an error of \texttt{\#erase}, for example \texttt{erasure failed: the constants ER.two\_\allowbreak{}u39 and ER.two' have the same kername}. So is an error of \texttt{erasePure}: \texttt{outOfFragment \textquotedbl{}unknown constant \ldots{}\textquotedbl{}}, \texttt{fuel}, or \texttt{failed}.
\item \textbf{Test outputs:} in \leanfile{tests/\allowbreak{}regress/\allowbreak{}entry\_\allowbreak{}route.lean}, \texttt{\#route} shows these routes:

\begin{itemize}
\item \texttt{useM}, \texttt{pidHR} and \texttt{@pid.\{1\}} take \texttt{viaPure ok}, and \texttt{eraseEntry} returns the pure result. For \texttt{useM} and \texttt{pidHR} this differs from \texttt{Meta}'s result; for \texttt{@pid.\{1\}} it is equal;
\item \texttt{useMN := mfirstN Nat.zero Nat.zero} takes \texttt{viaMeta (inductive type)}, and \texttt{@pid} takes \texttt{viaMeta (level metavariable)}. \texttt{eraseEntry} returns \texttt{Meta}'s result;
\item \texttt{both} takes \texttt{viaPure error}.
\end{itemize}

\texttt{useM.ast} keeps \texttt{mfirst one one}, and peregrine evaluates it to \texttt{\ensuremath{\lambda}s. \ensuremath{\lambda}z. s z}. \texttt{pidHR.ast} is \texttt{(Untyped () (Some tBox))}, and \texttt{\#erase ER.both} fails.
\item \textbf{Harness} (\texttt{scripts/\allowbreak{}pure-harness.sh}), 294 \texttt{\#erase} in all:

\begin{itemize}
\item 119 are out of the fragment;
\item 2 are \texttt{collectDeps} errors (\texttt{N\_\allowbreak{}collision}, both configurations);
\item 173 are in the fragment, and the file written by \texttt{\#erase} is the pure output for all 173:

\begin{itemize}
\item 161 are also byte-identical to \texttt{Meta}'s output;
\item 8 differ from it (\texttt{attrEx}, \texttt{irrAxiomArg}, \texttt{irrThmArg}, \texttt{pidHR}, both configurations);
\item 4 have no \texttt{Meta} output (\texttt{useFI}, \texttt{useLamHR}, both configurations).
\end{itemize}
\item The pure oracle makes 5751 calls and never fails. As in S-20, 10 of its answers differ from \texttt{Meta}'s.
\end{itemize}
\item \textbf{peregrine} on the in-fragment outputs (C11): on each of the 173 outputs of the harness, \texttt{peregrine validate} and \texttt{peregrine eval -{}-anf=false} were run on the corpus file, 346 checks in all. 336 pass. The 10 excluded checks fail as expected:

\begin{itemize}
\item \texttt{N\_\allowbreak{}quoteNs} (validate, eval): the namespace component \texttt{q\textquotedbl{}q} is printed without escaping (R-2), and peregrine's parser rejects the file;
\item \texttt{sixOnAxioms} (eval): \texttt{Axioms found};
\item \texttt{ugOne} (validate, eval): \texttt{Fixpoint body is not a lambda} (R-11).
\end{itemize}
\item Axioms: \texttt{Route} depends on none. \texttt{route}, \texttt{eraseEntry} and \texttt{EraseError.describe} depend on \texttt{propext}, \texttt{Classical.choice} and \texttt{Quot.sound}.
\end{itemize}
\item \textbf{Effect on emitted .ast (corpus):} the corpus has 712 files. \texttt{scripts/\allowbreak{}corpus-diff.sh} against the corpus of S-20 reports:

\begin{itemize}
\item 692 files identical. These include every benchmark file and every output of a program outside the fragment.
\item 8 files differing. Each is now the output of the pure path, and each equals the \texttt{erasePure} output of S-20's harness:

\begin{itemize}
\item \texttt{examples/\allowbreak{}Examples/\allowbreak{}attrEx.\{default,peano\}.ast}: \texttt{attrEx := mfirst inlTwo three}, where it was \texttt{inlTwo}, with the declarations of \texttt{three}, \texttt{add} and \texttt{mfirst} in addition. The \texttt{.inlinings} are identical;
\item \texttt{examples/\allowbreak{}Examples/\allowbreak{}irrAxiomArg.\{default,peano\}.ast}: \texttt{guardR \ensuremath{\Box} six}, where it was \texttt{guardR hR six}; the axiom \texttt{Ex.hR} is no longer declared;
\item \texttt{examples/\allowbreak{}Examples/\allowbreak{}irrThmArg.\{default,peano\}.ast}: \texttt{guardR2 \ensuremath{\Box} six}; the declaration \texttt{Ex.hR2 := \ensuremath{\Box}} is no longer emitted;
\item \texttt{examples/\allowbreak{}IrrAlias/\allowbreak{}pidHR.\{default,peano\}.ast}: \texttt{(Untyped () (Some tBox))}.
\end{itemize}
\item 4 files only in S-20: \texttt{examples/\allowbreak{}Names/\allowbreak{}N\_\allowbreak{}collision.\{default,peano\}.ast} and their \texttt{.inlinings}. \texttt{\#erase} fails on these programs now.
\item 12 files only in S-21:

\begin{itemize}
\item \texttt{examples/\allowbreak{}IrrAlias/\allowbreak{}useFI.\{default,peano\}.ast} and \texttt{useLamHR.\{default,peano\}.ast}, with their \texttt{.inlinings}. \texttt{\#erase} failed on these programs before;
\item \texttt{examples/\allowbreak{}Scope/\allowbreak{}pidU1.ast} and \texttt{pcompU111.ast}, with their \texttt{.inlinings}. These are new corpus programs.
\end{itemize}
\end{itemize}

The 6 outputs of \texttt{UnsafeRec.lean} are byte-identical. The corpus logs differ in three ways, and nowhere else:

\begin{itemize}
\item the messages above: the pure backend's missing log lines, the new \texttt{.mli} lines of \texttt{useFI} and \texttt{useLamHR} and of the two new \texttt{Scope} programs, and the missing \texttt{.mli} lines of \texttt{N\_\allowbreak{}collision};
\item job counts, from the two new modules;
\item backtrace frames of the 4 \texttt{PANIC} messages of \texttt{Defects.lean}. The \texttt{PANIC} counts are those of S-20.
\end{itemize}
\item \textbf{Regression test:} \leanfile{tests/\allowbreak{}regress/\allowbreak{}entry\_\allowbreak{}route.lean}. It defines a command \texttt{\#route t [config c]} and uses it to pin, for each program:

\begin{itemize}
\item the route (\texttt{viaPure ok}, \texttt{viaPure error (\ldots{})}, \texttt{viaMeta (\textless{}what\textgreater{})});
\item whether \texttt{eraseEntry} returns the result of the pure path or of the \texttt{Meta} path.
\end{itemize}

It pins the output of \texttt{\#erase} on:

\begin{itemize}
\item three programs in the fragment: \texttt{useM}, \texttt{pidHR}, \texttt{@pid.\{1\}}. peregrine validates and evaluates them;
\item two programs outside it: \texttt{useMN} (validated) and \texttt{@pid}.
\end{itemize}

It pins the error of \texttt{\#erase} on the colliding \texttt{both}. It fails before (25 errors; \texttt{useM.ast} and \texttt{pidHR.ast} differ, and \texttt{both.ast} is written) and passes after. \leanfile{tests/\allowbreak{}regress/\allowbreak{}pure\_\allowbreak{}backend.lean} and \texttt{pure\_\allowbreak{}oracle.lean} pin the new \texttt{\#erase} outputs of \texttt{useM} and \texttt{pidHR}, which now equal the pure path's. \texttt{scripts/\allowbreak{}pure-harness.sh} pins the summary above.
\end{itemize}

\section{\texorpdfstring{S-22: Comments, docstrings and the register cite only what the repository contains}{S-22: Comments, docstrings and the register cite only what the repository contains}}\label{reg:S-22}

\begin{itemize}
\item \textbf{Commit:} the commit whose subject starts with \texttt{shipping(S-22):} (\texttt{git log -{}-grep='\textasciicircum{}shipping(S-22):'}).
\item \textbf{Files and functions:} comments and docstrings only; no code changes.

\begin{itemize}
\item \leanfile{LeanToLambdaBox/\allowbreak{}Erasure.lean}: the docstring of \texttt{ErasureState} names \texttt{EraseT}, the monad that carries it, in place of \texttt{EraseM}.
\item \leanfile{LeanToLambdaBox/\allowbreak{}Erasure/\allowbreak{}Collect.lean}: the comment on the \texttt{DecidableEq} instances of \texttt{ModPath} and \texttt{Kername}, the docstrings of \texttt{EraseError}, \texttt{EnvView}, \texttt{EnvView.ofEnvironment} and \texttt{findConst}, and the heading of the \texttt{collectDeps} section.
\item \leanfile{LeanToLambdaBox/\allowbreak{}Erasure/\allowbreak{}Entry.lean}: the docstrings of \texttt{Route}, \texttt{route} and \texttt{eraseEntry}.
\item \leanfile{LeanToLambdaBox/\allowbreak{}Erasure/\allowbreak{}Pure.lean}: the headings of the oracle, backend and recursion-test sections, and the docstrings of \texttt{oracleFuel}, \texttt{PureCtx}, \texttt{PureState}, \texttt{PureM}, \texttt{PureM.findConst?}, \texttt{PureM.freshFVarId}, \texttt{PureM.instantiate1}, \texttt{PureM.isErasable}, \texttt{PureM.casesInfo?}, \texttt{PureM.ctorArity?}, \texttt{PureM.declInfo?}, \texttt{PureM.unsafeRecBase?}, \texttt{PureM.prepare}, the \texttt{Backend PureM} instance, \texttt{erasePure} and \texttt{nameOccurs}.
\item \leanfile{tests/\allowbreak{}corpus/\allowbreak{}Examples.lean}, \leanfile{tests/\allowbreak{}corpus/\allowbreak{}Names.lean}, \leanfile{tests/\allowbreak{}corpus/\allowbreak{}NearMiss.lean}: the module docstrings.
\item \leanfile{tests/\allowbreak{}regress/\allowbreak{}sparse\_\allowbreak{}discr.lean}: the module docstring cites S-14.
\item \texttt{doc/\allowbreak{}SHIPPING-CHANGES.md}: this entry; the section "The fixed corpus"; the title and files of S-15; the fields "Why necessary" of S-16 to S-21 and "Regression test" of S-19; the list of oracle differences in S-20; the references of S-13 and S-14 to the defect that S-14 fixes, whose entry R-29 under "Reported, not fixed" is removed.
\end{itemize}
\item \textbf{Why necessary:} these texts cited labels and documents that the repository does not contain: the labels S-A to S-F of the changes S-16 to S-21 before they had register ids, questions of a design document (\texttt{DESIGN.md Q2}, \texttt{Q6}, \texttt{Q7}, \texttt{Q8}, \texttt{Q13}), sections of a plan (\texttt{PLAN \S{}3}), a "checkpoint 1", its "decision 13" and a "scope note"; and the monad \texttt{EraseM}, which S-16 removed. A reader cannot resolve any of them. The register listed R-29 under "Reported, not fixed" with the note "fixed in S-14", although S-14 fixes it and describes it in full. The texts now cite the register's entries or state the reason themselves.
\item \textbf{Behaviour before:} \texttt{\#erase} output as in S-21. Reproduction: at the parent commit, \texttt{grep -cE '\textbackslash{}bS-[A-G]\textbackslash{}b\textbar{}DESIGN\textbar{}PLAN\textbar{}[Cc]heckpoint\textbar{}EraseM\textbackslash{}b\textbar{}R-29'} counts 2 lines in \leanfile{LeanToLambdaBox/\allowbreak{}Erasure.lean}, 6 in \texttt{Erasure/\allowbreak{}Collect.lean}, 3 in \texttt{Erasure/\allowbreak{}Entry.lean}, 20 in \texttt{Erasure/\allowbreak{}Pure.lean}, one in each of the three corpus files and one in \leanfile{tests/\allowbreak{}regress/\allowbreak{}sparse\_\allowbreak{}discr.lean}.
\item \textbf{Behaviour after:} the same outputs and log messages. The same \texttt{grep} counts one line: the commented-out \texttt{mkCase} in \leanfile{LeanToLambdaBox/\allowbreak{}Erasure.lean}, byte-identical to \texttt{main}, which is not a declaration.
\item \textbf{Effect on emitted .ast (corpus):} byte-identical for all 712 files; \texttt{scripts/\allowbreak{}corpus-diff.sh} against the corpus of S-21 reports 712 identical. The corpus logs are identical except for the addresses in the backtraces of the 4 \texttt{PANIC} messages of \texttt{Defects.lean}.
\item \textbf{Regression test:} the change does not alter behaviour, so no test fails before it, and no new test is added. The 20 tests of \texttt{tests/\allowbreak{}regress/} and \texttt{scripts/\allowbreak{}pure-harness.sh} pass unchanged.
\end{itemize}

\section{\texorpdfstring{S-23: CI builds and tests the \texttt{shipping} branch when it is pushed}{S-23: CI builds and tests the shipping branch when it is pushed}}\label{reg:S-23}

\begin{itemize}
\item \textbf{Commit:} the commit whose subject starts with \texttt{shipping(S-23):} (\texttt{git log -{}-grep='\textasciicircum{}shipping(S-23):'}).
\item \textbf{Files and functions:}

\begin{itemize}
\item \texttt{.github/\allowbreak{}workflows/\allowbreak{}build.yml}: the \texttt{push} trigger lists the branch \texttt{shipping} next to \texttt{main}.
\item \texttt{doc/\allowbreak{}SHIPPING-CHANGES.md}: this entry; the section "Regression tests" says when CI runs.
\end{itemize}
\item \textbf{Why necessary:} the workflow ran on pushes to \texttt{main} and on pull requests only, so a push of \texttt{shipping}, the branch that holds every change of this register, built nothing and ran neither \texttt{scripts/\allowbreak{}regress.sh} nor \texttt{scripts/\allowbreak{}pure-harness.sh}: a change that breaks the build, a regression test or the harness was caught only after \texttt{shipping} was merged elsewhere.
\item \textbf{Behaviour before:} the \texttt{on.push.branches} list of \texttt{build.yml} is \texttt{['main']}; GitHub starts no run of the workflow for a push of \texttt{shipping}.
\item \textbf{Behaviour after:} the list is \texttt{['main', 'shipping']}; a push of \texttt{shipping} runs the build, the regression tests and the harness, as a push of \texttt{main} does. The emitted programs do not change.
\item \textbf{Effect on emitted .ast (corpus):} byte-identical for all 712 files; \texttt{scripts/\allowbreak{}corpus-diff.sh} against the corpus of S-22 reports 712 identical. The corpus logs are identical except for the addresses in the backtraces of the 4 \texttt{PANIC} messages of \texttt{Defects.lean} and a job count of \texttt{lake} in the log of the pruned benchmarks (\texttt{[18/\allowbreak{}55]} becomes \texttt{[17/\allowbreak{}24]}: the build of S-22's corpus had recompiled the modules that S-22 edits).
\item \textbf{Regression test:} the change concerns CI only, so no test under \texttt{tests/\allowbreak{}regress/} can fail before it; the workflow's own run on a push of \texttt{shipping} checks it. The 20 regression tests and \texttt{scripts/\allowbreak{}pure-harness.sh} pass unchanged.
\end{itemize}

\section{\texorpdfstring{S-24: The docstring of \texttt{binderNameOf} gives the reason that holds}{S-24: The docstring of binderNameOf gives the reason that holds}}\label{reg:S-24}

\begin{itemize}
\item \textbf{Commit:} the commit whose subject starts with \texttt{shipping(S-24):} (\texttt{git log -{}-grep='\textasciicircum{}shipping(S-24):'}).
\item \textbf{Files and functions:}

\begin{itemize}
\item \leanfile{LeanToLambdaBox/\allowbreak{}Erasure.lean}: the docstring of \texttt{binderNameOf}. No code changes.
\item \texttt{doc/\allowbreak{}SHIPPING-CHANGES.md}: this entry.
\end{itemize}
\item \textbf{Why necessary:} the docstring said that a binder name is kept only when it is ASCII graphic "since the \ensuremath{\lambda}\ensuremath{\Box} parser rejects other names", the reason that \texttt{main}'s \texttt{fvar\_\allowbreak{}to\_\allowbreak{}name} gives in a comment. It is false: the printer writes a binder name as a string literal, \texttt{(nNamed \textquotedbl{}\textless{}name\textgreater{}\textquotedbl{})}, and peregrine reads non-ASCII characters in string literals (peregrine-tool \texttt{doc/\allowbreak{}format.md}, "Generic S-expression conventions"). The files \texttt{(Untyped () (Some (tLambda (nNamed \textquotedbl{}\ensuremath{\alpha}\ensuremath{_1}\textquotedbl{}) (tRel 0))))} and the same with \texttt{\textquotedbl{}a b\textquotedbl{}} pass \texttt{peregrine validate} and \texttt{peregrine eval}, and the first also \texttt{peregrine ocaml}, \texttt{c} and \texttt{wasm}. What the test does ensure is that no binder name contains \texttt{\textquotedbl{}} or \texttt{\textbackslash{}}, which the printer does not escape (R-2): \texttt{Name.toString} prints a name with either character with \texttt{\ensuremath{\langle\!\langle}\ensuremath{\rangle\!\rangle}} (\texttt{\ensuremath{\langle\!\langle}a\textquotedbl{}b\ensuremath{\rangle\!\rangle}}, \texttt{a.\ensuremath{\langle\!\langle}b\textquotedbl{}c\ensuremath{\rangle\!\rangle}}), which is not ASCII.
\item \textbf{Behaviour before:} \texttt{\#erase} output as in S-23; the docstring gives the parser as the reason.
\item \textbf{Behaviour after:} the same outputs; the docstring states the test and what it ensures.
\item \textbf{Effect on emitted .ast (corpus):} byte-identical for all 712 files; \texttt{scripts/\allowbreak{}corpus-diff.sh} against the corpus of S-23 reports 712 identical. The corpus logs are identical except for the addresses in the backtraces of the 4 \texttt{PANIC} messages of \texttt{Defects.lean} and a job count of \texttt{lake} in the log of the pruned benchmarks (\texttt{[17/\allowbreak{}24]} becomes \texttt{[18/\allowbreak{}55]}: the build recompiles the modules that import \leanfile{LeanToLambdaBox/\allowbreak{}Erasure.lean}). The docstring keeps its two lines, so the source positions in the \texttt{PANIC} messages do not move.
\item \textbf{Regression test:} the change does not alter behaviour, so no test fails before it, and no new test is added. \leanfile{tests/\allowbreak{}regress/\allowbreak{}traversal\_\allowbreak{}locals.lean}, which checks \texttt{binderNameOf} on \texttt{x}, \texttt{\ensuremath{\alpha}\ensuremath{_1}}, \texttt{a.b} and \texttt{\ensuremath{\langle\!\langle}a b\ensuremath{\rangle\!\rangle}}, passes unchanged, as do the other tests and \texttt{scripts/\allowbreak{}pure-harness.sh}.
\end{itemize}

\section{\texorpdfstring{Reported, not fixed}{Reported, not fixed}}\label{reg:reported-not-fixed}

Each entry below was reproduced on this branch. Unless stated otherwise, the reproduction is an example in \leanfile{tests/\allowbreak{}corpus/\allowbreak{}Defects.lean} whose output lies in \texttt{examples/\allowbreak{}Defects/} of a corpus, and \texttt{peregrine} is \texttt{../\allowbreak{}peregrine-tool/\allowbreak{}\_\allowbreak{}build/\allowbreak{}default/\allowbreak{}bin/\allowbreak{}main.exe}.

\subsection{\texorpdfstring{R-1: Primitive integers are printed as quoted strings}{R-1: Primitive integers are printed as quoted strings}}\label{reg:R-1}

\begin{itemize}
\item \textbf{What:} a machine-mode Nat literal is printed as \texttt{(tPrim (primInt \textquotedbl{}7\textquotedbl{}))}, a string atom. \texttt{peregrine-tool/\allowbreak{}doc/\allowbreak{}format.md} specifies a numeric atom, \texttt{(primInt 7)}, and peregrine's parser accepts only that (since peregrine commit \texttt{556fe09}).
\item \textbf{Where:} \leanfile{LeanToLambdaBox/\allowbreak{}Printing.lean}: instance \texttt{Serialize (BitVec n)}.
\item \textbf{Reproduction:} \texttt{primint.ast} (\texttt{def seven : Nat := 7}, default configuration): \texttt{peregrine validate primint.ast} fails with \texttt{could not read integral type, got a non-Num atom \textquotedbl{}7\textquotedbl{}}. All 40 benchmark \texttt{.ast} files contain such atoms.
\item \textbf{Impact:} no output of the default configuration (\texttt{nat := .machine}) that contains a Nat literal can be read by the current peregrine. Outputs with \texttt{nat := .peano} contain no primitive integer.
\item \textbf{Why not fixed:} not required by the verification goal unless it later becomes required.
\end{itemize}

\subsection{\texorpdfstring{R-2: String atoms are not escaped}{R-2: String atoms are not escaped}}\label{reg:R-2}

\begin{itemize}
\item \textbf{What:} string atoms are written between double quotes with no escaping. Kername identifiers are sanitized, but module-path components and inductive and constructor names are printed from the Lean name as they are, so a \texttt{\textquotedbl{}} or \texttt{\textbackslash{}} in them breaks the S-expression.
\item \textbf{Where:} \leanfile{LeanToLambdaBox/\allowbreak{}Printing.lean}: \texttt{quote\_\allowbreak{}atom}; \leanfile{LeanToLambdaBox/\allowbreak{}Basic.lean}: \texttt{toModPath}; \leanfile{LeanToLambdaBox/\allowbreak{}Erasure.lean}: \texttt{register\_\allowbreak{}inductive} (\texttt{toString ind\_\allowbreak{}name}, \texttt{toString ctor\_\allowbreak{}name}).
\item \textbf{Reproduction:} \texttt{quote\_\allowbreak{}modpath.ast} (\texttt{\ensuremath{\langle\!\langle}ns\textquotedbl{}q\ensuremath{\rangle\!\rangle}.f 1}): \texttt{peregrine validate} fails with \texttt{could not read 'prod', expected list of length 2, got list of a different length}; \texttt{quote\_\allowbreak{}inductive.ast} (\texttt{inductive \ensuremath{\langle\!\langle}Q\textquotedbl{}T\ensuremath{\rangle\!\rangle} where \textbar{} \ensuremath{\langle\!\langle}m\textquotedbl{}k\ensuremath{\rangle\!\rangle}}): \texttt{Invalid character}.
\item \textbf{Impact:} programs with such names produce unreadable files.
\item \textbf{Why not fixed:} not required by the verification goal unless it later becomes required.
\end{itemize}

\subsection{\texorpdfstring{R-3: Kernames are not injective}{R-3: Kernames are not injective}}\label{reg:R-3}

\begin{itemize}
\item \textbf{What:} \texttt{cleanIdent} replaces every character outside \texttt{[A-Za-z0-9\_\allowbreak{}]} by \texttt{\_\allowbreak{}u\textless{}code\textgreater{}}, so \texttt{a b} and \texttt{a\_\allowbreak{}u32b} get the same identifier; numeric name components and digit strings coincide; a mutual block's kername is the concatenation of its members' names (\texttt{[A, B]} and \texttt{AB} coincide); and the module (\texttt{DirPath}) is always empty.
\item \textbf{Where:} \leanfile{LeanToLambdaBox/\allowbreak{}Basic.lean}: \texttt{cleanIdent}, \texttt{toKername}, \texttt{toModPath}, \texttt{rootKername}; \leanfile{LeanToLambdaBox/\allowbreak{}Erasure.lean}: \texttt{register\_\allowbreak{}inductive} (\texttt{mutualBlockName}).
\item \textbf{Reproduction:} \texttt{collide.ast} (\texttt{colA.\ensuremath{\langle\!\langle}a b\ensuremath{\rangle\!\rangle} + colA.a\_\allowbreak{}u32b}): \texttt{peregrine validate} fails with \texttt{Duplicate definition .Defects.colA.a\_\allowbreak{}u32b}. A program in the fragment of the pure path is not affected: \texttt{collectDeps} finds the collision and \texttt{\#erase} fails (S-21), e.g. on \texttt{Nm.both} in \leanfile{tests/\allowbreak{}corpus/\allowbreak{}Names.lean} with \texttt{erasure failed: the constants Nm.two\_\allowbreak{}u39 and Nm.two' have the same kername}.
\item \textbf{Impact:} distinct Lean constants can be emitted under one \ensuremath{\lambda}\ensuremath{\Box} name, for programs outside the fragment.
\item \textbf{Why not fixed:} not required by the verification goal unless it later becomes required.
\end{itemize}

\subsection{\texorpdfstring{R-4: Unsupported literals are erased to \ensuremath{\Box} after a panic}{R-4: Unsupported literals are erased to   after a panic}}\label{reg:R-4}

\begin{itemize}
\item \textbf{What:} string literals, and Nat literals above \texttt{2\textasciicircum{}62 - 1} in machine mode, reach \texttt{panic!}. A Lean panic prints \texttt{PANIC ...} and continues with the default value, here \texttt{\ensuremath{\Box}}, so \texttt{\#erase} succeeds, writes the file and exits with status 0.
\item \textbf{Where:} \leanfile{LeanToLambdaBox/\allowbreak{}Erasure.lean}: \texttt{visitLiteral}.
\item \textbf{Reproduction:} \texttt{strlit.ast} (\texttt{\#erase \textquotedbl{}abc\textquotedbl{}}) is \texttt{(Untyped () (Some tBox))}; in \texttt{biglit.ast} (\texttt{5000000000000000000 : Nat}) the literal is \texttt{tBox}. The log contains \texttt{PANIC at Erasure.visitLiteral}. Both files pass \texttt{peregrine validate}.
\item \textbf{Impact:} silent miscompilation.
\item \textbf{Why not fixed:} not required by the verification goal unless it later becomes required.
\end{itemize}

\subsection{\texorpdfstring{R-5: An alternative whose \ensuremath{\Pi}-type is not syntactic gets a branch without binders}{R-5: An alternative whose  -type is not syntactic gets a branch without binders}}\label{reg:R-5}

\begin{itemize}
\item \textbf{What:} a \texttt{casesOn} alternative that is not a \ensuremath{\lambda} is \ensuremath{\eta}-expanded through its inferred type with \texttt{forallMonocular}, which expects a syntactic \texttt{\ensuremath{\forall}}. If the type is a \ensuremath{\Pi} only after unfolding, it hits \texttt{unreachable!}; the panic's default yields a branch with no binders, which no longer matches the constructor's arity.
\item \textbf{Where:} \leanfile{LeanToLambdaBox/\allowbreak{}Erasure.lean}: \texttt{forallMonocular}, \texttt{lambdaMonocularOrIntro}, \texttt{lambdaOrIntroToArity} (called by \texttt{visitAlt}); \texttt{withAppEtaToMinArity} has the same limitation.
\item \textbf{Reproduction:} \texttt{nonsyn\_\allowbreak{}pi.ast} (\texttt{viaCases 3}, where the successor alternative is \texttt{g : MyFun} with \texttt{def MyFun := Nat \ensuremath{\to} Nat}; Lean value 2). The log contains \texttt{PANIC at Erasure.forallMonocular}; \texttt{peregrine validate} accepts the file; \texttt{peregrine eval} prints \texttt{constr con\_\allowbreak{}15} / \texttt{constr con\_\allowbreak{}106} instead of \texttt{Nat.succ (Nat.succ Nat.zero)}.
\item \textbf{Impact:} silent miscompilation that \texttt{peregrine validate} does not detect.
\item \textbf{Why not fixed:} not required by the verification goal unless it later becomes required.
\end{itemize}

\subsection{\texorpdfstring{R-6: Peano literals exhaust the stack}{R-6: Peano literals exhaust the stack}}\label{reg:R-6}

\begin{itemize}
\item \textbf{What:} with \texttt{nat := .peano}, a literal \texttt{n} is translated by \texttt{n} nested calls (\texttt{visitLiteral} \ensuremath{\to} \texttt{visitConstructor} \ensuremath{\to} \texttt{visitAppArgs} \ensuremath{\to} \texttt{visitExpr}).
\item \textbf{Where:} \leanfile{LeanToLambdaBox/\allowbreak{}Erasure.lean}: \texttt{visitLiteral} (Peano branch).
\item \textbf{Reproduction:} a file containing \texttt{import LeanToLambdaBox} and \texttt{\#erase (1000000 : Nat) config \{nat := .peano, extern := .preferLogical\} to \textquotedbl{}p.ast\textquotedbl{}} makes Lean abort (exit status 134) with \texttt{deep recursion was detected at 'interpreter'}; \texttt{200000} succeeds and writes 23 601 044 bytes (\texttt{5000}: 591 044 bytes). It is not in the corpus because it aborts the elaboration.
\item \textbf{Impact:} Peano mode fails on large literals, and the output grows linearly with the value.
\item \textbf{Why not fixed:} not required by the verification goal unless it later becomes required.
\end{itemize}

\subsection{\texorpdfstring{R-7: Inductive types are always declared non-propositional}{R-7: Inductive types are always declared non-propositional}}\label{reg:R-7}

\begin{itemize}
\item \textbf{What:} every inductive body is printed with \texttt{ind\_\allowbreak{}propositional = false} and \texttt{IntoAny}, also for inductives in Prop (\texttt{And}, \texttt{Eq}, \texttt{Acc}, \ldots{}). A match on a proof of a large-eliminating Prop inductive becomes a \texttt{tCase} on \texttt{\ensuremath{\Box}}. MetaRocq evaluates such a case only for inductives flagged propositional (\texttt{eval\_\allowbreak{}iota\_\allowbreak{}sing}, \texttt{erasure/\allowbreak{}theories/\allowbreak{}EWcbvEval.v}), and removes such cases only for them (\texttt{EOptimizePropDiscr.v}).
\item \textbf{Where:} \leanfile{LeanToLambdaBox/\allowbreak{}Basic.lean}: \texttt{OneInductiveBody} (defaults \texttt{propositional := false}, \texttt{kelim := .IntoAny}); \leanfile{LeanToLambdaBox/\allowbreak{}Erasure.lean}: \texttt{register\_\allowbreak{}inductive}.
\item \textbf{Reproduction:} \texttt{and\_\allowbreak{}elim.ast} (\texttt{andElim True True \ensuremath{\langle}trivial, trivial\ensuremath{\rangle}}, Lean value 7): \texttt{peregrine eval} fails with \texttt{Case: \textless{}15\textgreater{} branch not found}. With the flag of \texttt{And} set to \texttt{true} by hand in a copy, it evaluates to 7.
\item \textbf{Impact:} programs that eliminate proofs of \texttt{And}, \texttt{Eq}, \texttt{Acc}, \ldots{} into data do not evaluate.
\item \textbf{Why not fixed:} not required by the verification goal unless it later becomes required.
\end{itemize}

\subsection{\texorpdfstring{R-8: Informative index fields of large-eliminating Prop inductives are lost}{R-8: Informative index fields of large-eliminating Prop inductives are lost}}\label{reg:R-8}

\begin{itemize}
\item \textbf{What:} Lean lets a Prop inductive with one constructor eliminate into any sort when each non-Prop field of the constructor also appears as an index of its result type (for instance \texttt{Foo.mk (n : Nat) : Foo 0 n}, or \texttt{Acc.intro}, whose \texttt{x} is an index). The branch of such a match receives the field. The eraser turns the proof into \texttt{\ensuremath{\Box}} and binds the field in the branch; the value exists only in the index argument, which is dropped. Setting the propositional flag (R-7) does not help, since MetaRocq's rule for a case on \texttt{\ensuremath{\Box}} fills the fields with \texttt{\ensuremath{\Box}}.
\item \textbf{Where:} \leanfile{LeanToLambdaBox/\allowbreak{}Erasure.lean}: \texttt{visitCases} (generic path).
\item \textbf{Reproduction:} \texttt{index\_\allowbreak{}field.ast} (\texttt{getN 1 (Foo.mk 1)}, \texttt{\#reduce} gives 1): \texttt{peregrine eval} fails with \texttt{Case: \textless{}15\textgreater{} branch not found}; with the flag of \texttt{Foo} set to \texttt{true} by hand in a copy, it prints \texttt{constr con\_\allowbreak{}15} / \texttt{constr con\_\allowbreak{}107} instead of \texttt{Nat.succ Nat.zero}.
\item \textbf{Impact:} such programs do not evaluate, and would compute wrong values once R-7 is fixed alone.
\item \textbf{Why not fixed:} not required by the verification goal unless it later becomes required.
\end{itemize}

\subsection{\texorpdfstring{R-9: Recursors and other constants without a value become axioms}{R-9: Recursors and other constants without a value become axioms}}\label{reg:R-9}

\begin{itemize}
\item \textbf{What:} a constant without a value (every recursor \texttt{T.rec}, the \texttt{Quot} primitives, Lean axioms) is emitted as \texttt{ConstantDecl None}. Casts along equalities (\texttt{h \ensuremath{\blacktriangleright} x}, matches on \texttt{rfl}) go through \texttt{Eq.rec}, well-founded recursion reaches \texttt{False.rec}, and noncomputable recursion reaches \texttt{T.rec}; none of them gets a \ensuremath{\lambda}\ensuremath{\Box} definition.
\item \textbf{Where:} \leanfile{LeanToLambdaBox/\allowbreak{}Erasure.lean}: \texttt{visitMutual} (the branch where \texttt{ci.value?} is \texttt{none}), \texttt{addAxiom}.
\item \textbf{Reproduction:} \texttt{eq\_\allowbreak{}cast.ast} (\texttt{castTo Nat rfl 3}): \texttt{peregrine eval} fails with \texttt{Axioms found ... .Eq.rec}; \texttt{wf\_\allowbreak{}rec.ast} (\texttt{log2 8}, well-founded recursion): \texttt{Axioms found ... .False.rec}.
\item \textbf{Impact:} such programs run only where the axioms are implemented outside \ensuremath{\lambda}\ensuremath{\Box}; the benchmark runtime implements \texttt{Eq.rec} only.
\item \textbf{Why not fixed:} not required by the verification goal unless it later becomes required.
\end{itemize}

\subsection{\texorpdfstring{R-10: Recursive structures get no projection declarations}{R-10: Recursive structures get no projection declarations}}\label{reg:R-10}

\begin{itemize}
\item \textbf{What:} projection declarations are generated only for a non-recursive inductive with one constructor alone in its block, but \texttt{Expr.proj} on any structure is translated to \texttt{tProj}.
\item \textbf{Where:} \leanfile{LeanToLambdaBox/\allowbreak{}Erasure.lean}: \texttt{register\_\allowbreak{}inductive} (\texttt{is\_\allowbreak{}struct}), \texttt{visitProj}.
\item \textbf{Reproduction:} \texttt{rec\_\allowbreak{}struct\_\allowbreak{}proj.ast} (\texttt{structure RS where val : Nat; kids : List RS}, \texttt{rsVal \ensuremath{\langle}3, []\ensuremath{\rangle}}): \texttt{peregrine validate} fails with \texttt{Projection .Defects\_\allowbreak{}u46RS,0:0,0 not found}.
\item \textbf{Impact:} programs that project out of recursive structures are rejected.
\item \textbf{Why not fixed:} not required by the verification goal unless it later becomes required.
\end{itemize}

\subsection{\texorpdfstring{R-11: Fixpoints have recursive-argument index 0 and bodies that are not \ensuremath{\eta}-expanded}{R-11: Fixpoints have recursive-argument index 0 and bodies that are not  -expanded}}\label{reg:R-11}

\begin{itemize}
\item \textbf{What:} every fixpoint definition is printed with \texttt{rarg = 0}, and its body is used as it is (the code has a TODO to \ensuremath{\eta}-expand it). MetaRocq's fixpoint well-formedness requires a \ensuremath{\lambda} body and its guarded-fixpoint evaluation reads \texttt{rarg}. Current outputs work because the compiler's pre-definitions are \ensuremath{\lambda}s and evaluation is call by value.
\item \textbf{Where:} \leanfile{LeanToLambdaBox/\allowbreak{}Basic.lean}: \texttt{FixDef} (\texttt{principalArgIdx := 0}); \leanfile{LeanToLambdaBox/\allowbreak{}Erasure.lean}: \texttt{visitMutual}, \texttt{mkDef}.
\item \textbf{Reproduction:} every recursive definition, e.g. corpus \texttt{examples/\allowbreak{}PortProbe/\allowbreak{}fact.default.ast}: \texttt{(def (nNamed \textquotedbl{}Tiny.fact\textquotedbl{}) (tLambda ...) 0)}. A body that is not a \ensuremath{\lambda}: corpus \texttt{examples/\allowbreak{}UnsafeRec/\allowbreak{}ugOne.default.ast}, from \texttt{unsafe def ug : CN.\{0\} \ensuremath{\to} CN.\{0\} := (fun \_\allowbreak{} n =\textgreater{} n) (fun x =\textgreater{} ug x)}, has \texttt{(def (nNamed \textquotedbl{}URec.ug\textquotedbl{}) (tApp ...) 0)}, and \texttt{peregrine validate} rejects it: \texttt{Error while checking .URec.ug: Fixpoint body is not a lambda}.
\item \textbf{Impact:} none observed on recursive definitions whose value is a \ensuremath{\lambda}, which includes every compiler pre-definition; a recursive \texttt{unsafe def} whose value is not a \ensuremath{\lambda} gives a program that peregrine rejects.
\item \textbf{Why not fixed:} not required by the verification goal unless it later becomes required.
\end{itemize}

\subsection{\texorpdfstring{R-12: Mutual blocks skip the \texttt{@[extern]} and \texttt{@[inline]} handling}{R-12: Mutual blocks skip the @[extern] and @[inline] handling}}\label{reg:R-12}

\begin{itemize}
\item \textbf{What:} the rules that turn \texttt{@[extern]} constants into axioms and record \texttt{@[inline]} / \texttt{@[always\_\allowbreak{}inline]} constants run only for blocks with a single declaration.
\item \textbf{Where:} \leanfile{LeanToLambdaBox/\allowbreak{}Erasure.lean}: \texttt{visitMutual}.
\item \textbf{Reproduction:} \texttt{mutual\_\allowbreak{}inline.ast.inlinings} does not list \texttt{fooI}, which is \texttt{@[inline]} in a mutual block. For \texttt{@[extern]}: with the default configuration, a recursive \texttt{@[extern \textquotedbl{}f\textquotedbl{}] def} becomes an axiom, while the same definition inside a \texttt{mutual} block is erased to its \texttt{tFix} (this case is not in the corpus).
\item \textbf{Impact:} inconsistent treatment of attributes; low.
\item \textbf{Why not fixed:} not required by the verification goal unless it later becomes required.
\end{itemize}

\subsection{\texorpdfstring{R-13: \texttt{@[implemented\_\allowbreak{}by]} is ignored}{R-13: @[implemented\_by] is ignored}}\label{reg:R-13}

\begin{itemize}
\item \textbf{What:} the eraser always uses a constant's logical definition, while Lean's compiler runs its \texttt{@[implemented\_\allowbreak{}by]} implementation. When the two disagree, the erased program and Lean's compiled code compute different values.
\item \textbf{Where:} \leanfile{LeanToLambdaBox/\allowbreak{}Erasure.lean}: \texttt{visitMutual} (uses \texttt{ci.value?}), \texttt{visitConstApp} (no \texttt{implemented\_\allowbreak{}by} check).
\item \textbf{Reproduction:} \texttt{implemented\_\allowbreak{}by.ast} (\texttt{fastId 2}, where \texttt{fastId n := n + 1} is implemented by \texttt{slowId n := n}): \texttt{\#eval fastId 2} prints 2; \texttt{peregrine eval implemented\_\allowbreak{}by.ast} prints 3.
\item \textbf{Impact:} differs from Lean's compiled code for such constants; the logical definition is what a proof about the kernel term refers to.
\item \textbf{Why not fixed:} not required by the verification goal unless it later becomes required.
\end{itemize}

\subsection{\texorpdfstring{R-14: Erasability is decided by the elaborator at default transparency}{R-14: Erasability is decided by the elaborator at default transparency}}\label{reg:R-14}

\begin{itemize}
\item \textbf{What:} \texttt{isErasable} uses \texttt{Meta.inferType}, \texttt{Meta.isProp} and \texttt{Meta.isTypeFormerType}; the last one weak-head normalizes at default transparency, which does not unfold \texttt{@[irreducible]} definitions. A type behind an irreducible alias is therefore kept as a relevant term, and so is a proof whose proposition's sort is behind such an alias. Where \texttt{Meta.inferType} needs to unfold the alias to a \ensuremath{\Pi}-type or a sort, \texttt{\#erase} fails.
\item \textbf{Where:} \leanfile{LeanToLambdaBox/\allowbreak{}Erasure.lean}: \texttt{isErasable}.
\item \textbf{Reproduction:} \texttt{irreducible\_\allowbreak{}alias.ast}: \texttt{mkT : MyType}, with \texttt{@[irreducible] def MyType := Type}, is emitted as the declaration \texttt{mkT := id \ensuremath{\Box} \ensuremath{\Box}} instead of being erased. Programs in the fragment of the pure path are not affected: \texttt{\#erase} erases them with the pure oracle, which unfolds \texttt{@[irreducible]} definitions (S-19, S-21). In the corpus files \leanfile{tests/\allowbreak{}corpus/\allowbreak{}IrrAlias.lean} and \leanfile{tests/\allowbreak{}corpus/\allowbreak{}Examples.lean}, with \texttt{@[irreducible] def IProp : Type := Prop}, \texttt{axiom R : IProp} and \texttt{axiom hR : R}, \texttt{examples/\allowbreak{}IrrAlias/\allowbreak{}pidHR.default.ast} (\texttt{pid hR}) is \texttt{\ensuremath{\Box}}, \texttt{examples/\allowbreak{}Examples/\allowbreak{}irrAxiomArg.default.ast} (\texttt{guardR hR six}) passes \texttt{\ensuremath{\Box}} for \texttt{hR}, and \texttt{Irr.useFI := fI two} (with \texttt{@[irreducible] def EndoC : Type 1 := CNat \ensuremath{\to} CNat}) and \texttt{Irr.useLamHR := guardR (lamHR two) six} are erased; with the \texttt{Meta} oracle, the first two keep \texttt{hR} and the last two fail (\texttt{function expected}, \texttt{type expected}).
\item \textbf{Impact:} under-erasure, for programs outside the fragment. In \texttt{irreducible\_\allowbreak{}alias.ast} the value computed is unaffected; a kept proof axiom stops \texttt{peregrine eval}; and \texttt{\#erase} refuses some programs.
\item \textbf{Why not fixed:} not required by the verification goal unless it later becomes required.
\end{itemize}

\subsection{\texorpdfstring{R-15: \texttt{\#erase} without \texttt{to} logs the program in place of the attributes}{R-15: \#erase without to logs the program in place of the attributes}}\label{reg:R-15}

\begin{itemize}
\item \textbf{What:} without an output path, the branch meant to log the attributes configuration logs the program a second time.
\item \textbf{Where:} \leanfile{LeanToLambdaBox/\allowbreak{}Erasure/\allowbreak{}Command.lean}: \texttt{eraseElab} (the \texttt{.none} case of the \texttt{.inlinings} match).
\item \textbf{Reproduction:} a file containing \texttt{import LeanToLambdaBox}, \texttt{def seven : Nat := 7} and \texttt{\#erase seven}: the messages contain the program twice and no \texttt{(attributes\_\allowbreak{}config ...)}.
\item \textbf{Impact:} cosmetic.
\item \textbf{Why not fixed:} not required by the verification goal unless it later becomes required.
\end{itemize}

\subsection{\texorpdfstring{R-16: A log message lacks string interpolation}{R-16: A log message lacks string interpolation}}\label{reg:R-16}

\begin{itemize}
\item \textbf{What:} the message for \texttt{@[extern]} constructors is a plain string literal, so it prints the placeholders literally.
\item \textbf{Where:} \leanfile{LeanToLambdaBox/\allowbreak{}Erasure.lean}: \texttt{register\_\allowbreak{}inductive}.
\item \textbf{Reproduction:} after \texttt{scripts/\allowbreak{}corpus.sh OUT}, \texttt{OUT/\allowbreak{}\_\allowbreak{}meta/\allowbreak{}logs/\allowbreak{}benchmarks-prune.log} contains \texttt{Constructor \{ctor\_\allowbreak{}name\} of type \{ind\_\allowbreak{}name\} is marked @[extern], emitting axiom.}
\item \textbf{Impact:} cosmetic.
\item \textbf{Why not fixed:} not required by the verification goal unless it later becomes required.
\end{itemize}

\subsection{\texorpdfstring{R-17: The \texttt{.mli} signature falls back to \texttt{unit}}{R-17: The .mli signature falls back to unit}}\label{reg:R-17}

\begin{itemize}
\item \textbf{What:} \texttt{to\_\allowbreak{}ml\_\allowbreak{}type} handles \texttt{Nat}, \texttt{Int}, \texttt{Unit}/\texttt{PUnit}, \texttt{Bool}, \texttt{List}, \texttt{Option}, \texttt{Array}, \texttt{Prod} and arrows; any other type, for instance a user inductive or \texttt{String} (S-5), is reported with a warning and printed as \texttt{unit}, which does not describe the value. The file has no final newline.
\item \textbf{Where:} \leanfile{LeanToLambdaBox/\allowbreak{}Erasure.lean}: \texttt{to\_\allowbreak{}ml\_\allowbreak{}type}, \texttt{gen\_\allowbreak{}mli}; \leanfile{LeanToLambdaBox/\allowbreak{}Erasure/\allowbreak{}Command.lean}: \texttt{eraseElab}.
\item \textbf{Reproduction:} \texttt{mli\_\allowbreak{}fallback.mli} (\texttt{toMy : Nat \ensuremath{\to} MyNat}, a user inductive) is \texttt{val main: Z.t -\textgreater{} unit}, with the warning \texttt{failed to translate Defects.MyNat into ML type, emitting unit instead}.
\item \textbf{Impact:} an OCaml harness written against such a signature is type-unsound.
\item \textbf{Why not fixed:} not required by the verification goal unless it later becomes required.
\end{itemize}

\subsection{\texorpdfstring{R-18: The benchmark Makefile does not re-erase after a frontend change}{R-18: The benchmark Makefile does not re-erase after a frontend change}}\label{reg:R-18}

\begin{itemize}
\item \textbf{What:} the rule that produces \texttt{\%.ast}, \texttt{\%.ast.inlinings} and \texttt{\%.mli} depends on the generated \texttt{.lean} file and on \texttt{../\allowbreak{}FromLeanCommon*}, not on the frontend sources. A build directory keeps the outputs of the frontend that first produced them.
\item \textbf{Where:} \texttt{benchmarks/\allowbreak{}via\_\allowbreak{}malfunction/\allowbreak{}Makefile}: rule \texttt{\$(build)/\allowbreak{}\%.ast \$(build)/\allowbreak{}\%.ast.inlinings \$(build)/\allowbreak{}\%.mli: \$(build)/\allowbreak{}\%.lean ../\allowbreak{}FromLeanCommon.lean ../\allowbreak{}FromLeanCommon/}.
\item \textbf{Reproduction:} in \texttt{benchmarks/\allowbreak{}via\_\allowbreak{}malfunction}, after \texttt{make build/\allowbreak{}\textless{}id\textgreater{}/\allowbreak{}even.ast}, edit \leanfile{LeanToLambdaBox/\allowbreak{}Erasure.lean} and run the same command: make prints \texttt{make: 'build/\allowbreak{}\textless{}id\textgreater{}/\allowbreak{}even.ast' is up to date.} \texttt{scripts/\allowbreak{}corpus.sh} deletes the outputs before erasing for this reason.
\item \textbf{Impact:} stale benchmark outputs after frontend changes.
\item \textbf{Why not fixed:} not required by the verification goal unless it later becomes required.
\end{itemize}

\subsection{\texorpdfstring{R-19: \texttt{lake build} in \texttt{benchmarks/} builds nothing}{R-19: lake build in benchmarks/ builds nothing}}\label{reg:R-19}

\begin{itemize}
\item \textbf{What:} \texttt{benchmarks/\allowbreak{}lakefile.toml} sets \texttt{default\_\allowbreak{}targets}; Lake's key is \texttt{defaultTargets}, so the package has no default target.
\item \textbf{Where:} \texttt{benchmarks/\allowbreak{}lakefile.toml}.
\item \textbf{Reproduction:} in a copy of \texttt{benchmarks/} without \texttt{.lake}, \texttt{lake build} prints \texttt{warning: no targets specified and no default targets configured} and \texttt{Nothing to build.}, and creates no \texttt{.lake}.
\item \textbf{Impact:} low; the benchmark Makefile builds what it needs through \texttt{lake lean}.
\item \textbf{Why not fixed:} not required by the verification goal unless it later becomes required.
\end{itemize}

\subsection{\texorpdfstring{R-20: The benchmark runtime disagrees with Lean on some operations}{R-20: The benchmark runtime disagrees with Lean on some operations}}\label{reg:R-20}

\begin{itemize}
\item \textbf{What:}

\begin{itemize}
\item \texttt{def\_\allowbreak{}\_\allowbreak{}Nat\_\allowbreak{}pred = Z.pred} gives \texttt{Nat.pred 0 = -1}; Lean gives 0.
\item \texttt{def\_\allowbreak{}\_\allowbreak{}Int\_\allowbreak{}sub = Z.add} in \texttt{int.ml} (the inline variant \texttt{int-inline.ml} is correct).
\item \texttt{nat.mli} types the second argument of \texttt{Nat.shiftl}, \texttt{Nat.shiftr} and \texttt{Nat.pow} as OCaml \texttt{int}, while the erased program passes a \texttt{Z.t}.
\item \texttt{def\_\allowbreak{}\_\allowbreak{}Array\_\allowbreak{}get\_\allowbreak{}u33Internal} (\texttt{Array.get!}) out of bounds runs \texttt{assert false}; Lean panics and returns the default element.
\end{itemize}
\item \textbf{Where:} \texttt{benchmarks/\allowbreak{}via\_\allowbreak{}malfunction/\allowbreak{}nat.ml:25} and \texttt{nat-inline.ml:58} (\texttt{def\_\allowbreak{}\_\allowbreak{}Nat\_\allowbreak{}pred}), \texttt{int.ml:9} (\texttt{def\_\allowbreak{}\_\allowbreak{}Int\_\allowbreak{}sub}), \texttt{nat.mli:16-18}, \texttt{JCFArray.ml:168-170}.
\item \textbf{Reproduction:} the cited lines. No benchmark program references \texttt{Nat.pred} or \texttt{Int.sub} (none of the 40 benchmark \texttt{.ast} files names them).
\item \textbf{Impact:} wrong results or aborts for programs that use these operations on these inputs.
\item \textbf{Why not fixed:} not required by the verification goal unless it later becomes required.
\end{itemize}

\subsection{\texorpdfstring{R-21: Axioms without an OCaml implementation fail only at link time}{R-21: Axioms without an OCaml implementation fail only at link time}}\label{reg:R-21}

\begin{itemize}
\item \textbf{What:} the eraser emits an axiom for every constant without a value and, under \texttt{extern := .preferAxiom}, for every \texttt{@[extern]} constant; the benchmark runtime implements a fixed set (\texttt{Nat}, \texttt{Int}, \texttt{Eq.rec}, \texttt{Decidable}, eight \texttt{Array} operations). Any other axiom (for instance \texttt{False.rec}, \texttt{Quot.lift}, \texttt{Array.pop}, \texttt{UInt32} operations) erases without error and fails when the OCaml program is linked.
\item \textbf{Where:} \leanfile{LeanToLambdaBox/\allowbreak{}Erasure.lean}: \texttt{visitMutual}, \texttt{addAxiom}; \texttt{benchmarks/\allowbreak{}via\_\allowbreak{}malfunction/\allowbreak{}axioms.ml}.
\item \textbf{Reproduction:} \texttt{wf\_\allowbreak{}rec.ast} references the axiom \texttt{False.rec}; no file in \texttt{benchmarks/\allowbreak{}via\_\allowbreak{}malfunction/} defines \texttt{def\_\allowbreak{}\_\allowbreak{}False\_\allowbreak{}rec}.
\item \textbf{Impact:} missing implementations are reported by the OCaml linker, not by the eraser.
\item \textbf{Why not fixed:} not required by the verification goal unless it later becomes required.
\end{itemize}

\subsection{\texorpdfstring{R-22: The benchmark \texttt{list\_\allowbreak{}sum\_\allowbreak{}rev} does not reverse}{R-22: The benchmark list\_sum\_rev does not reverse}}\label{reg:R-22}

\begin{itemize}
\item \textbf{What:} \texttt{list\_\allowbreak{}sum\_\allowbreak{}rev n := List.replicate n 1 \textbar{}\textgreater{}.foldl Nat.add 0} is the same program as \texttt{list\_\allowbreak{}sum\_\allowbreak{}foldl}; the reversing variant is \texttt{shared\_\allowbreak{}list\_\allowbreak{}sum\_\allowbreak{}rev}, which is not in \texttt{benchmarks/\allowbreak{}TESTS}.
\item \textbf{Where:} \leanfile{benchmarks/\allowbreak{}FromLeanCommon.lean}: \texttt{list\_\allowbreak{}sum\_\allowbreak{}rev}.
\item \textbf{Reproduction:} corpus \texttt{benchmarks/\allowbreak{}prune/\allowbreak{}list\_\allowbreak{}sum\_\allowbreak{}rev.ast} equals \texttt{list\_\allowbreak{}sum\_\allowbreak{}foldl.ast} byte for byte once the name \texttt{list\_\allowbreak{}sum\_\allowbreak{}rev} is replaced by \texttt{list\_\allowbreak{}sum\_\allowbreak{}foldl}.
\item \textbf{Impact:} the benchmark measures the left fold twice.
\item \textbf{Why not fixed:} not required by the verification goal unless it later becomes required.
\end{itemize}

\subsection{\texorpdfstring{R-23: Documentation drift}{R-23: Documentation drift}}\label{reg:R-23}

\begin{itemize}
\item \textbf{What:} \texttt{README.md} shows \texttt{import Erasure}; the module is \texttt{LeanToLambdaBox}. \texttt{benchmarks/\allowbreak{}via\_\allowbreak{}malfunction/\allowbreak{}README.md} says "These four switches" and lists three.
\item \textbf{Where:} \texttt{README.md:7}; \texttt{benchmarks/\allowbreak{}via\_\allowbreak{}malfunction/\allowbreak{}README.md:12}.
\item \textbf{Reproduction:} a file whose first line is \texttt{import Erasure} fails with \texttt{unknown module prefix 'Erasure'}.
\item \textbf{Impact:} documentation only.
\item \textbf{Why not fixed:} not required by the verification goal unless it later becomes required.
\end{itemize}

\subsection{\texorpdfstring{R-24: Auto-inlining cannot be turned on in the benchmarks and does not shrink the \texttt{.ast}}{R-24: Auto-inlining cannot be turned on in the benchmarks and does not shrink the .ast}}\label{reg:R-24}

\begin{itemize}
\item \textbf{What:} \texttt{ErasureConfig.auto\_\allowbreak{}inline\_\allowbreak{}typeclass\_\allowbreak{}dispatch} (defect D10 of S-2) is off by default, and the benchmark Makefile has no variable for it: the only way is to override \texttt{ERASURE\_\allowbreak{}CONFIG} on the \texttt{make} command line, which also replaces the \texttt{remove\_\allowbreak{}irrel\_\allowbreak{}constr\_\allowbreak{}args} setting that \texttt{PRUNE\_\allowbreak{}CONSTRUCTORS} adds. Commit \texttt{94a429c} expected the option to "shrink the AST bloat from typeclass dispatch" (2 to 4 times); the option changes only \texttt{.ast.inlinings}, not the \texttt{.ast}, and the constants it marks are shrunk only by peregrine's inlining, which keeps their declarations.
\item \textbf{Where:} \leanfile{LeanToLambdaBox/\allowbreak{}Erasure.lean}: \texttt{ErasureConfig.auto\_\allowbreak{}inline\_\allowbreak{}typeclass\_\allowbreak{}dispatch}, \texttt{visitMutual}; \texttt{benchmarks/\allowbreak{}via\_\allowbreak{}malfunction/\allowbreak{}Makefile} (\texttt{ERASURE\_\allowbreak{}CONFIG}).
\item \textbf{Reproduction:} \leanfile{tests/\allowbreak{}regress/\allowbreak{}auto\_\allowbreak{}inline.lean}: \texttt{on.ast} equals \texttt{default.ast}, and only \texttt{on.ast.inlinings} differs. \texttt{grep -rn auto\_\allowbreak{}inline benchmarks/} finds nothing.
\item \textbf{Impact:} informational: the option has no effect on any benchmark, and the size of the emitted program is unchanged.
\item \textbf{Why not fixed:} not a defect of the eraser's output. A benchmark switch is a feature of the benchmark pipeline that the verification goal does not require.
\end{itemize}

\subsection{\texorpdfstring{R-25: The inlining diagnosis note is not committed, and its candidate causes are wrong}{R-25: The inlining diagnosis note is not committed, and its candidate causes are wrong}}\label{reg:R-25}

\begin{itemize}
\item \textbf{What:} commit \texttt{94a429c} lists \texttt{references/\allowbreak{}inlining\_\allowbreak{}diagnosis.md} (defect D11 of S-2) among its changes, but \texttt{references/\allowbreak{}.gitignore} (\texttt{*}) excludes the file and no commit contains it. The note's conclusion holds: the kername of \texttt{instDecidableEqNat} is the same in its call sites, its declaration and the \texttt{.ast.inlinings} file, so the symptom of Zulip (Feb 17: the constant is listed but not inlined) lies in peregrine. Its three candidate causes (phase ordering, a filter on \ensuremath{\lambda}-shaped bodies, an allow-list) and its statement that inlining removes the declaration are wrong. Peregrine's inlining pass, MetaRocq 1.5.1 \texttt{EInlining.inline}, does not rewrite \texttt{tCase} scrutinees, and \texttt{if n = 0} puts \texttt{instDecidableEqNat n 0} in a scrutinee; inlining keeps all declarations. MetaRocq changed the \texttt{tCase} case to inline the scrutinee in commit \texttt{6b4d5ebc} (\texttt{erasure/\allowbreak{}theories/\allowbreak{}EInlining.v:30}), on its \texttt{9.1} branch and in no release; peregrine requires \texttt{rocq-metarocq-erasure-plugin} 1.5.1.
\item \textbf{Where:} \texttt{references/\allowbreak{}inlining\_\allowbreak{}diagnosis.md} (ignored, not in the repository); \texttt{peregrine-tool/\allowbreak{}\_\allowbreak{}build/\allowbreak{}default/\allowbreak{}src/\allowbreak{}extraction/\allowbreak{}EInlining.ml:32-33} (extracted from MetaRocq 1.5.1 \texttt{EInlining.v}).
\item \textbf{Reproduction:} \texttt{git log -{}-all -{}- references/\allowbreak{}inlining\_\allowbreak{}diagnosis.md} prints nothing, and \texttt{git check-ignore -v references/\allowbreak{}inlining\_\allowbreak{}diagnosis.md} prints \texttt{references/\allowbreak{}.gitignore:1:*}. For the program \texttt{scr} of S-2 (claim C5), compiled with \texttt{peregrine compile scr.ast unbox.config -{}-attributes=scr.ast.inlinings}, the \texttt{.mlf} keeps \texttt{(let (\$discr (apply \$def\_\allowbreak{}\_\allowbreak{}instDecidableEqNat \$n \ldots{})) (switch \$discr \ldots{}))} in \texttt{isZeroIf}, while the call outside a scrutinee, in \texttt{isZeroDec}, becomes \texttt{\$def\_\allowbreak{}\_\allowbreak{}Nat\_\allowbreak{}decEq}.
\item \textbf{Impact:} informational: a constant marked for inlining that occurs in the scrutinee of a \texttt{match} or \texttt{if}, such as \texttt{instDecidableEqNat}, stays a call in the compiled code.
\item \textbf{Why not fixed:} the cause is in peregrine's MetaRocq dependency, outside this repository, and the note is not part of the repository.
\end{itemize}

\subsection{\texorpdfstring{R-26: \texttt{[@inlined hint]} silences warning 55 without inlining anything}{R-26: [@inlined hint] silences warning 55 without inlining anything}}\label{reg:R-26}

\begin{itemize}
\item \textbf{What:} the Zarith calls of \texttt{nat-inline.ml} and \texttt{int-inline.ml} carry \texttt{[@inlined hint]} (defect D12 of S-2). With OCaml 4.14.2 without flambda, the generated code is the same as with \texttt{[@inlined]}: every Zarith call stays an out-of-line call; only warning 55 (\texttt{Cannot inline: Function information unavailable}) is no longer printed. The question of Zulip (Feb 17) whether these inlinings work has the answer no, for this compiler, and the warning that showed it is gone.
\item \textbf{Where:} \texttt{benchmarks/\allowbreak{}via\_\allowbreak{}malfunction/\allowbreak{}nat-inline.ml}, \texttt{benchmarks/\allowbreak{}via\_\allowbreak{}malfunction/\allowbreak{}int-inline.ml}.
\item \textbf{Reproduction:} in \texttt{benchmarks/\allowbreak{}via\_\allowbreak{}malfunction}, \texttt{make build/\allowbreak{}\textless{}id\textgreater{}/\allowbreak{}axioms.cmx FLAMBDA=0 MALFUNCTION\_\allowbreak{}NO\_\allowbreak{}FLAMBDA\_\allowbreak{}SWITCH=peregrine ARRAYML=JCFArrayOCaml4.ml} prints warning 55 20 times before S-2 and never after. Compiling the copied \texttt{nat.ml} and \texttt{int.ml} with \texttt{ocamlfind ocamlopt -package zarith -c -dcmm} gives the same Cmm before and after up to source locations, with 12 and 8 calls \texttt{app \textquotedbl{}camlZ\_\allowbreak{}\_\allowbreak{}\ldots{}\textquotedbl{}}.
\item \textbf{Impact:} informational: no change in the generated code.
\item \textbf{Why not fixed:} not a defect: the change does what its commit message says.
\end{itemize}

\subsection{\texorpdfstring{R-27: The benchmark Makefile builds \texttt{axioms.cmx} without waiting for \texttt{LeanArray.cmx}}{R-27: The benchmark Makefile builds axioms.cmx without waiting for LeanArray.cmx}}\label{reg:R-27}

\begin{itemize}
\item \textbf{What:} \texttt{axioms.ml} contains \texttt{include LeanArray}, but the rule for \texttt{\$(build)/\allowbreak{}axioms.cmx} lists \texttt{nat.cmx}, \texttt{int.cmx} and \texttt{eq.cmx} as prerequisites, not \texttt{LeanArray.cmx}. When \texttt{axioms.cmx} is compiled before \texttt{LeanArray.cmx}, OCaml prints \texttt{Warning 58 [no-cmx-file]: no cmx file was found in path for module LeanArray, and its interface was not compiled with -opaque} and compiles \texttt{axioms.cmx} without the optimization information of \texttt{LeanArray}.
\item \textbf{Where:} \texttt{benchmarks/\allowbreak{}via\_\allowbreak{}malfunction/\allowbreak{}Makefile}: rule \texttt{\$(build)/\allowbreak{}axioms.cmx}.
\item \textbf{Reproduction:} in \texttt{benchmarks/\allowbreak{}via\_\allowbreak{}malfunction}, with a build directory \texttt{B} that does not exist yet, \texttt{make build=B FLAMBDA=0 MALFUNCTION\_\allowbreak{}NO\_\allowbreak{}FLAMBDA\_\allowbreak{}SWITCH=peregrine ARRAYML=JCFArrayOCaml4.ml B/\allowbreak{}axioms.cmx} prints the warning. The target \texttt{B/\allowbreak{}\textless{}test\textgreater{}.cmx} in a fresh directory prints it too.
\item \textbf{Impact:} a serial build of \texttt{bin/\allowbreak{}\textless{}test\textgreater{}} is not affected, since the link rule lists \texttt{LeanArray.cmx} before \texttt{axioms.cmx}. A parallel build (\texttt{make -j}), or a build of \texttt{axioms.cmx} or \texttt{\textless{}test\textgreater{}.cmx} on its own, can compile \texttt{axioms.cmx} without cross-module information for \texttt{LeanArray}. That can change the machine code, and so the timings, of programs that use \texttt{Array}, not their results.
\item \textbf{Why not fixed:} not required by the verification goal unless it later becomes required; it concerns the build of the benchmark runtime, not the emitted programs.
\end{itemize}

\subsection{\texorpdfstring{R-28: The benchmark Makefile needs GNU make 4.4 and does not say so}{R-28: The benchmark Makefile needs GNU make 4.4 and does not say so}}\label{reg:R-28}

\begin{itemize}
\item \textbf{What:} \texttt{benchmarks/\allowbreak{}via\_\allowbreak{}malfunction/\allowbreak{}Makefile} uses two features that are new in GNU make 4.4: the function \texttt{\$(let ...)} in \texttt{register\_\allowbreak{}test}, and the special target \texttt{.NOTINTERMEDIATE}. Older versions do not reject them. There \texttt{\$(let ...)} expands to nothing, so the variable \texttt{benches} is empty and no rule \texttt{\$(build)/\allowbreak{}\textless{}test\textgreater{}\_\allowbreak{}main.ml} exists; and \texttt{.NOTINTERMEDIATE} is an ordinary target, so a file that make creates through a chain of pattern rules is removed as an intermediate file at the end of the run. Neither the Makefile nor \texttt{benchmarks/\allowbreak{}via\_\allowbreak{}malfunction/\allowbreak{}README.md} states the requirement.
\item \textbf{Where:} \texttt{benchmarks/\allowbreak{}via\_\allowbreak{}malfunction/\allowbreak{}Makefile}: \texttt{define register\_\allowbreak{}test} and \texttt{.NOTINTERMEDIATE:}; \texttt{benchmarks/\allowbreak{}via\_\allowbreak{}malfunction/\allowbreak{}README.md}.
\item \textbf{Reproduction:} with GNU make 4.3, in \texttt{benchmarks/\allowbreak{}via\_\allowbreak{}malfunction}, \texttt{make -n FLAMBDA=0 MALFUNCTION\_\allowbreak{}NO\_\allowbreak{}FLAMBDA\_\allowbreak{}SWITCH=peregrine ARRAYML=JCFArrayOCaml4.ml build=B bin=C C/\allowbreak{}even} stops with \texttt{No rule to make target 'C/\allowbreak{}even'}, where make 4.4.1 lists the whole build; \texttt{make -p} shows \texttt{benches :=} empty with make 4.3 and 20 programs with make 4.4.1. In the setting of \leanfile{tests/\allowbreak{}regress/\allowbreak{}makefile\_\allowbreak{}inlinings.lean}, make 4.3 ends the run with \texttt{rm \$(build)/\allowbreak{}rInl.ast.inlinings}: the file regenerated for peregrine is removed again.
\item \textbf{Impact:} low. With make older than 4.4 the benchmark binaries cannot be built, and the message does not name the cause. The dry-run tests \texttt{makefile\_\allowbreak{}cmi} and \texttt{makefile\_\allowbreak{}inlinings} pass with make 4.3, the version on CI's \texttt{ubuntu-latest}.
\item \textbf{Why not fixed:} not required by the verification goal unless it later becomes required.
\end{itemize}

